% Options for packages loaded elsewhere
\PassOptionsToPackage{unicode}{hyperref}
\PassOptionsToPackage{hyphens}{url}
\documentclass[
  a4paper]{report}
\usepackage{xcolor}
\usepackage[top=3.5cm,bottom=3cm,left=3.5cm,right=2cm]{geometry}
\usepackage{amsmath,amssymb}
\setcounter{secnumdepth}{5}
\usepackage{iftex}
\ifPDFTeX
  \usepackage[T1]{fontenc}
  \usepackage[utf8]{inputenc}
  \usepackage{textcomp} % provide euro and other symbols
\else % if luatex or xetex
  \usepackage{unicode-math} % this also loads fontspec
  \defaultfontfeatures{Scale=MatchLowercase}
  \defaultfontfeatures[\rmfamily]{Ligatures=TeX,Scale=1}
\fi
\usepackage{lmodern}
\ifPDFTeX\else
  % xetex/luatex font selection
  \setmainfont[]{Times New Roman}
\fi
% Use upquote if available, for straight quotes in verbatim environments
\IfFileExists{upquote.sty}{\usepackage{upquote}}{}
\IfFileExists{microtype.sty}{% use microtype if available
  \usepackage[]{microtype}
  \UseMicrotypeSet[protrusion]{basicmath} % disable protrusion for tt fonts
}{}
\usepackage{setspace}
\makeatletter
\@ifundefined{KOMAClassName}{% if non-KOMA class
  \IfFileExists{parskip.sty}{%
    \usepackage{parskip}
  }{% else
    \setlength{\parindent}{0pt}
    \setlength{\parskip}{6pt plus 2pt minus 1pt}}
}{% if KOMA class
  \KOMAoptions{parskip=half}}
\makeatother
\usepackage{longtable,booktabs,array}
\newcounter{none} % for unnumbered tables
\usepackage{calc} % for calculating minipage widths
% Correct order of tables after \paragraph or \subparagraph
\usepackage{etoolbox}
\makeatletter
\patchcmd\longtable{\par}{\if@noskipsec\mbox{}\fi\par}{}{}
\makeatother
% Allow footnotes in longtable head/foot
\IfFileExists{footnotehyper.sty}{\usepackage{footnotehyper}}{\usepackage{footnote}}
\makesavenoteenv{longtable}
\setlength{\emergencystretch}{3em} % prevent overfull lines
\providecommand{\tightlist}{%
  \setlength{\itemsep}{0pt}\setlength{\parskip}{0pt}}
\usepackage{bookmark}
\IfFileExists{xurl.sty}{\usepackage{xurl}}{} % add URL line breaks if available
\urlstyle{same}
\hypersetup{
  pdftitle={Application of AI in Detecting and Preventing Brute-Force Attacks on SSH Systems with Early Prediction},
  pdfauthor={FPT University - Information Assurance},
  hidelinks,
  pdfcreator={LaTeX via pandoc}}

\title{Application of AI in Detecting and Preventing Brute-Force Attacks
on SSH Systems with Early Prediction}
\author{FPT University - Information Assurance}
\date{2026}

\begin{document}
\maketitle

{
\setcounter{tocdepth}{2}
\tableofcontents
}
\setstretch{1.5}
\chapter{FRONT MATTER}\label{front-matter}

MINISTRY OF EDUCATION AND TRAINING

\textbf{FPT UNIVERSITY}

\textbf{CAPSTONE PROJECT}

\chapter{DESIGN AND IMPLEMENTATION OF AN INTELLIGENT SSH BRUTE-FORCE
ATTACK DETECTION AND PREVENTION SYSTEM USING ISOLATION FOREST WITH
EWMA-ADAPTIVE PERCENTILE DYNAMIC
THRESHOLDING}\label{design-and-implementation-of-an-intelligent-ssh-brute-force-attack-detection-and-prevention-system-using-isolation-forest-with-ewma-adaptive-percentile-dynamic-thresholding}

\textbf{Major:} Information Assurance

\textbf{Capstone Project Code:}

\textbf{Student Name:}

\textbf{Student ID:}

\textbf{Supervisor:}

\textbf{Hanoi, 2025}

\section{ABSTRACT}\label{abstract}

The Secure Shell (SSH) protocol is the predominant mechanism for remote
administration of Linux and Unix server infrastructure worldwide. SSH
brute-force attacks --- in which adversaries systematically attempt
large numbers of username--password combinations to gain unauthorized
access --- constitute one of the most persistent and prevalent threats
to Internet-facing servers. Conventional defense tools such as Fail2Ban
employ static, count-based thresholds (e.g., blocking an IP address
after five failed login attempts within ten minutes) that are
fundamentally unable to detect sophisticated attack variants, including
low-and-slow brute-force attacks, distributed attacks that spread
attempts across multiple source IP addresses, and credential stuffing
attacks that leverage credentials stolen from unrelated breaches. This
thesis presents the design, implementation, and comprehensive evaluation
of an intelligent SSH brute-force attack detection and prevention system
that leverages unsupervised machine learning, specifically the Isolation
Forest algorithm, combined with a novel EWMA-Adaptive Percentile dynamic
thresholding mechanism to enable early prediction of attacks before they
reach full intensity.

The proposed system was evaluated using a hybrid dataset of 174,250 SSH
log lines comprising real-world attack data collected from an SSH
honeypot (119,729 log lines from 679 unique IP addresses across 5 days
of operation) and simulated normal behavior data (54,521 log lines from
64 user accounts). A set of 14 behavioral features was engineered and
extracted over 5-minute sliding windows per source IP address, capturing
temporal, authentication, connection, and session characteristics. Three
unsupervised anomaly detection algorithms were systematically trained,
optimized, and compared: Isolation Forest (IF), Local Outlier Factor
(LOF), and One-Class Support Vector Machine (OCSVM). After
hyperparameter optimization (contamination = 0.01, max\_features = 0.75,
max\_samples = 512, n\_estimators = 500), the Isolation Forest model
achieved an accuracy of 90.31\%, an F1-score of 93.74\%, a recall of
96.75\%, and a false positive rate of 29.00\%. The OCSVM model achieved
the highest accuracy of 91.38\% and F1-score of 94.55\%, while LOF
achieved perfect recall of 100\% but exhibited the highest false
positive rate of 67.10\%. Feature importance analysis revealed that
temporal features --- session\_duration\_mean (5.50\%),
min\_inter\_attempt\_time (3.86\%), and mean\_inter\_attempt\_time
(2.61\%) --- dominate discrimination between automated attacks and
normal activity, outperforming the count-based features traditionally
relied upon by tools such as Fail2Ban.

The system was deployed as a complete end-to-end solution using a
Docker-based microservices architecture comprising 9 services: FastAPI
(API server and ML inference), React (monitoring dashboard),
Elasticsearch (storage and indexing), Logstash (log normalization),
Kibana (visualization), Fail2Ban (automated IP blocking), Redis
(caching), PostgreSQL (configuration persistence), and Nginx (reverse
proxy). The five-stage real-time detection pipeline achieves an
end-to-end latency of under 2 seconds, with AI processing stages
completing in under 100 milliseconds. The EWMA-Adaptive Percentile
dynamic threshold (alpha = 0.3, base\_percentile = 95,
sensitivity\_factor = 1.5) enables the system to adapt to evolving
traffic patterns and provides early warning capabilities, detecting
low-and-slow attacks within 3--5 attempts (approximately 2--3 minutes),
significantly outperforming traditional static threshold methods. Five
attack scenarios of escalating difficulty were designed and tested,
demonstrating robust detection across basic brute-force, distributed,
dictionary, credential stuffing, and low-and-slow attack types. The
system was validated with 51 passing unit tests across all components
and 5 integrated attack scenario evaluations.

\textbf{Keywords:} SSH brute-force detection, Isolation Forest, anomaly
detection, dynamic threshold, EWMA, unsupervised machine learning, ELK
Stack, intrusion detection system, cybersecurity, early prediction,
Docker microservices, Fail2Ban

\section{ACKNOWLEDGEMENT}\label{acknowledgement}

I would like to express my sincere gratitude to all those who have
contributed to the completion of this thesis.

First and foremost, I extend my deepest thanks to my thesis supervisor
for their invaluable guidance, rigorous feedback, and consistent
encouragement throughout the entire research process. Their expertise in
cybersecurity and machine learning has been instrumental in shaping the
direction, methodology, and quality of this work.

I am profoundly grateful to the faculty members of the Information
Assurance program at FPT University for providing a solid foundation in
cybersecurity principles, software engineering, and research
methodology. The knowledge and analytical skills I acquired during my
studies were essential to the successful execution of this project.

I would also like to thank the members of the thesis examination
committee for their time, careful reading, and constructive suggestions,
which have strengthened the rigor and clarity of this research.

Special appreciation goes to my classmates and colleagues who provided
technical assistance during the implementation phase, shared their
experiences with related tools and frameworks, and offered moral support
during the more challenging phases of the project. The collaborative and
intellectually stimulating environment at FPT University has been a
constant source of motivation.

I gratefully acknowledge the open-source community for developing and
maintaining the tools and frameworks that formed the technical backbone
of this research, including scikit-learn, FastAPI, the Elastic Stack,
Docker, React, and Fail2Ban. Without these freely available,
well-documented, and actively maintained projects, this work would not
have been possible.

Finally, I owe my deepest gratitude to my family for their unconditional
love, patience, and unwavering support throughout my entire academic
journey. Their belief in my abilities has been a constant source of
strength and inspiration, and this thesis is dedicated to them.

\section{TABLE OF CONTENTS}\label{table-of-contents}

\textbf{ABSTRACT}

\textbf{ACKNOWLEDGEMENT}

\textbf{TABLE OF CONTENTS}

\textbf{LIST OF FIGURES}

\textbf{LIST OF TABLES}

\textbf{ABBREVIATIONS}

\textbf{CHAPTER 1: INTRODUCTION}
\ldots\ldots\ldots\ldots\ldots\ldots\ldots\ldots\ldots\ldots\ldots\ldots\ldots\ldots\ldots\ldots..
1

1.1 Background and Motivation
\ldots\ldots\ldots\ldots\ldots\ldots\ldots\ldots\ldots\ldots\ldots\ldots\ldots\ldots\ldots\ldots..
1

1.2 Problem Statement
\ldots\ldots\ldots\ldots\ldots\ldots\ldots\ldots\ldots\ldots\ldots\ldots\ldots\ldots\ldots\ldots..
3

1.3 Research Questions
\ldots\ldots\ldots\ldots\ldots\ldots\ldots\ldots\ldots\ldots\ldots\ldots\ldots\ldots\ldots\ldots..
5

1.4 Research Objectives
\ldots\ldots\ldots\ldots\ldots\ldots\ldots\ldots\ldots\ldots\ldots\ldots\ldots\ldots\ldots\ldots..
6

1.5 Scope and Delimitations
\ldots\ldots\ldots\ldots\ldots\ldots\ldots\ldots\ldots\ldots\ldots\ldots\ldots\ldots\ldots\ldots..
7

1.6 Thesis Contributions
\ldots\ldots\ldots\ldots\ldots\ldots\ldots\ldots\ldots\ldots\ldots\ldots\ldots\ldots\ldots\ldots..
8

1.7 Thesis Organization
\ldots\ldots\ldots\ldots\ldots\ldots\ldots\ldots\ldots\ldots\ldots\ldots\ldots\ldots\ldots\ldots..
9

\textbf{CHAPTER 2: LITERATURE REVIEW}
\ldots\ldots\ldots\ldots\ldots\ldots\ldots\ldots\ldots\ldots\ldots\ldots\ldots\ldots\ldots\ldots..
11

2.1 The SSH Protocol
\ldots\ldots\ldots\ldots\ldots\ldots\ldots\ldots\ldots\ldots\ldots\ldots\ldots\ldots\ldots\ldots..
11

2.1.1 Protocol Architecture and Layers
\ldots\ldots\ldots\ldots\ldots\ldots\ldots\ldots\ldots\ldots\ldots\ldots\ldots\ldots\ldots\ldots..
11

2.1.2 Authentication Mechanisms
\ldots\ldots\ldots\ldots\ldots\ldots\ldots\ldots\ldots\ldots\ldots\ldots\ldots\ldots\ldots\ldots..
13

2.1.3 SSH Brute-Force Attack Variants
\ldots\ldots\ldots\ldots\ldots\ldots\ldots\ldots\ldots\ldots\ldots\ldots\ldots\ldots\ldots\ldots..
14

2.2 Traditional Detection and Prevention Methods
\ldots\ldots\ldots\ldots\ldots\ldots\ldots\ldots\ldots\ldots\ldots\ldots\ldots\ldots\ldots\ldots..
16

2.2.1 Fail2Ban and Rate-Limiting Approaches
\ldots\ldots\ldots\ldots\ldots\ldots\ldots\ldots\ldots\ldots\ldots\ldots\ldots\ldots\ldots\ldots..
16

2.2.2 Port Knocking and Other Hardening Techniques
\ldots\ldots\ldots\ldots\ldots\ldots\ldots\ldots\ldots\ldots\ldots\ldots\ldots\ldots\ldots\ldots..
17

2.2.3 Limitations of Traditional Approaches
\ldots\ldots\ldots\ldots\ldots\ldots\ldots\ldots\ldots\ldots\ldots\ldots\ldots\ldots\ldots\ldots..
18

2.3 Machine Learning for Intrusion Detection
\ldots\ldots\ldots\ldots\ldots\ldots\ldots\ldots\ldots\ldots\ldots\ldots\ldots\ldots\ldots\ldots..
19

2.3.1 Supervised vs.~Unsupervised Approaches
\ldots\ldots\ldots\ldots\ldots\ldots\ldots\ldots\ldots\ldots\ldots\ldots\ldots\ldots\ldots\ldots..
19

2.3.2 Isolation Forest
\ldots\ldots\ldots\ldots\ldots\ldots\ldots\ldots\ldots\ldots\ldots\ldots\ldots\ldots\ldots\ldots..
21

2.3.3 Local Outlier Factor
\ldots\ldots\ldots\ldots\ldots\ldots\ldots\ldots\ldots\ldots\ldots\ldots\ldots\ldots\ldots\ldots..
23

2.3.4 One-Class Support Vector Machine
\ldots\ldots\ldots\ldots\ldots\ldots\ldots\ldots\ldots\ldots\ldots\ldots\ldots\ldots\ldots\ldots..
24

2.4 Dynamic Thresholding Methods
\ldots\ldots\ldots\ldots\ldots\ldots\ldots\ldots\ldots\ldots\ldots\ldots\ldots\ldots\ldots\ldots..
25

2.4.1 EWMA for Anomaly Detection
\ldots\ldots\ldots\ldots\ldots\ldots\ldots\ldots\ldots\ldots\ldots\ldots\ldots\ldots\ldots\ldots..
25

2.4.2 Adaptive Percentile Methods
\ldots\ldots\ldots\ldots\ldots\ldots\ldots\ldots\ldots\ldots\ldots\ldots\ldots\ldots\ldots\ldots..
26

2.5 ELK Stack for Security Monitoring
\ldots\ldots\ldots\ldots\ldots\ldots\ldots\ldots\ldots\ldots\ldots\ldots\ldots\ldots\ldots\ldots..
27

2.6 Related Work
\ldots\ldots\ldots\ldots\ldots\ldots\ldots\ldots\ldots\ldots\ldots\ldots\ldots\ldots\ldots\ldots..
29

2.7 Research Gap and Positioning
\ldots\ldots\ldots\ldots\ldots\ldots\ldots\ldots\ldots\ldots\ldots\ldots\ldots\ldots\ldots\ldots..
32

\textbf{CHAPTER 3: METHODOLOGY}
\ldots\ldots\ldots\ldots\ldots\ldots\ldots\ldots\ldots\ldots\ldots\ldots\ldots\ldots\ldots\ldots..
34

3.1 System Architecture
\ldots\ldots\ldots\ldots\ldots\ldots\ldots\ldots\ldots\ldots\ldots\ldots\ldots\ldots\ldots\ldots..
34

3.1.1 Overall Architecture
\ldots\ldots\ldots\ldots\ldots\ldots\ldots\ldots\ldots\ldots\ldots\ldots\ldots\ldots\ldots\ldots..
34

3.1.2 Docker-Based Microservices Deployment
\ldots\ldots\ldots\ldots\ldots\ldots\ldots\ldots\ldots\ldots\ldots\ldots\ldots\ldots\ldots\ldots..
36

3.1.3 Data Flow Pipeline
\ldots\ldots\ldots\ldots\ldots\ldots\ldots\ldots\ldots\ldots\ldots\ldots\ldots\ldots\ldots\ldots..
38

3.2 Data Collection and Preparation
\ldots\ldots\ldots\ldots\ldots\ldots\ldots\ldots\ldots\ldots\ldots\ldots\ldots\ldots\ldots\ldots..
40

3.2.1 Honeypot Data Collection
\ldots\ldots\ldots\ldots\ldots\ldots\ldots\ldots\ldots\ldots\ldots\ldots\ldots\ldots\ldots\ldots..
40

3.2.2 Normal Behavior Simulation
\ldots\ldots\ldots\ldots\ldots\ldots\ldots\ldots\ldots\ldots\ldots\ldots\ldots\ldots\ldots\ldots..
42

3.2.3 Data Labeling and Splitting
\ldots\ldots\ldots\ldots\ldots\ldots\ldots\ldots\ldots\ldots\ldots\ldots\ldots\ldots\ldots\ldots..
44

3.3 Feature Engineering
\ldots\ldots\ldots\ldots\ldots\ldots\ldots\ldots\ldots\ldots\ldots\ldots\ldots\ldots\ldots\ldots..
45

3.3.1 Sliding Window Method
\ldots\ldots\ldots\ldots\ldots\ldots\ldots\ldots\ldots\ldots\ldots\ldots\ldots\ldots\ldots\ldots..
45

3.3.2 Feature Definition
\ldots\ldots\ldots\ldots\ldots\ldots\ldots\ldots\ldots\ldots\ldots\ldots\ldots\ldots\ldots\ldots..
47

3.3.3 Feature Normalization
\ldots\ldots\ldots\ldots\ldots\ldots\ldots\ldots\ldots\ldots\ldots\ldots\ldots\ldots\ldots\ldots..
50

3.3.4 Model Selection and Configuration
\ldots\ldots\ldots\ldots\ldots\ldots\ldots\ldots\ldots\ldots\ldots\ldots\ldots\ldots\ldots\ldots..
51

3.4 Dynamic Threshold Design
\ldots\ldots\ldots\ldots\ldots\ldots\ldots\ldots\ldots\ldots\ldots\ldots\ldots\ldots\ldots\ldots..
53

3.4.1 EWMA-Adaptive Percentile Algorithm
\ldots\ldots\ldots\ldots\ldots\ldots\ldots\ldots\ldots\ldots\ldots\ldots\ldots\ldots\ldots\ldots..
53

3.4.2 Parameter Selection
\ldots\ldots\ldots\ldots\ldots\ldots\ldots\ldots\ldots\ldots\ldots\ldots\ldots\ldots\ldots\ldots..
55

3.4.3 Two-Level Detection Mechanism
\ldots\ldots\ldots\ldots\ldots\ldots\ldots\ldots\ldots\ldots\ldots\ldots\ldots\ldots\ldots\ldots..
56

3.5 Real-Time Detection Pipeline
\ldots\ldots\ldots\ldots\ldots\ldots\ldots\ldots\ldots\ldots\ldots\ldots\ldots\ldots\ldots\ldots..
57

3.6 Evaluation Methodology
\ldots\ldots\ldots\ldots\ldots\ldots\ldots\ldots\ldots\ldots\ldots\ldots\ldots\ldots\ldots\ldots..
59

3.6.1 Evaluation Metrics
\ldots\ldots\ldots\ldots\ldots\ldots\ldots\ldots\ldots\ldots\ldots\ldots\ldots\ldots\ldots\ldots..
59

3.6.2 Attack Scenario Design
\ldots\ldots\ldots\ldots\ldots\ldots\ldots\ldots\ldots\ldots\ldots\ldots\ldots\ldots\ldots\ldots..
61

3.6.3 System Performance Evaluation
\ldots\ldots\ldots\ldots\ldots\ldots\ldots\ldots\ldots\ldots\ldots\ldots\ldots\ldots\ldots\ldots..
62

\textbf{CHAPTER 4: RESULTS AND ANALYSIS}
\ldots\ldots\ldots\ldots\ldots\ldots\ldots\ldots\ldots\ldots\ldots\ldots\ldots\ldots\ldots\ldots..
64

4.1 Dataset Analysis
\ldots\ldots\ldots\ldots\ldots\ldots\ldots\ldots\ldots\ldots\ldots\ldots\ldots\ldots\ldots\ldots..
64

4.1.1 Overall Statistics
\ldots\ldots\ldots\ldots\ldots\ldots\ldots\ldots\ldots\ldots\ldots\ldots\ldots\ldots\ldots\ldots..
64

4.1.2 Honeypot Data Analysis
\ldots\ldots\ldots\ldots\ldots\ldots\ldots\ldots\ldots\ldots\ldots\ldots\ldots\ldots\ldots\ldots..
65

4.1.3 Simulation Data Analysis
\ldots\ldots\ldots\ldots\ldots\ldots\ldots\ldots\ldots\ldots\ldots\ldots\ldots\ldots\ldots\ldots..
67

4.2 Feature Analysis
\ldots\ldots\ldots\ldots\ldots\ldots\ldots\ldots\ldots\ldots\ldots\ldots\ldots\ldots\ldots\ldots..
68

4.2.1 Feature Distributions
\ldots\ldots\ldots\ldots\ldots\ldots\ldots\ldots\ldots\ldots\ldots\ldots\ldots\ldots\ldots\ldots..
68

4.2.2 Correlation Analysis
\ldots\ldots\ldots\ldots\ldots\ldots\ldots\ldots\ldots\ldots\ldots\ldots\ldots\ldots\ldots\ldots..
70

4.2.3 Feature Importance
\ldots\ldots\ldots\ldots\ldots\ldots\ldots\ldots\ldots\ldots\ldots\ldots\ldots\ldots\ldots\ldots..
71

4.3 Model Performance
\ldots\ldots\ldots\ldots\ldots\ldots\ldots\ldots\ldots\ldots\ldots\ldots\ldots\ldots\ldots\ldots..
73

4.3.1 Baseline Results
\ldots\ldots\ldots\ldots\ldots\ldots\ldots\ldots\ldots\ldots\ldots\ldots\ldots\ldots\ldots\ldots..
73

4.3.2 Anomaly Score Distributions
\ldots\ldots\ldots\ldots\ldots\ldots\ldots\ldots\ldots\ldots\ldots\ldots\ldots\ldots\ldots\ldots..
75

4.3.3 ROC and Precision-Recall Analysis
\ldots\ldots\ldots\ldots\ldots\ldots\ldots\ldots\ldots\ldots\ldots\ldots\ldots\ldots\ldots\ldots..
76

4.3.4 Confusion Matrix Analysis
\ldots\ldots\ldots\ldots\ldots\ldots\ldots\ldots\ldots\ldots\ldots\ldots\ldots\ldots\ldots\ldots..
77

4.3.5 Baseline Summary
\ldots\ldots\ldots\ldots\ldots\ldots\ldots\ldots\ldots\ldots\ldots\ldots\ldots\ldots\ldots\ldots..
78

4.4 Interpretation of Results
\ldots\ldots\ldots\ldots\ldots\ldots\ldots\ldots\ldots\ldots\ldots\ldots\ldots\ldots\ldots\ldots..
79

4.4.1 Optimized Model Performance
\ldots\ldots\ldots\ldots\ldots\ldots\ldots\ldots\ldots\ldots\ldots\ldots\ldots\ldots\ldots\ldots..
79

4.4.2 Dynamic Threshold Evaluation
\ldots\ldots\ldots\ldots\ldots\ldots\ldots\ldots\ldots\ldots\ldots\ldots\ldots\ldots\ldots\ldots..
81

4.4.3 Comparison with Related Studies
\ldots\ldots\ldots\ldots\ldots\ldots\ldots\ldots\ldots\ldots\ldots\ldots\ldots\ldots\ldots\ldots..
83

4.5 Attack Scenario Evaluation
\ldots\ldots\ldots\ldots\ldots\ldots\ldots\ldots\ldots\ldots\ldots\ldots\ldots\ldots\ldots\ldots..
84

4.5.1 Scenario Descriptions
\ldots\ldots\ldots\ldots\ldots\ldots\ldots\ldots\ldots\ldots\ldots\ldots\ldots\ldots\ldots\ldots..
84

4.5.2 Detection Results
\ldots\ldots\ldots\ldots\ldots\ldots\ldots\ldots\ldots\ldots\ldots\ldots\ldots\ldots\ldots\ldots..
86

4.5.3 Low-and-Slow Detection Analysis
\ldots\ldots\ldots\ldots\ldots\ldots\ldots\ldots\ldots\ldots\ldots\ldots\ldots\ldots\ldots\ldots..
87

4.6 System Performance
\ldots\ldots\ldots\ldots\ldots\ldots\ldots\ldots\ldots\ldots\ldots\ldots\ldots\ldots\ldots\ldots..
89

4.6.1 Latency Analysis
\ldots\ldots\ldots\ldots\ldots\ldots\ldots\ldots\ldots\ldots\ldots\ldots\ldots\ldots\ldots\ldots..
89

4.6.2 Throughput Analysis
\ldots\ldots\ldots\ldots\ldots\ldots\ldots\ldots\ldots\ldots\ldots\ldots\ldots\ldots\ldots\ldots..
90

4.6.3 Resource Utilization
\ldots\ldots\ldots\ldots\ldots\ldots\ldots\ldots\ldots\ldots\ldots\ldots\ldots\ldots\ldots\ldots..
91

\textbf{CHAPTER 5: DISCUSSION}
\ldots\ldots\ldots\ldots\ldots\ldots\ldots\ldots\ldots\ldots\ldots\ldots\ldots\ldots\ldots\ldots..
93

5.1 Restatement of Research Problem
\ldots\ldots\ldots\ldots\ldots\ldots\ldots\ldots\ldots\ldots\ldots\ldots\ldots\ldots\ldots\ldots..
93

5.2 Achievement of Research Objectives
\ldots\ldots\ldots\ldots\ldots\ldots\ldots\ldots\ldots\ldots\ldots\ldots\ldots\ldots\ldots\ldots..
95

5.2.1 Objective 1: Behavioral Feature Set Design
\ldots\ldots\ldots\ldots\ldots\ldots\ldots\ldots\ldots\ldots\ldots\ldots\ldots\ldots\ldots\ldots..
95

5.2.2 Objective 2: Model Training and Evaluation
\ldots\ldots\ldots\ldots\ldots\ldots\ldots\ldots\ldots\ldots\ldots\ldots\ldots\ldots\ldots\ldots..
96

5.2.3 Objective 3: Dynamic Threshold
\ldots\ldots\ldots\ldots\ldots\ldots\ldots\ldots\ldots\ldots\ldots\ldots\ldots\ldots\ldots\ldots..
97

5.2.4 Objective 4: Integrated System
\ldots\ldots\ldots\ldots\ldots\ldots\ldots\ldots\ldots\ldots\ldots\ldots\ldots\ldots\ldots\ldots..
98

5.2.5 Objective 5: Attack Scenario Evaluation
\ldots\ldots\ldots\ldots\ldots\ldots\ldots\ldots\ldots\ldots\ldots\ldots\ldots\ldots\ldots\ldots..
99

5.3 Interpretation of Key Findings
\ldots\ldots\ldots\ldots\ldots\ldots\ldots\ldots\ldots\ldots\ldots\ldots\ldots\ldots\ldots\ldots..
100

5.3.1 Why Timing Features Dominate
\ldots\ldots\ldots\ldots\ldots\ldots\ldots\ldots\ldots\ldots\ldots\ldots\ldots\ldots\ldots\ldots..
100

5.3.2 Why Isolation Forest Is Chosen Despite OCSVM Having Higher F1
\ldots\ldots\ldots\ldots\ldots\ldots\ldots\ldots\ldots\ldots\ldots\ldots\ldots\ldots\ldots\ldots..
102

5.3.3 Dynamic Threshold vs.~Static Threshold Advantage
\ldots\ldots\ldots\ldots\ldots\ldots\ldots\ldots\ldots\ldots\ldots\ldots\ldots\ldots\ldots\ldots..
103

5.3.4 FPR Analysis and Acceptable Trade-off in Security Context
\ldots\ldots\ldots\ldots\ldots\ldots\ldots\ldots\ldots\ldots\ldots\ldots\ldots\ldots\ldots\ldots..
105

5.4 Comparison with Existing Literature
\ldots\ldots\ldots\ldots\ldots\ldots\ldots\ldots\ldots\ldots\ldots\ldots\ldots\ldots\ldots\ldots..
107

5.5 Practical Implications
\ldots\ldots\ldots\ldots\ldots\ldots\ldots\ldots\ldots\ldots\ldots\ldots\ldots\ldots\ldots\ldots..
110

5.5.1 Deployment Considerations
\ldots\ldots\ldots\ldots\ldots\ldots\ldots\ldots\ldots\ldots\ldots\ldots\ldots\ldots\ldots\ldots..
110

5.5.2 Scalability Considerations
\ldots\ldots\ldots\ldots\ldots\ldots\ldots\ldots\ldots\ldots\ldots\ldots\ldots\ldots\ldots\ldots..
112

5.5.3 Cost-Benefit Analysis
\ldots\ldots\ldots\ldots\ldots\ldots\ldots\ldots\ldots\ldots\ldots\ldots\ldots\ldots\ldots\ldots..
113

5.6 Limitations and Threats to Validity
\ldots\ldots\ldots\ldots\ldots\ldots\ldots\ldots\ldots\ldots\ldots\ldots\ldots\ldots\ldots\ldots..
114

5.7 Chapter Summary
\ldots\ldots\ldots\ldots\ldots\ldots\ldots\ldots\ldots\ldots\ldots\ldots\ldots\ldots\ldots\ldots..
116

\textbf{CHAPTER 6: CONCLUSION AND FUTURE WORK}
\ldots\ldots\ldots\ldots\ldots\ldots\ldots\ldots\ldots\ldots\ldots\ldots\ldots\ldots\ldots\ldots..
117

6.1 Conclusion
\ldots\ldots\ldots\ldots\ldots\ldots\ldots\ldots\ldots\ldots\ldots\ldots\ldots\ldots\ldots\ldots..
117

6.2 Contributions Summary
\ldots\ldots\ldots\ldots\ldots\ldots\ldots\ldots\ldots\ldots\ldots\ldots\ldots\ldots\ldots\ldots..
120

6.3 Future Work
\ldots\ldots\ldots\ldots\ldots\ldots\ldots\ldots\ldots\ldots\ldots\ldots\ldots\ldots\ldots\ldots..
122

6.3.1 Deep Learning for Improved Detection
\ldots\ldots\ldots\ldots\ldots\ldots\ldots\ldots\ldots\ldots\ldots\ldots\ldots\ldots\ldots\ldots..
122

6.3.2 Online Learning and Continuous Adaptation
\ldots\ldots\ldots\ldots\ldots\ldots\ldots\ldots\ldots\ldots\ldots\ldots\ldots\ldots\ldots\ldots..
123

6.3.3 Multi-Protocol Extension
\ldots\ldots\ldots\ldots\ldots\ldots\ldots\ldots\ldots\ldots\ldots\ldots\ldots\ldots\ldots\ldots..
124

6.3.4 Federated Detection
\ldots\ldots\ldots\ldots\ldots\ldots\ldots\ldots\ldots\ldots\ldots\ldots\ldots\ldots\ldots\ldots..
125

6.3.5 Threat Intelligence Integration
\ldots\ldots\ldots\ldots\ldots\ldots\ldots\ldots\ldots\ldots\ldots\ldots\ldots\ldots\ldots\ldots..
126

6.3.6 SOAR Integration
\ldots\ldots\ldots\ldots\ldots\ldots\ldots\ldots\ldots\ldots\ldots\ldots\ldots\ldots\ldots\ldots..
127

6.3.7 Explainability and Interpretability
\ldots\ldots\ldots\ldots\ldots\ldots\ldots\ldots\ldots\ldots\ldots\ldots\ldots\ldots\ldots\ldots..
128

6.3.8 Long-Term IP Profiling
\ldots\ldots\ldots\ldots\ldots\ldots\ldots\ldots\ldots\ldots\ldots\ldots\ldots\ldots\ldots\ldots..
129

6.3.9 Real-World Deployment Validation
\ldots\ldots\ldots\ldots\ldots\ldots\ldots\ldots\ldots\ldots\ldots\ldots\ldots\ldots\ldots\ldots..
130

\textbf{REFERENCES}
\ldots\ldots\ldots\ldots\ldots\ldots\ldots\ldots\ldots\ldots\ldots\ldots\ldots\ldots\ldots\ldots..
131

\textbf{APPENDICES}
\ldots\ldots\ldots\ldots\ldots\ldots\ldots\ldots\ldots\ldots\ldots\ldots\ldots\ldots\ldots\ldots..
137

Appendix A: Feature Extraction Code Snippets
\ldots\ldots\ldots\ldots\ldots\ldots\ldots\ldots\ldots\ldots\ldots\ldots\ldots\ldots\ldots\ldots..
137

Appendix B: Docker Compose Configuration
\ldots\ldots\ldots\ldots\ldots\ldots\ldots\ldots\ldots\ldots\ldots\ldots\ldots\ldots\ldots\ldots..
139

Appendix C: Isolation Forest Hyperparameter Tuning Results
\ldots\ldots\ldots\ldots\ldots\ldots\ldots\ldots\ldots\ldots\ldots\ldots\ldots\ldots\ldots\ldots..
141

Appendix D: Attack Simulation Scripts
\ldots\ldots\ldots\ldots\ldots\ldots\ldots\ldots\ldots\ldots\ldots\ldots\ldots\ldots\ldots\ldots..
143

Appendix E: Logstash Pipeline Configuration
\ldots\ldots\ldots\ldots\ldots\ldots\ldots\ldots\ldots\ldots\ldots\ldots\ldots\ldots\ldots\ldots..
145

Appendix F: Dynamic Threshold Parameter Sensitivity Analysis
\ldots\ldots\ldots\ldots\ldots\ldots\ldots\ldots\ldots\ldots\ldots\ldots\ldots\ldots\ldots\ldots..
147

\section{LIST OF FIGURES}\label{list-of-figures}

\begin{itemize}
\tightlist
\item
  Figure 1.1: Global trends in SSH brute-force attacks (2018--2025)
\item
  Figure 1.2: Comparison of traditional vs.~AI-based attack detection
  approaches
\item
  Figure 1.3: Overview of the proposed system architecture (block
  diagram)
\item
  Figure 2.1: Layered architecture of the SSH-2 protocol
\item
  Figure 2.2: Illustration of static threshold vs.~dynamic threshold
  (EWMA, Adaptive Percentile, and hybrid) on the same anomaly score
  series
\item
  Figure 2.3: Architecture for integrating AI models with the ELK Stack
  for SSH monitoring
\item
  Figure 2.4: Venn diagram showing the research positioning at the
  intersection of three domains
\item
  Figure 3.1: Overall system architecture for SSH brute-force attack
  detection
\item
  Figure 3.2: Data flow pipeline of the system
\item
  Figure 3.3: Illustration of the sliding window method (window = 5 min,
  stride = 1 min)
\item
  Figure 3.4: Model training and evaluation workflow
\item
  Figure 3.5: Illustration of the EWMA-Adaptive Percentile dynamic
  threshold algorithm
\item
  Figure 3.6: Five-stage real-time detection pipeline
\item
  Figure 4.1: Distribution of failed authentication attempts per IP
  (honeypot)
\item
  Figure 4.2: Comparison of hourly activity distribution between
  honeypot (attack) and simulation (normal)
\item
  Figure 4.3: Boxplot comparison of key feature distributions between
  Normal and Attack classes
\item
  Figure 4.4: Histogram distribution of 14 features, separated by
  normal/attack labels
\item
  Figure 4.5: Pearson correlation matrix heatmap of 14 features
\item
  Figure 4.6: Bar chart of feature importance for 14 features
  (permutation importance on Isolation Forest)
\item
  Figure 4.7: Anomaly score distribution on the test set for three
  models (IF, LOF, OCSVM)
\item
  Figure 4.8: Violin plot comparing anomaly score distributions between
  normal and attack for three models
\item
  Figure 4.9: Radar chart comparing five performance metrics across
  three models
\item
  Figure 4.10: ROC curves for three models on the test set
\item
  Figure 4.11: Precision-Recall curves for three models on the test set
\item
  Figure 4.12: Anomaly score and dynamic threshold evolution over time
  on test data
\item
  Figure 4.13: Effect of alpha, base\_percentile, and
  sensitivity\_factor on F1-Score (parameter sensitivity analysis)
\item
  Figure 4.14: Heatmap of detection rates by attack scenario and model
\item
  Figure 5.1: Comparative positioning of the proposed system against
  existing studies
\end{itemize}

\section{LIST OF TABLES}\label{list-of-tables}

\begin{itemize}
\tightlist
\item
  Table 1.1: Summary of research challenges and proposed solutions
\item
  Table 1.2: Research objectives and corresponding approaches
\item
  Table 2.1: Comparison of characteristics of SSH brute-force attack
  variants
\item
  Table 2.2: Comparison of advantages and disadvantages of traditional
  detection methods
\item
  Table 2.3: Comparison of Isolation Forest, LOF, and One-Class SVM
  characteristics
\item
  Table 2.4: Comparative summary of related research works
\item
  Table 3.1: List of 9 Docker services in the system with ports and
  roles
\item
  Table 3.2: Summary of two data sources (honeypot and simulation)
\item
  Table 3.3: Dataset split summary (training, testing, normal, attack)
\item
  Table 3.4: Summary of 14 behavioral features and their significance
\item
  Table 3.5: Comparison of characteristics of three anomaly detection
  models
\item
  Table 3.6: Dynamic threshold algorithm parameters and their meanings
\item
  Table 3.7: Summary of evaluation metrics and their security
  significance
\item
  Table 4.1: Overall dataset statistics (log lines, samples, labels)
\item
  Table 4.2: Detailed honeypot data statistics (IPs, events, time range)
\item
  Table 4.3: Detailed simulation data statistics (users, events, session
  types)
\item
  Table 4.4: Descriptive statistics of key features by label (normal
  vs.~attack)
\item
  Table 4.5: Feature importance ranking (Top 5 features by permutation
  importance)
\item
  Table 4.6: Performance comparison of three models on the test set
  (baseline configuration)
\item
  Table 4.7: Confusion matrix estimates for three models (TP, TN, FP,
  FN)
\item
  Table 4.8: Comparison with related studies (method, dataset, metrics)
\item
  Table 4.9: Dynamic threshold vs.~static threshold comparison (OCSVM
  model)
\item
  Table 4.10: Description of 5 attack scenarios (type, speed, IP count,
  difficulty)
\item
  Table 4.11: Low-and-Slow detection results by model (detection rate,
  time to first alert)
\item
  Table 4.12: Summary of detection results across 5 attack scenarios
  (all models)
\item
  Table 4.13: Test environment configuration (hardware, software,
  versions)
\item
  Table 4.14: Latency analysis by pipeline stage (ingestion,
  aggregation, scoring, decision, action)
\item
  Table 4.15: Processing throughput by model (events/second, scoring
  latency)
\item
  Table 4.16: Resource utilization of main Docker services (CPU, RAM,
  disk)
\item
  Table 5.1: Optimized model performance comparison (IF, LOF, OCSVM
  after tuning)
\item
  Table 5.2: Comprehensive comparison with related works (12 studies)
\end{itemize}

\section{ABBREVIATIONS}\label{abbreviations}

{\def\LTcaptype{none} % do not increment counter
\begin{longtable}[]{@{}ll@{}}
\toprule\noalign{}
Abbreviation & Full Form \\
\midrule\noalign{}
\endhead
\bottomrule\noalign{}
\endlastfoot
AI & Artificial Intelligence \\
API & Application Programming Interface \\
AUC & Area Under the Curve \\
BST & Binary Search Tree \\
CPU & Central Processing Unit \\
CSS & Cascading Style Sheets \\
DNS & Domain Name System \\
ELK & Elasticsearch, Logstash, Kibana \\
EWMA & Exponentially Weighted Moving Average \\
FN & False Negative \\
FP & False Positive \\
FPR & False Positive Rate \\
FTP & File Transfer Protocol \\
GDPR & General Data Protection Regulation \\
GPU & Graphics Processing Unit \\
HTML & Hypertext Markup Language \\
HTTP & Hypertext Transfer Protocol \\
IDS & Intrusion Detection System \\
IF & Isolation Forest \\
IETF & Internet Engineering Task Force \\
IoC & Indicator of Compromise \\
IoT & Internet of Things \\
IP & Internet Protocol \\
IQR & Interquartile Range \\
ISO & International Organization for Standardization \\
LOF & Local Outlier Factor \\
LSTM & Long Short-Term Memory \\
ML & Machine Learning \\
MTTD & Mean Time to Detection \\
MTTR & Mean Time to Respond \\
NCSC & National Cyber Security Center \\
NIDS & Network Intrusion Detection System \\
OCSVM & One-Class Support Vector Machine \\
PAM & Pluggable Authentication Module \\
RAM & Random Access Memory \\
RBF & Radial Basis Function \\
RDP & Remote Desktop Protocol \\
REST & Representational State Transfer \\
RFC & Request for Comments \\
ROC & Receiver Operating Characteristic \\
SCP & Secure Copy Protocol \\
SFTP & SSH File Transfer Protocol \\
SHAP & SHapley Additive exPlanations \\
SIEM & Security Information and Event Management \\
SMTP & Simple Mail Transfer Protocol \\
SOAR & Security Orchestration, Automation and Response \\
SOC & Security Operations Center \\
SSH & Secure Shell \\
STIX & Structured Threat Information eXpression \\
SVM & Support Vector Machine \\
TAXII & Trusted Automated eXchange of Intelligence Information \\
TCP & Transmission Control Protocol \\
TN & True Negative \\
TP & True Positive \\
TPR & True Positive Rate \\
\end{longtable}
}

\chapter{CHAPTER 1: INTRODUCTION}\label{chapter-1-introduction}

\section{1.1 Background}\label{background}

The Secure Shell (SSH) protocol, standardized through RFC 4251 by Ylonen
and Lonvick {[}1{]}, remains the predominant mechanism for remote
administration of server infrastructure worldwide. Originally conceived
as a secure replacement for plaintext protocols such as Telnet and rsh,
SSH has evolved into the de facto standard for encrypted remote access,
file transfer, and tunneling across enterprise, academic, and government
networks {[}2{]}. The protocol's ubiquity, however, has rendered it one
of the most targeted attack surfaces in the contemporary threat
landscape. Wu et al.~{[}3{]} analyzed data from the National Center for
Supercomputing Applications (NCSA) honeypot infrastructure and
documented approximately 11 billion SSH brute-force login attempts over
a multi-year observation period, underscoring the industrial scale at
which these attacks are conducted.

The magnitude of the SSH brute-force problem is further corroborated by
industry threat intelligence. The Verizon Data Breach Investigations
Report (DBIR) {[}4{]} consistently identifies stolen credentials as the
single most common attack vector in confirmed data breaches, accounting
for approximately 29\% of all breaches analyzed. Because SSH password
authentication provides a direct path to credential theft through
systematic enumeration, SSH brute-force attacks constitute a substantial
proportion of this attack category. Hellemons et al.~{[}5{]}
demonstrated through their SSHCure three-phase model---comprising
scanning, brute-force, and compromise phases---that SSH attacks follow
predictable temporal patterns that, in principle, are amenable to
automated detection if sufficiently discriminative features are
extracted from authentication logs.

Traditional countermeasures against SSH brute-force attacks have relied
predominantly on static, rule-based mechanisms. Fail2Ban {[}6{]}, the
most widely deployed open-source intrusion prevention tool for SSH,
monitors authentication log files and blocks offending IP addresses via
firewall rules when the number of failed login attempts exceeds a
preconfigured threshold (typically 5 failures within a 600-second
window). While effective against naive, high-speed attacks, this
approach exhibits well-documented limitations against sophisticated
adversaries who employ slow-rate, distributed, or credential-stuffing
strategies {[}7{]}. The fundamental inadequacy of static thresholds in
the face of adaptive adversaries has motivated a growing body of
research into machine-learning-based detection approaches.

The application of machine learning to network intrusion detection has
been the subject of extensive investigation over the past two decades.
Chandola, Banerjee, and Kumar {[}7{]} provided a seminal survey of
anomaly detection techniques, establishing a taxonomy that distinguishes
among classification-based, nearest-neighbor-based, clustering-based,
and isolation-based approaches. Within this taxonomy, unsupervised and
semi-supervised methods hold particular appeal for intrusion detection
because they do not require labeled attack data for training---a
significant practical advantage given the difficulty and expense of
obtaining comprehensive, accurately labeled attack datasets in
operational environments {[}12{]}. Buczak and Guven {[}12{]} surveyed
data mining and machine learning methods for cybersecurity and concluded
that anomaly-detection-based approaches offer superior generalization to
novel attack variants compared to signature-based methods, albeit at the
cost of elevated false positive rates.

Among unsupervised anomaly detection algorithms, Isolation Forest (IF),
proposed by Liu, Ting, and Zhou {[}8, 9{]}, has attracted particular
attention for network security applications due to its linear time
complexity, absence of distributional assumptions, and native
suitability for high-dimensional data. Unlike density-based or
distance-based methods, Isolation Forest detects anomalies through the
principle of isolation: anomalous data points, by virtue of their
atypical feature values, require fewer random partitions to be separated
from the remainder of the dataset. This computational paradigm yields
anomaly scores that are both interpretable and efficiently computable,
properties that are essential for real-time intrusion detection systems.
Complementary algorithms including Local Outlier Factor (LOF) {[}10{]}
and One-Class Support Vector Machine (OCSVM) {[}11{]} provide
alternative detection paradigms based on local density estimation and
maximum-margin separation, respectively, enabling rigorous comparative
evaluation.

The convergence of three developments---the escalating scale and
sophistication of SSH brute-force attacks {[}3, 4{]}, the maturation of
efficient anomaly detection algorithms {[}8, 9, 10, 11{]}, and the
availability of scalable log analytics infrastructure such as the ELK
Stack (Elasticsearch, Logstash, Kibana)---creates a compelling
opportunity for the development of intelligent, adaptive SSH security
systems. This thesis capitalizes on this convergence by proposing a
semi-supervised anomaly detection system that combines Isolation Forest
with a novel EWMA-Adaptive Percentile dynamic thresholding mechanism
{[}13, 14{]}, trained exclusively on normal SSH behavior and evaluated
against real-world attack data collected from a production honeypot. The
system is fully integrated with ELK Stack for log ingestion and
visualization, Fail2Ban for automated response, and Docker for
containerized deployment, thereby bridging the persistent gap between
research prototypes and production-ready security tools {[}12, 15{]}.

The growing reliance on cloud computing and Infrastructure-as-a-Service
(IaaS) models has further amplified the urgency of SSH security
research. All major cloud providers use SSH as the default protocol for
initial access to virtual machine instances, and the proliferation of
containerized microservice architectures has multiplied the number of
SSH-accessible endpoints within typical enterprise environments {[}1,
2{]}. The financial consequences of successful SSH compromises extend
far beyond the immediate cost of the intrusion: once an attacker gains
SSH access, the compromise typically escalates through lateral movement,
data exfiltration, ransomware deployment, or the installation of
cryptocurrency miners {[}4{]}. The Verizon DBIR {[}4{]} estimates that
compromised credentials are the most common initial attack vector,
responsible for the largest share of all confirmed breaches. These
cascading consequences underscore the need for proactive, predictive
defense mechanisms that can identify attack intent before credential
compromise occurs.

The dataset employed in this research comprises 174,250 SSH log lines:
119,729 lines of real attack traffic collected from a Virtual Private
Server (VPS) honeypot (encompassing 679 unique attacking IP addresses,
29,301 failed password attempts, and 532 accepted root logins) and
54,521 lines of simulated normal behavior (64 user accounts performing
routine SSH operations). This combination of authentic attack data and
controlled normal-behavior data provides a realistic foundation for
model training and evaluation that is absent from studies relying
exclusively on synthetic benchmark datasets {[}24{]}. Ring et
al.~{[}24{]} surveyed 34 datasets for network-based intrusion detection
and found that many widely used benchmarks do not adequately represent
modern attack characteristics, recommending the use of
application-specific data collected from controlled environments or
honeypots. The system extracts 14 behavioral features per IP address per
5-minute sliding window and applies the Isolation Forest algorithm with
dynamic thresholding to achieve an F1-score of 93.74\% and an accuracy
of 90.31\%, with comparative benchmarks from LOF (F1 = 89.94\%) and
OCSVM (F1 = 94.55\%) providing context for the primary model's
performance.

\section{1.2 Problem Statement}\label{problem-statement}

Despite decades of research and the deployment of numerous
countermeasures, SSH brute-force attacks remain a persistent and growing
threat to networked systems. Analysis of the literature and empirical
observation of real-world attack traffic reveal four critical gaps in
the current state of SSH brute-force defense that this thesis seeks to
address.

\textbf{Gap 1: Static thresholds cannot adapt to dynamic traffic
patterns.} Traditional detection tools, exemplified by Fail2Ban {[}6{]},
rely on fixed thresholds---for example, blocking an IP address after a
predetermined number of failed login attempts within a fixed time
window. This approach creates an inherent dilemma: thresholds set too
low generate excessive false positives during periods of legitimate
high-activity (e.g., system maintenance windows, shift changes), while
thresholds set too high permit sophisticated slow-rate attacks to
proceed undetected {[}7, 12{]}. SSH traffic exhibits pronounced temporal
variability---business hours versus off-hours, weekdays versus
weekends---that static thresholds are structurally incapable of
accommodating. Lucas and Saccucci {[}13{]} demonstrated in the context
of statistical process control that exponentially weighted moving
average (EWMA) schemes provide superior sensitivity to small, persistent
shifts in process mean compared to fixed-limit schemes, a finding
directly applicable to the anomaly-score time series generated by SSH
monitoring systems.

\textbf{Gap 2: Absence of predictive, early-warning capability.} The
vast majority of deployed SSH defense systems operate in a purely
reactive mode: they detect and respond to attacks only after a
sufficient volume of malicious activity has already occurred {[}5, 7{]}.
Hellemons et al.~{[}5{]} documented the three-phase structure of SSH
attacks (scanning, brute-force, compromise) and noted that the scanning
and early brute-force phases exhibit distinctive behavioral signatures
that could, in principle, enable detection before the attack reaches its
critical phase. However, existing systems rarely exploit these
early-phase signatures. Javed and Paxson {[}23{]} studied stealthy SSH
brute-forcing campaigns that unfold over periods of days to weeks,
demonstrating that attacks operating below conventional detection
thresholds can achieve compromise rates comparable to aggressive
campaigns while remaining invisible to standard monitoring tools. The
absence of early-warning capability means that defenders forfeit the
opportunity to intervene during the reconnaissance phase, when the cost
of mitigation is lowest.

\textbf{Gap 3: Insufficient adaptation to evolving attack behaviors.}
Modern SSH brute-force attacks have evolved far beyond simple high-speed
password enumeration. Bezerra et al.~{[}18{]} analyzed SSH attack
durations and identified significant variability in attack temporal
profiles, ranging from sub-second bursts to campaigns spanning multiple
days. Owens and Matthews {[}19{]} demonstrated that credential-stuffing
attacks---which exploit username-password pairs leaked from prior data
breaches---exhibit fundamentally different behavioral signatures than
traditional dictionary attacks, including higher per-attempt success
rates and lower attempt volumes. Park et al.~{[}22{]} documented SSH
attacks specifically targeting network infrastructure devices (routers,
switches), which exhibit distinct authentication patterns compared to
attacks against general-purpose servers. Traditional rule-based systems,
designed around a single attack archetype, lack the flexibility to
accommodate this diversity {[}7, 12{]}. Machine learning approaches that
construct behavioral profiles from multiple feature dimensions offer a
principled mechanism for detecting diverse attack variants without
requiring explicit rules for each {[}8, 9{]}.

\textbf{Gap 4: Disconnect between algorithmic research and deployable
systems.} A substantial body of research has demonstrated the
effectiveness of machine learning algorithms for SSH-related anomaly
detection {[}20, 21, 25{]}. However, the overwhelming majority of these
studies evaluate algorithms in isolation on benchmark datasets without
addressing the engineering challenges of integration into operational
infrastructure---log ingestion pipelines, real-time feature extraction,
automated response mechanisms, visualization dashboards, and
containerized deployment {[}12, 15{]}. Sperotto et al.~{[}15{]} noted
that the transition from research prototype to production system remains
one of the most significant barriers to the adoption of
machine-learning-based intrusion detection. This research-to-deployment
gap means that organizations cannot readily benefit from algorithmic
advances, and the practical impact of the research remains limited.

\subsection{Research Questions}\label{research-questions}

In light of the four gaps identified above, this thesis is guided by the
following research questions:

\textbf{RQ1:} How effectively does the Isolation Forest algorithm detect
SSH brute-force attacks in a semi-supervised setting (trained
exclusively on normal behavior data), and how does its performance
compare to Local Outlier Factor (LOF) and One-Class Support Vector
Machine (OCSVM) on the same real-world honeypot dataset?

\textbf{RQ2:} To what extent does the proposed EWMA-Adaptive Percentile
dynamic thresholding mechanism improve detection performance---measured
by accuracy, F1-score, and false positive rate---compared to static
threshold approaches across varying traffic conditions?

\textbf{RQ3:} What degree of early prediction capability does the system
achieve across five distinct attack scenarios (basic brute-force,
distributed attack, slow brute-force, credential stuffing, and
dictionary attack), and which behavioral features contribute most to
early detection?

\textbf{RQ4:} Can the integrated system architecture---combining ELK
Stack for log analytics, Isolation Forest for anomaly detection,
Fail2Ban for automated response, and Docker for containerized
deployment---meet the requirements for real-time monitoring and
automated response in a practical operational environment?

\section{1.3 Research Objectives}\label{research-objectives}

This research aims to design, implement, and rigorously evaluate an
intelligent SSH brute-force attack detection and prevention system that
integrates semi-supervised anomaly detection, adaptive dynamic
thresholding, and automated response within a modern, containerized
infrastructure. The following five specific objectives operationalize
the research questions stated above.

\textbf{Objective 1: Design a comprehensive SSH behavioral feature set.}
Develop and validate a set of 14 behavioral features extracted from SSH
authentication logs over 5-minute sliding windows per source IP address
{[}8, 9{]}. The feature set encompasses four categories: frequency
features (attempt counts, failure rates), temporal features
(inter-attempt timing, session duration), authentication features
(unique usernames, password diversity), and connection features (port
entropy, geographic indicators). The feature set must capture sufficient
behavioral information to discriminate between normal SSH activity and
the five attack variants under study.

\textbf{Objective 2: Train and comparatively evaluate semi-supervised
anomaly detection models.} Implement and systematically compare three
anomaly detection algorithms---Isolation Forest {[}8, 9{]}, Local
Outlier Factor {[}10{]}, and One-Class SVM {[}11{]}---in a
semi-supervised configuration where models are trained exclusively on
normal-behavior data. Establish performance benchmarks using accuracy,
precision, recall, F1-score, and false positive rate (FPR), with target
performance thresholds of F1 \textgreater= 85\% and recall \textgreater=
95\% to ensure high detection rates with acceptable false alarm levels
{[}20, 21{]}.

\textbf{Objective 3: Develop and validate an EWMA-Adaptive Percentile
dynamic thresholding mechanism.} Design a hybrid dynamic threshold that
combines the trend-tracking properties of Exponentially Weighted Moving
Average (EWMA) {[}13{]} with the distribution-free adaptiveness of
percentile-based methods {[}14{]}. The mechanism must automatically
adjust the anomaly detection threshold in response to temporal
variations in SSH traffic patterns, reducing false positive rates during
high-activity periods while maintaining sensitivity during low-activity
periods.

\textbf{Objective 4: Build a complete, containerized, end-to-end
system.} Design and implement an integrated system architecture
comprising nine Docker services: ELK Stack (Elasticsearch, Logstash,
Kibana) for log ingestion, storage, and visualization; a FastAPI backend
for model serving and API endpoints; a React frontend for security
monitoring dashboards; and Fail2Ban for automated IP blocking upon
attack detection {[}6, 15{]}. The architecture must support real-time
processing of SSH authentication events with end-to-end latency suitable
for operational deployment.

\textbf{Objective 5: Evaluate system performance against five realistic
attack scenarios.} Design and execute five simulated SSH brute-force
attack scenarios---basic brute-force, distributed attack, slow
brute-force, credential stuffing, and dictionary attack---that
collectively test the system's detection capability, early prediction
ability, and automated response across the full spectrum of modern
attack techniques {[}5, 18, 19, 23{]}. Each scenario targets a specific
system capability: basic and dictionary scenarios test standard
detection, distributed and credential-stuffing scenarios test resilience
to per-IP evasion, and slow brute-force tests early-warning
effectiveness.

{[}Table 1.1: Mapping of research questions to objectives and evaluation
criteria{]}

\section{1.4 Significance of the Study}\label{significance-of-the-study}

\subsection{1.4.1 Practical Significance}\label{practical-significance}

The system developed in this research provides a directly deployable
solution for organizations seeking to enhance their SSH security posture
beyond the capabilities of traditional tools. The Docker-based
containerized architecture enables deployment on any Linux server with
minimal configuration, while the use of exclusively open-source
components (scikit-learn, FastAPI, Elasticsearch, Kibana, Fail2Ban)
eliminates licensing costs that would otherwise be associated with
commercial Security Information and Event Management (SIEM) solutions
{[}6, 15{]}. The system's real-time monitoring capability, delivered
through both Kibana dashboards and a custom React-based interface,
provides security administrators with actionable visibility into SSH
authentication patterns. Integration with Fail2Ban ensures that detected
attacks trigger automated IP blocking, reducing mean time to respond
(MTTR) from minutes (manual intervention) to seconds (automated
response).

The early-warning capability represents a particularly significant
practical advance. Rather than waiting for an attacker to execute
hundreds or thousands of login attempts before triggering a static
threshold, the system can identify attack intent from behavioral
features extracted during the initial probing phase {[}5, 23{]}. In the
slow brute-force scenario, the system issues an EARLY\_WARNING alert
after only 3--5 attempts (approximately 2--3 minutes into the attack),
whereas Fail2Ban with default settings would fail to detect the attack
entirely if the attacker maintains a sufficiently low attempt rate. This
proactive detection capability provides defenders with a critical time
advantage for investigation and response.

\subsection{1.4.2 Research Significance}\label{research-significance}

This thesis contributes to the academic literature in three distinct
areas. First, it provides a rigorous comparative evaluation of three
anomaly detection algorithms (IF, LOF, OCSVM) on real-world SSH honeypot
data in a semi-supervised configuration {[}8, 9, 10, 11{]}. While these
algorithms have been extensively studied on benchmark datasets such as
NSL-KDD and CICIDS2017, their comparative performance on authentic SSH
attack data---with its characteristic noise, class imbalance, and
behavioral diversity---remains underexplored {[}24, 25{]}. Second, the
proposed EWMA-Adaptive Percentile hybrid thresholding mechanism {[}13,
14{]} represents a novel contribution to the literature on dynamic
thresholds for anomaly detection, combining the complementary strengths
of trend-following (EWMA) and distribution-free (percentile) approaches.
Third, the finding that temporal features---specifically
session\_duration\_mean (importance: 5.50\%), min\_inter\_attempt\_time
(3.86\%), and mean\_inter\_attempt\_time (2.61\%)---outperform
traditional frequency-based features (e.g., fail\_count) as
discriminators between normal and attack behavior has implications for
the design of future SSH monitoring systems {[}8, 15{]}.

\subsection{1.4.3 Empirical Significance}\label{empirical-significance}

The dataset constructed for this research---174,250 log lines combining
119,729 lines of real honeypot attack data (679 IPs, 29,301 failed
passwords, 532 accepted root logins) with 54,521 lines of simulated
normal behavior (64 users)---provides an empirical foundation that is
both more realistic and more comprehensive than the benchmark datasets
commonly used in the literature {[}24{]}. The honeypot data captures the
full behavioral diversity of real-world SSH attackers, including
geographic distribution, temporal patterns, username preferences, and
attack-speed profiles, while the simulated normal data provides
controlled ground truth for model training. The 14 behavioral features
extracted per IP per 5-minute window, the feature importance rankings,
and the comparative algorithm performance metrics constitute empirical
contributions that can inform and guide future research in this domain
{[}20, 21{]}.

\section{1.5 Scope and Limitations}\label{scope-and-limitations}

\subsection{1.5.1 Scope}\label{scope}

The scope of this research is defined along four dimensions.
\textbf{Protocol and attack type:} The study focuses exclusively on SSH
version 2 (SSH-2) and password-based brute-force attacks in five
variants: basic brute-force, dictionary attack, slow brute-force,
distributed attack, and credential stuffing {[}1, 2{]}. Other SSH-based
attacks (man-in-the-middle, session hijacking, protocol-level exploits)
fall outside the scope. \textbf{Data:} The research employs two data
sources: (i) real attack data collected from an SSH honeypot deployed on
a production VPS, comprising 119,729 log lines from 679 unique IP
addresses with 29,301 failed password attempts and 532 accepted root
logins; and (ii) simulated normal-behavior data comprising 54,521 log
lines from 64 user accounts performing routine SSH operations. The
combined dataset totals 174,250 log lines {[}3, 24{]}.
\textbf{Algorithms:} Three semi-supervised anomaly detection algorithms
are implemented and evaluated: Isolation Forest (primary) {[}8, 9{]},
LOF (benchmark) {[}10{]}, and OCSVM (benchmark) {[}11{]}. Supervised and
deep learning approaches are excluded by design, as the semi-supervised
paradigm addresses the practical constraint of unavailable labeled
attack data. \textbf{Infrastructure:} The system is built on ELK Stack
(Elasticsearch 8.x, Logstash, Kibana), FastAPI, React, Fail2Ban, and
Docker Compose (9 services), deployed on a Linux-based server
environment.

\subsection{1.5.2 Limitations}\label{limitations}

Four limitations must be acknowledged. First, the normal-behavior data
was generated through simulation rather than collected from a production
environment; while the simulation covers a broad range of legitimate SSH
usage patterns (64 user profiles), certain environment-specific
behaviors may be underrepresented {[}24{]}. Second, the system has been
evaluated at moderate scale; scalability to environments processing
authentication events from hundreds of simultaneous SSH servers has not
been empirically verified, although the ELK Stack architecture supports
horizontal scaling {[}15{]}. Third, the Isolation Forest model was
trained on data from a specific VPS environment and may require
retraining when deployed to environments with substantially different
baseline SSH usage patterns {[}8, 9{]}. Fourth, the false positive rate
of 29.00\% for Isolation Forest, while acceptable for a semi-supervised
system operating without labeled attack training data, indicates room
for improvement through feature refinement, ensemble methods, or
incorporation of limited labeled data in future work {[}20, 21{]}.
Goldstein and Uchida {[}20{]} observed similar false positive rates for
unsupervised methods across multiple benchmark datasets, suggesting that
this is a characteristic limitation of the one-class learning paradigm
rather than a deficiency specific to this implementation.

\section{1.6 Thesis Structure}\label{thesis-structure}

This thesis is organized into six chapters, each addressing a distinct
phase of the research process:

\textbf{Chapter 1: Introduction.} Establishes the research context,
articulates the four-gap problem statement, formulates four research
questions (RQ1--RQ4), specifies five research objectives, and delineates
the scope and limitations of the study.

\textbf{Chapter 2: Literature Review.} Provides a comprehensive
synthesis of the theoretical foundations and related work, including:
the SSH protocol and its security landscape {[}1, 2{]}; brute-force
attack taxonomy and evolution {[}3, 4, 5{]}; traditional defense
mechanisms and their limitations {[}6, 7{]}; anomaly detection paradigms
{[}7, 12{]}; the mathematical foundations of Isolation Forest {[}8,
9{]}, LOF {[}10{]}, and OCSVM {[}11{]}; dynamic thresholding mechanisms
{[}13, 14{]}; and a systematic comparison of related studies. The
chapter concludes with the identification of research gaps and the
articulation of this thesis's contributions.

\textbf{Chapter 3: Methodology.} Describes the research methodology,
including system architecture design, data collection and preprocessing
procedures, the 14-feature engineering pipeline, model configuration and
training protocols, the EWMA-Adaptive Percentile dynamic threshold
mechanism, and the five attack scenario simulation designs.

\textbf{Chapter 4: Experimental Results.} Presents the experimental
outcomes, including dataset statistics, model training and
hyperparameter optimization results, comparative algorithm performance,
dynamic threshold evaluation, feature importance analysis, and attack
scenario test results.

\textbf{Chapter 5: Discussion.} Interprets the experimental results in
the context of the research questions, compares findings with related
work, analyzes the significance of temporal features as top
discriminators, and discusses the practical implications of the
integrated system architecture.

\textbf{Chapter 6: Conclusion and Future Work.} Summarizes the principal
contributions and findings, evaluates the degree to which each research
objective was achieved, and proposes directions for future research
including deep learning extensions, federated learning for multi-site
deployment, and adversarial robustness evaluation.

{[}Figure 1.1: Global trends in SSH brute-force attacks and the
evolution of defense mechanisms (2015--2025){]}

{[}Figure 1.2: Conceptual comparison of static threshold versus dynamic
threshold detection approaches{]}

{[}Figure 1.3: High-level architecture of the proposed system showing
the nine Docker services and data flow{]}

{[}Table 1.1: Mapping of research questions to objectives and evaluation
criteria{]}

{[}Table 1.2: Summary of research gaps and corresponding contributions
of this thesis{]}

\chapter{CHAPTER 2: LITERATURE
REVIEW}\label{chapter-2-literature-review}

\section{2.1 Secure Shell (SSH) Protocol and Its Security
Landscape}\label{secure-shell-ssh-protocol-and-its-security-landscape}

The Secure Shell (SSH) protocol was originally developed in 1995 by Tatu
Ylonen at the Helsinki University of Technology as a secure replacement
for unencrypted remote access protocols including Telnet, rlogin, and
rsh {[}1{]}. The protocol was subsequently standardized by the Internet
Engineering Task Force (IETF) through the RFC 4250--4256 document
series, with RFC 4251 defining the overall protocol architecture
{[}1{]}. Barrett, Silverman, and Byrnes {[}2{]} provide a comprehensive
treatment of SSH implementation and deployment, documenting the
protocol's evolution from a research tool to the universal standard for
secure remote system administration. As of 2024, SSH is deployed on
virtually every Linux and Unix server connected to the Internet, and all
major cloud infrastructure providers---Amazon Web Services, Google Cloud
Platform, and Microsoft Azure---use SSH key pairs as the default
mechanism for initial access to virtual machine instances {[}1, 2{]}.

The SSH-2 protocol architecture is organized into three layered
components, each addressing a distinct security function {[}1, 2{]}. The
Transport Layer Protocol (RFC 4253) establishes server authentication,
data confidentiality through symmetric encryption, and data integrity
through message authentication codes. The handshake process involves
protocol version exchange, algorithm negotiation, Diffie-Hellman key
exchange, and server host key verification. The User Authentication
Protocol (RFC 4252) handles client identity verification through one of
several methods: password authentication, public key authentication,
host-based authentication, or keyboard-interactive authentication. The
Connection Protocol (RFC 4254) multiplexes multiple logical channels
over a single encrypted connection, supporting remote command execution,
port forwarding, and secure file transfer via SFTP/SCP.

Password-based authentication, despite being the least secure of the
available methods, remains widely deployed due to its simplicity and the
absence of prerequisite key management infrastructure {[}2, 4{]}. The
default OpenSSH configuration permits up to six authentication attempts
per connection (MaxAuthTries = 6) and imposes no limit on the number of
concurrent connections from a single source IP address {[}1, 6{]}. This
permissive default configuration, combined with the protocol's lack of
any built-in anti-automation mechanism (unlike web applications that can
employ CAPTCHAs), creates a structurally exploitable attack surface for
automated brute-force tools {[}2, 5{]}.

The scale of SSH-targeted attacks is documented by multiple empirical
studies. Wu et al.~{[}3{]} analyzed data from the National Center for
Supercomputing Applications (NCSA) honeypot deployment and recorded
approximately 11 billion SSH brute-force login attempts, demonstrating
that SSH brute-forcing operates at an industrial scale driven by
globally distributed botnets. The Verizon Data Breach Investigations
Report {[}4{]} identifies stolen credentials as the most prevalent
initial attack vector, accounting for approximately 29\% of confirmed
data breaches, with brute-force and credential-stuffing attacks
representing the primary credential-theft mechanisms. Hellemons et
al.~{[}5{]} developed the SSHCure system and characterized SSH attacks
as following a three-phase temporal structure: an initial scanning phase
in which the attacker probes target reachability, a brute-force phase
involving systematic credential enumeration, and a compromise phase in
which successful credentials are exploited for lateral movement or data
exfiltration.

SSH servers record authentication events in system log files (typically
\texttt{/var/log/auth.log} on Debian-based systems or
\texttt{/var/log/secure} on Red Hat-based systems), with each event
containing a timestamp, source IP address, attempted username,
authentication result (success or failure), authentication method, and
source port {[}1, 5{]}. These structured log fields provide the raw data
from which behavioral features can be extracted for
machine-learning-based anomaly detection, a capability that this thesis
exploits through a 14-feature extraction pipeline operating on 5-minute
sliding windows per source IP address.

{[}Figure 2.1: Layered architecture of the SSH-2 protocol and
authentication flow{]}

\section{2.2 Brute-Force Attacks: Taxonomy, Patterns, and
Evolution}\label{brute-force-attacks-taxonomy-patterns-and-evolution}

Brute-force attacks against SSH servers can be formally defined as
systematic attempts to discover valid authentication credentials through
exhaustive or guided enumeration of the credential space {[}3, 4{]}.
Given a character alphabet of size \(|A|\) and a target password of
length \(L\), the theoretical maximum search space is \(S = |A|^L\). For
a password composed of lowercase letters (26) and digits (10) with
length 8, this yields \(S = 36^8 \approx 2.82 \times 10^{12}\)
combinations. In practice, however, attackers rarely perform exhaustive
enumeration; instead, they exploit the highly non-uniform distribution
of human-chosen passwords to reduce the effective search space by orders
of magnitude {[}4, 19{]}.

The taxonomy of SSH brute-force attacks has evolved substantially over
the past decade, driven by the arms race between attackers and
defenders. Based on the literature {[}3, 5, 18, 19, 23{]} and empirical
analysis of the honeypot data collected for this research (679 unique
attacking IPs, 29,301 failed password attempts), five principal attack
variants can be identified:

\textbf{Classic brute-force.} The attacker attempts credential
combinations at maximum speed, typically targeting high-value accounts
such as root, admin, and user {[}3, 5{]}. Characteristic behavioral
signatures include extremely high login attempt frequency (hundreds to
thousands per minute), near-uniform inter-attempt intervals (reflecting
automated tool cadence), and rapid cycling through multiple usernames.
Wu et al.~{[}3{]} found that the root account is targeted in over 80\%
of SSH brute-force sessions observed in their honeypot data.

\textbf{Dictionary attack.} This variant employs curated password lists
(wordlists) rather than exhaustive enumeration, exploiting the
observation that a small number of passwords account for a
disproportionate fraction of real-world credential choices {[}19{]}.
Owens and Matthews {[}19{]} analyzed SSH honeypot data and found that
the top 20 most-attempted passwords accounted for over 40\% of all login
attempts, with ``123456'', ``password'', and ``admin'' consistently
appearing in the top positions. Widely used wordlists include RockYou
(approximately 14 million entries) and the SecLists collection, and
attackers increasingly generate target-specific lists incorporating
organizational names, service identifiers, and locale-specific terms.

\textbf{Slow brute-force (low-and-slow attack).} The attacker
deliberately reduces the attempt rate to evade frequency-based detection
thresholds, executing only a few attempts per minute or per hour
{[}23{]}. Javed and Paxson {[}23{]} conducted a systematic study of
stealthy SSH brute-forcing and demonstrated that attacks operating at
rates as low as 1--2 attempts per hour can achieve compromise rates
comparable to aggressive campaigns while remaining invisible to
conventional monitoring tools configured with standard thresholds. This
attack variant is particularly insidious because its behavioral profile
closely resembles that of a legitimate user who has forgotten their
password {[}5, 7{]}.

\textbf{Distributed brute-force.} This variant leverages botnets or
proxy networks to distribute the attack across hundreds or thousands of
source IP addresses, with each individual IP performing only a small
number of attempts {[}3, 5{]}. Per-IP failure-count-based detection
methods are structurally incapable of identifying distributed attacks
because no single IP exceeds the detection threshold. Wu et al.~{[}3{]}
documented coordinated distributed attacks originating from
geographically diverse IP addresses that collectively executed millions
of attempts while maintaining per-IP attempt counts below typical
Fail2Ban thresholds.

\textbf{Credential stuffing.} This variant exploits username-password
pairs leaked from data breaches at other services, capitalizing on the
widespread practice of password reuse across platforms {[}4, 19{]}.
Credential-stuffing attacks are particularly dangerous because the
per-attempt success rate is substantially higher than that of random or
dictionary-based attacks. The Verizon DBIR {[}4{]} notes that
credential-stuffing attacks have increased significantly following major
data breaches, with automated tools such as Sentry MBA and OpenBullet
enabling attackers to test millions of leaked credentials against SSH
servers.

Bezerra et al.~{[}18{]} analyzed SSH attack durations across a
large-scale honeypot dataset and identified significant variability in
temporal profiles: attack sessions ranged from sub-second bursts
(automated scanners testing a single credential) to sustained campaigns
spanning multiple days (coordinated, multi-phase attacks). This temporal
diversity imposes a requirement for detection systems to operate
effectively across multiple time scales---a requirement that motivates
the 5-minute sliding window approach adopted in this thesis.

{[}Table 2.1: Taxonomy of SSH brute-force attack variants with
characteristic behavioral features{]}

{\def\LTcaptype{none} % do not increment counter
\begin{longtable}[]{@{}
  >{\raggedright\arraybackslash}p{(\linewidth - 8\tabcolsep) * \real{0.1200}}
  >{\raggedright\arraybackslash}p{(\linewidth - 8\tabcolsep) * \real{0.1733}}
  >{\raggedright\arraybackslash}p{(\linewidth - 8\tabcolsep) * \real{0.1733}}
  >{\raggedright\arraybackslash}p{(\linewidth - 8\tabcolsep) * \real{0.2533}}
  >{\raggedright\arraybackslash}p{(\linewidth - 8\tabcolsep) * \real{0.2800}}@{}}
\toprule\noalign{}
\begin{minipage}[b]{\linewidth}\raggedright
Variant
\end{minipage} & \begin{minipage}[b]{\linewidth}\raggedright
Attempt Rate
\end{minipage} & \begin{minipage}[b]{\linewidth}\raggedright
IP Diversity
\end{minipage} & \begin{minipage}[b]{\linewidth}\raggedright
Password Strategy
\end{minipage} & \begin{minipage}[b]{\linewidth}\raggedright
Detection Difficulty
\end{minipage} \\
\midrule\noalign{}
\endhead
\bottomrule\noalign{}
\endlastfoot
Classic brute-force & Very high (\textgreater100/min) & Single IP &
Exhaustive or top-N & Low \\
Dictionary attack & High (10--100/min) & Single IP & Wordlist-based &
Low--Medium \\
Slow brute-force & Very low (\textless1/min) & Single IP & Targeted &
High \\
Distributed attack & Low per IP & Many IPs (\textgreater10) &
Coordinated & High \\
Credential stuffing & Medium & Variable & Leaked credentials & Very
High \\
\end{longtable}
}

\section{2.3 Traditional Defense Mechanisms and Their
Limitations}\label{traditional-defense-mechanisms-and-their-limitations}

Traditional SSH brute-force defense mechanisms can be categorized into
four classes, each with characteristic strengths and limitations {[}6,
7, 12{]}.

\textbf{Threshold-based intrusion prevention.} Fail2Ban {[}6{]} is the
most widely deployed tool in this category. It monitors SSH
authentication logs and applies firewall rules (via iptables or
nftables) to block IP addresses that exceed a configured failure
threshold within a specified time window. The default SSH jail
configuration uses maxretry = 5, findtime = 600 seconds, and bantime =
600 seconds. While effective against naive high-speed attacks,
Fail2Ban's static threshold architecture creates an inherent
precision-sensitivity tradeoff: lowering the threshold increases false
positives (blocking legitimate users who mistype passwords), while
raising it increases false negatives (permitting slow or distributed
attacks) {[}6, 7{]}. Chandola et al.~{[}7{]} identified this
static-threshold limitation as a fundamental challenge in anomaly
detection, noting that real-world data streams exhibit non-stationary
behavior that fixed thresholds cannot accommodate.

\textbf{IP reputation and blacklisting.} Services such as AbuseIPDB,
Spamhaus, and Blocklist.de maintain databases of IP addresses associated
with malicious activity {[}4, 6{]}. Blacklist-based approaches are
inherently reactive---an IP is listed only after it has been observed
conducting attacks elsewhere---and are easily circumvented by attackers
who rotate through fresh IP addresses or employ residential proxy
networks {[}3, 7{]}. Furthermore, blacklists suffer from both false
positives (legitimate services operating behind shared NAT addresses
that have been erroneously listed) and incomplete coverage (newly
provisioned attack infrastructure that has not yet been reported).

\textbf{Signature-based intrusion detection.} Network intrusion
detection systems such as Snort and Suricata can identify SSH
brute-force attacks by matching traffic patterns against predefined rule
signatures {[}7, 15{]}. However, signature-based detection is
structurally limited to known attack patterns for which rules have been
written and cannot detect novel attack variants {[}12{]}. Moreover,
because SSH traffic is encrypted end-to-end, deep packet inspection of
SSH session content is infeasible; detection must rely on metadata
features (connection frequency, flow duration, packet sizes) rather than
payload analysis {[}15{]}.

\textbf{Authentication hardening.} Measures such as public key
authentication, two-factor authentication (via PAM modules),
non-standard port assignment, and port knocking reduce the attack
surface by either eliminating password authentication entirely or
raising the barrier to initial connection establishment {[}1, 2{]}.
While highly effective when deployed, these measures are not universally
applicable---legacy systems, multi-user environments with diverse client
configurations, and cloud instances requiring initial password-based
bootstrapping all present practical obstacles to full authentication
hardening {[}2, 6{]}.

The limitations of these traditional mechanisms motivate the application
of machine learning approaches that can learn complex behavioral
patterns from data, adapt to changing traffic conditions, and detect
novel attack variants without requiring explicit rule definitions {[}7,
12{]}. Buczak and Guven {[}12{]} concluded their comprehensive survey by
noting that the integration of machine learning with traditional
security infrastructure represents the most promising direction for
next-generation intrusion detection systems.

{[}Table 2.2: Comparison of traditional SSH defense mechanisms{]}

{\def\LTcaptype{none} % do not increment counter
\begin{longtable}[]{@{}
  >{\raggedright\arraybackslash}p{(\linewidth - 6\tabcolsep) * \real{0.1569}}
  >{\raggedright\arraybackslash}p{(\linewidth - 6\tabcolsep) * \real{0.2157}}
  >{\raggedright\arraybackslash}p{(\linewidth - 6\tabcolsep) * \real{0.2549}}
  >{\raggedright\arraybackslash}p{(\linewidth - 6\tabcolsep) * \real{0.3725}}@{}}
\toprule\noalign{}
\begin{minipage}[b]{\linewidth}\raggedright
Method
\end{minipage} & \begin{minipage}[b]{\linewidth}\raggedright
Strengths
\end{minipage} & \begin{minipage}[b]{\linewidth}\raggedright
Limitations
\end{minipage} & \begin{minipage}[b]{\linewidth}\raggedright
Evasion Techniques
\end{minipage} \\
\midrule\noalign{}
\endhead
\bottomrule\noalign{}
\endlastfoot
Fail2Ban (threshold) & Simple, low overhead & Static threshold, no
adaptation & Slow-rate, distributed \\
IP blacklisting & No local computation & Reactive, incomplete coverage &
IP rotation, proxies \\
Signature IDS & Known-attack detection & Cannot detect novel attacks &
Encryption, polymorphism \\
Auth hardening & Eliminates attack surface & Not universally deployable
& N/A (preventive) \\
\end{longtable}
}

\section{2.4 Anomaly Detection for Intrusion Detection
Systems}\label{anomaly-detection-for-intrusion-detection-systems}

Anomaly detection, defined as the identification of data patterns that
deviate significantly from expected normal behavior, constitutes one of
the foundational paradigms in intrusion detection {[}7{]}. Chandola,
Banerjee, and Kumar {[}7{]} provided a seminal taxonomy of anomaly
detection approaches, categorizing them along two principal dimensions:
the nature of the detection model (classification-based,
nearest-neighbor-based, clustering-based, statistical,
information-theoretic, spectral, and isolation-based) and the
availability of labeled data (supervised, semi-supervised, and
unsupervised).

In the context of network intrusion detection, the semi-supervised
anomaly detection paradigm---in which models are trained exclusively on
data representing normal behavior and subsequently identify deviations
from the learned normal profile as potential attacks---has attracted
substantial research attention {[}7, 12, 17{]}. This paradigm offers
three key advantages for SSH brute-force detection. First, it does not
require labeled attack data for training, circumventing the practical
challenge of obtaining comprehensive, accurately labeled datasets in
operational environments {[}17, 24{]}. Pimentel et al.~{[}17{]} provided
a comprehensive review of novelty detection---the closely related
problem of identifying previously unseen patterns---and noted that
one-class learning approaches (trained on a single class of normal data)
provide the most robust framework for detecting genuinely novel
anomalies. Second, semi-supervised approaches can detect zero-day
attacks and novel attack variants that have never been observed in
training data, because detection is based on deviation from normality
rather than similarity to known attacks {[}7, 12{]}. Third, in SSH
environments, normal behavior data is readily available from production
logs (by selecting periods of known-clean operation) or from controlled
simulation, making the semi-supervised training paradigm practically
feasible {[}15, 24{]}.

The choice among specific anomaly detection algorithms involves
tradeoffs along multiple dimensions: computational complexity (critical
for real-time processing), sensitivity to hyperparameter selection,
ability to handle high-dimensional feature spaces, robustness to noise
and outliers in training data, and interpretability of anomaly scores
{[}7, 20{]}. Goldstein and Uchida {[}20{]} conducted a comprehensive
comparative evaluation of unsupervised anomaly detection algorithms on
standardized benchmark datasets and found that no single algorithm
dominates across all datasets and evaluation metrics, underscoring the
importance of algorithm selection based on the specific characteristics
of the target domain. Nassif et al.~{[}21{]} extended this analysis to
the cybersecurity domain specifically, surveying machine learning
approaches for anomaly-based intrusion detection and identifying
Isolation Forest, LOF, and OCSVM as the three most widely studied
unsupervised algorithms for network security applications.

For this thesis, three algorithms were selected for implementation and
comparative evaluation: Isolation Forest {[}8, 9{]} as the primary
detection algorithm, and Local Outlier Factor {[}10{]} and One-Class SVM
{[}11{]} as benchmarks. This selection is motivated by the complementary
detection paradigms they represent (isolation-based, density-based, and
boundary-based, respectively), their established track records in
network intrusion detection research {[}20, 21{]}, and their
availability in mature, well-tested implementations within the
scikit-learn library. The following three sections present the
mathematical foundations of each algorithm.

\section{2.5 Isolation Forest and Its
Extensions}\label{isolation-forest-and-its-extensions}

\subsection{2.5.1 Algorithmic Foundation}\label{algorithmic-foundation}

Isolation Forest (IF) was proposed by Liu, Ting, and Zhou {[}8{]} at the
IEEE International Conference on Data Mining (ICDM) in 2008 and
subsequently formalized in a comprehensive journal publication in ACM
Transactions on Knowledge Discovery from Data (TKDD) in 2012 {[}9{]}.
The algorithm represents a paradigm shift in anomaly detection: whereas
traditional methods detect anomalies as data points that are distant
from (distance-based methods) or sparse relative to (density-based
methods) the majority of the data, Isolation Forest detects anomalies as
data points that are easy to isolate through random partitioning {[}8,
9{]}.

The fundamental insight underlying Isolation Forest is that anomalous
data points, by virtue of their atypical feature values, require fewer
random partitions to be separated from the rest of the dataset than
normal points {[}8{]}. This property arises because anomalies are
typically few in number and possess feature values that are
substantially different from those of the majority population, making
them susceptible to isolation by a small number of random splits.

\subsection{2.5.2 Isolation Tree
Construction}\label{isolation-tree-construction}

An Isolation Tree (iTree) is constructed through a recursive random
partitioning process {[}8, 9{]}. Given a dataset
\(X = \{x_1, x_2, \ldots, x_n\}\) with \(d\) features, the construction
proceeds as follows:

\begin{enumerate}
\def\labelenumi{\arabic{enumi}.}
\tightlist
\item
  Randomly select a feature \(q\) from the \(d\) available features.
\item
  Randomly select a split value \(p\) uniformly from the interval
  \([\min(X_q), \max(X_q)]\), where \(X_q\) denotes the values of
  feature \(q\) in the current node's data.
\item
  Partition the data into two subsets: the left branch contains points
  where \(x_q < p\), and the right branch contains points where
  \(x_q \geq p\).
\item
  Recurse on each branch until one of the stopping conditions is met:
  the node contains a single data point (isolated), or the tree reaches
  a predefined maximum depth \(l = \lceil \log_2 \psi \rceil\), where
  \(\psi\) is the sub-sampling size.
\end{enumerate}

\subsection{2.5.3 Anomaly Score
Computation}\label{anomaly-score-computation}

The anomaly score of a data point \(x\) is derived from its average path
length across an ensemble of \(t\) Isolation Trees {[}8, 9{]}. The path
length \(h(x)\) of a point \(x\) in a single iTree is defined as the
number of edges traversed from the root node to the terminal node
(external node) at which \(x\) is isolated, plus an adjustment factor
\(c(k)\) for the unbuilt subtree when the maximum depth is reached,
where \(k\) is the number of data points in the terminal node.

The anomaly score is defined as:

\[s(x, n) = 2^{-\frac{E[h(x)]}{c(n)}}\]

where \(E[h(x)]\) is the expected (average) path length of \(x\) across
\(t\) Isolation Trees, and \(c(n)\) is the average path length of an
unsuccessful search in a Binary Search Tree (BST) constructed from \(n\)
data points, serving as a normalization factor {[}8, 9{]}. The
normalization factor is computed as:

\[c(n) = 2H(n-1) - \frac{2(n-1)}{n}\]

where \(H(i) = \ln(i) + \gamma\) is the harmonic number approximation
and \(\gamma \approx 0.5772\) is the Euler-Mascheroni constant {[}8{]}.

The anomaly score has the following interpretation {[}8, 9{]}: -
\(s(x, n) \to 1\): the point has a short average path length, indicating
high susceptibility to isolation and therefore high anomaly likelihood.
- \(s(x, n) \approx 0.5\): the average path length is close to the
expected path length for the dataset, indicating no clear anomaly
signal. - \(s(x, n) \to 0\): the point has a long average path length,
indicating deep embedding within the normal data distribution.

\subsection{2.5.4 Computational Complexity and Practical
Advantages}\label{computational-complexity-and-practical-advantages}

The time complexity of Isolation Forest is
\(O(t \cdot \psi \cdot \log \psi)\) for training and
\(O(t \cdot \log \psi)\) for scoring a single data point, where \(t\) is
the number of trees and \(\psi\) is the sub-sampling size {[}8, 9{]}.
This linear-logarithmic complexity represents a significant advantage
over density-based methods (LOF: \(O(n^2)\)) and kernel-based methods
(OCSVM: \(O(n^2)\) to \(O(n^3)\)), making Isolation Forest particularly
suitable for real-time applications where scoring latency is a critical
constraint {[}8, 20{]}.

Additional practical advantages include: (i) robustness to the swamping
and masking effects that degrade the performance of distance-based and
density-based methods when anomalies are numerous or clustered {[}9{]};
(ii) absence of distributional assumptions about the data, unlike
statistical methods that assume Gaussian or other parametric
distributions {[}8, 17{]}; (iii) effective handling of high-dimensional
feature spaces without suffering from the curse of dimensionality that
affects distance-based methods {[}9, 20{]}; and (iv) interpretability of
the anomaly score, which has a natural probabilistic interpretation as
the likelihood of isolation under random partitioning {[}8{]}.

\subsection{2.5.5 Extended Isolation
Forest}\label{extended-isolation-forest}

Hariri, Kind, and Brunner {[}16{]} identified a limitation of the
original Isolation Forest algorithm: because axis-parallel splits are
used exclusively, the algorithm can produce artifacts in the anomaly
score landscape, assigning anomalous scores to normal points that lie
along the coordinate axes in regions of feature space dominated by
anomalies. The Extended Isolation Forest (EIF) addresses this limitation
by replacing axis-parallel splits with hyperplane splits of arbitrary
orientation, defined by a random normal vector and an intercept point
{[}16{]}. While EIF provides more accurate anomaly scores in datasets
with complex geometric structures, the original Isolation Forest remains
the more widely deployed variant for intrusion detection due to its
simplicity, interpretability, and lower computational overhead {[}16,
20{]}. This thesis employs the original Isolation Forest algorithm for
the primary detection model.

{[}Figure 2.2: Illustration of Isolation Forest partitioning: anomalous
points (short path lengths) versus normal points (long path lengths){]}

\section{2.6 Local Outlier Factor (LOF)}\label{local-outlier-factor-lof}

\subsection{2.6.1 Algorithmic
Foundation}\label{algorithmic-foundation-1}

Local Outlier Factor (LOF) was proposed by Breunig, Kriegel, Ng, and
Sander {[}10{]} at the ACM SIGMOD International Conference on Management
of Data in 2000. LOF introduced the concept of local density-based
anomaly detection, addressing a fundamental limitation of global
approaches: a data point may be anomalous relative to its local
neighborhood while appearing normal in a global context, or vice versa
{[}10{]}. By comparing the local density of each data point with the
local densities of its \(k\)-nearest neighbors, LOF produces a
degree-of-outlierness score that captures the relative isolation of a
point within its local context.

\subsection{2.6.2 Mathematical
Formulation}\label{mathematical-formulation}

The LOF computation involves several intermediate definitions {[}10{]}.
Given a dataset \(X\) and a positive integer \(k\):

\textbf{\(k\)-distance.} The \(k\)-distance of a point \(x\), denoted
\(d_k(x)\), is the distance between \(x\) and its \(k\)-th nearest
neighbor in \(X\).

\textbf{\(k\)-distance neighborhood.} The \(k\)-distance neighborhood of
\(x\), denoted \(N_k(x)\), is the set of all points whose distance from
\(x\) is at most \(d_k(x)\):
\(N_k(x) = \{y \in X \setminus \{x\} : d(x, y) \leq d_k(x)\}\).

\textbf{Reachability distance.} The reachability distance of \(x\) with
respect to a point \(o\) is:

\[\text{reach-dist}_k(x, o) = \max\{d_k(o), d(x, o)\}\]

This definition smooths the density estimation by replacing distances
smaller than \(d_k(o)\) with \(d_k(o)\), reducing the sensitivity of the
density estimate to statistical fluctuations among the nearest neighbors
{[}10{]}.

\textbf{Local reachability density.} The local reachability density of
\(x\) is the inverse of the average reachability distance from \(x\) to
its \(k\)-nearest neighbors:

\[\text{lrd}_k(x) = \left( \frac{\sum_{o \in N_k(x)} \text{reach-dist}_k(x, o)}{|N_k(x)|} \right)^{-1}\]

\textbf{Local Outlier Factor.} The LOF score of \(x\) is the average
ratio of the local reachability densities of \(x\)'s neighbors to
\(x\)'s own local reachability density:

\[\text{LOF}_k(x) = \frac{\sum_{o \in N_k(x)} \frac{\text{lrd}_k(o)}{\text{lrd}_k(x)}}{|N_k(x)|}\]

Interpretation: \(\text{LOF}_k(x) \approx 1\) indicates that \(x\) has a
local density similar to its neighbors (normal).
\(\text{LOF}_k(x) \gg 1\) indicates that \(x\) has a significantly lower
local density than its neighbors (anomalous). \(\text{LOF}_k(x) < 1\)
indicates that \(x\) is denser than its neighborhood (deeply embedded in
a cluster) {[}10{]}.

\subsection{2.6.3 Strengths and Limitations for SSH
Detection}\label{strengths-and-limitations-for-ssh-detection}

LOF excels at detecting local anomalies---points that are anomalous
relative to their immediate neighborhood but not necessarily in a global
context {[}10, 20{]}. This property is valuable for SSH intrusion
detection in environments where different user groups exhibit distinct
behavioral patterns (e.g., system administrators versus application
developers), as LOF can identify behavior that is anomalous within a
specific user group without requiring global normalization {[}10, 17{]}.

However, LOF has significant practical limitations for real-time SSH
monitoring {[}10, 20, 21{]}. The computational complexity for computing
\(k\)-nearest neighbors is \(O(n^2)\) for brute-force search (or
\(O(n \log n)\) with spatial indexing structures such as KD-trees, which
degrade in high-dimensional spaces). Additionally, LOF requires storing
the entire training dataset in memory for neighbor queries during the
scoring phase, and the algorithm's performance is sensitive to the
choice of \(k\), which may require domain-specific tuning {[}10, 20{]}.
Goldstein and Uchida {[}20{]} found that LOF's performance is
competitive with Isolation Forest on low-dimensional benchmark datasets
but degrades more rapidly as dimensionality increases.

\section{2.7 One-Class Support Vector Machine
(OCSVM)}\label{one-class-support-vector-machine-ocsvm}

\subsection{2.7.1 Algorithmic
Foundation}\label{algorithmic-foundation-2}

One-Class Support Vector Machine (OCSVM) was proposed by Scholkopf,
Platt, Shawe-Taylor, Smola, and Williamson {[}11{]} in Neural
Computation in 2001. The algorithm extends the traditional two-class SVM
framework to the one-class (novelty detection) setting by finding a
maximal-margin hyperplane in a kernel-induced feature space that
separates the training data from the origin {[}11{]}. Data points that
fall on the origin side of the hyperplane are classified as anomalous
(novel).

\subsection{2.7.2 Optimization Problem}\label{optimization-problem}

The OCSVM optimization problem is formulated as follows {[}11{]}. Given
training data \(\{x_1, x_2, \ldots, x_n\}\) drawn from the normal class,
a kernel function \(\Phi\) that maps data into a high-dimensional
feature space, and a parameter \(\nu \in (0, 1]\) that controls the
tradeoff between the margin and the fraction of training points
permitted to fall on the anomalous side of the hyperplane:

\[\min_{w, \xi, \rho} \frac{1}{2} \|w\|^2 + \frac{1}{\nu n} \sum_{i=1}^{n} \xi_i - \rho\]

subject to:

\[w \cdot \Phi(x_i) \geq \rho - \xi_i, \quad \xi_i \geq 0, \quad i = 1, \ldots, n\]

where \(w\) is the normal vector to the separating hyperplane, \(\rho\)
is the offset, and \(\xi_i\) are slack variables that permit soft-margin
violations {[}11{]}. The parameter \(\nu\) has a dual interpretation: it
is an upper bound on the fraction of training points classified as
outliers and a lower bound on the fraction of support vectors {[}11{]}.

The dual formulation, solved via quadratic programming, yields the
decision function:

\[f(x) = \text{sgn}\left( \sum_{i=1}^{n} \alpha_i K(x_i, x) - \rho \right)\]

where \(\alpha_i\) are the dual variables (Lagrange multipliers) and
\(K(x_i, x) = \Phi(x_i) \cdot \Phi(x)\) is the kernel function {[}11{]}.
The most commonly used kernel for anomaly detection is the Radial Basis
Function (RBF) kernel:

\[K(x_i, x) = \exp\left( -\gamma \|x_i - x\|^2 \right)\]

where \(\gamma > 0\) controls the kernel bandwidth. Points for which
\(f(x) < 0\) are classified as anomalous {[}11{]}.

\subsection{2.7.3 Strengths and Limitations for SSH
Detection}\label{strengths-and-limitations-for-ssh-detection-1}

OCSVM has a rigorous theoretical foundation rooted in statistical
learning theory and convex optimization, providing principled
generalization guarantees {[}11, 17{]}. The \(\nu\) parameter offers
intuitive control over the expected anomaly proportion, which can be set
based on domain knowledge about the expected attack prevalence {[}11,
20{]}. Pimentel et al.~{[}17{]} noted that OCSVM is one of the
best-studied novelty detection algorithms, with well-understood
theoretical properties.

The primary limitations of OCSVM for real-time SSH detection are
computational: the training complexity is \(O(n^2)\) to \(O(n^3)\)
depending on the solver, and the scoring complexity is \(O(n_{sv})\)
where \(n_{sv}\) is the number of support vectors, which can be a
substantial fraction of \(n\) {[}11, 20, 21{]}. Additionally, OCSVM
performance is highly sensitive to the choice of kernel function and the
hyperparameters \(\nu\) and \(\gamma\), requiring careful tuning via
cross-validation or domain-specific heuristics {[}11, 20{]}. Goldstein
and Uchida {[}20{]} observed that OCSVM achieves competitive detection
performance when properly tuned but exhibits higher variance across
hyperparameter settings compared to Isolation Forest.

{[}Table 2.3: Comparative summary of Isolation Forest, LOF, and OCSVM{]}

{\def\LTcaptype{none} % do not increment counter
\begin{longtable}[]{@{}
  >{\raggedright\arraybackslash}p{(\linewidth - 6\tabcolsep) * \real{0.1786}}
  >{\raggedright\arraybackslash}p{(\linewidth - 6\tabcolsep) * \real{0.4286}}
  >{\raggedright\arraybackslash}p{(\linewidth - 6\tabcolsep) * \real{0.1786}}
  >{\raggedright\arraybackslash}p{(\linewidth - 6\tabcolsep) * \real{0.2143}}@{}}
\toprule\noalign{}
\begin{minipage}[b]{\linewidth}\raggedright
Property
\end{minipage} & \begin{minipage}[b]{\linewidth}\raggedright
Isolation Forest {[}8, 9{]}
\end{minipage} & \begin{minipage}[b]{\linewidth}\raggedright
LOF {[}10{]}
\end{minipage} & \begin{minipage}[b]{\linewidth}\raggedright
OCSVM {[}11{]}
\end{minipage} \\
\midrule\noalign{}
\endhead
\bottomrule\noalign{}
\endlastfoot
Detection paradigm & Isolation-based & Local density-based &
Boundary-based \\
Training complexity & \(O(t \cdot \psi \log \psi)\) & \(O(n^2)\) &
\(O(n^2)\)--\(O(n^3)\) \\
Scoring complexity & \(O(t \log \psi)\) & \(O(n \cdot k)\) &
\(O(n_{sv})\) \\
Memory requirement & Model only (\(t\) trees) & Full training set &
Support vectors \\
Distributional assumptions & None & None & Implicit (via kernel) \\
Key hyperparameters & \(t\), \(\psi\), contamination & \(k\) & \(\nu\),
\(\gamma\), kernel \\
Sensitivity to dimensionality & Low & High & Moderate \\
Score interpretability & High (probability-like) & Moderate (ratio) &
Low (signed distance) \\
Real-time suitability & High & Low--Moderate & Moderate \\
\end{longtable}
}

\section{2.8 Feature Selection and Data Preprocessing for Intrusion
Detection
Systems}\label{feature-selection-and-data-preprocessing-for-intrusion-detection-systems}

The effectiveness of any machine-learning-based intrusion detection
system is fundamentally constrained by the quality and informativeness
of the features extracted from raw data {[}7, 12, 15{]}. Buczak and
Guven {[}12{]} emphasized that feature engineering---the process of
transforming raw log data into meaningful numerical representations---is
often the single most impactful component of the machine learning
pipeline for cybersecurity applications, outweighing the choice of
algorithm in many practical scenarios.

Features for SSH brute-force detection can be categorized into four
principal classes {[}5, 12, 15{]}:

\textbf{Frequency features} capture the volume and rate of
authentication activity. Examples include the total number of login
attempts, the number of failed attempts, the failure ratio
(failed/total), and the number of unique usernames attempted within a
time window {[}5, 15{]}. Frequency features are effective for detecting
high-speed brute-force attacks but provide limited discriminative power
against slow-rate attacks that operate below conventional rate
thresholds {[}23{]}.

\textbf{Temporal features} capture the timing characteristics of
authentication events. Examples include the mean, minimum, and standard
deviation of inter-attempt intervals, session duration statistics, and
the distribution of attempts across time-of-day bins {[}5, 18{]}.
Bezerra et al.~{[}18{]} demonstrated that temporal features provide
significant discriminative power for distinguishing automated attacks
from human-initiated sessions, because automated tools produce highly
regular temporal patterns (near-constant inter-attempt intervals) that
differ markedly from the irregular timing of human interactions.
Sperotto et al.~{[}15{]} confirmed this finding in their flow-based SSH
analysis, reporting that temporal features contributed more to detection
accuracy than frequency features alone.

\textbf{Authentication features} capture the credential-level
characteristics of login behavior. Examples include the number of unique
usernames attempted, the presence of root/admin login attempts, the
diversity of attempted passwords (measurable through entropy or
unique-count metrics), and the ratio of successful to total attempts
{[}5, 19{]}. Owens and Matthews {[}19{]} found that attack sessions
exhibit characteristic username and password distributions (heavily
skewed toward common accounts and weak passwords) that differ
systematically from legitimate usage patterns.

\textbf{Connection features} capture network-level properties of SSH
sessions. Examples include the number of concurrent connections, source
port entropy, geographic origin diversity, and the use of known
malicious network ranges {[}15, 24{]}. Ring et al.~{[}24{]} surveyed
datasets for network-based intrusion detection and identified
connection-level features as particularly valuable for detecting
distributed attacks where individual-IP behavioral features may appear
benign.

This thesis extracts 14 features per source IP address per 5-minute
sliding window, spanning all four feature categories. The feature set
was designed to balance comprehensiveness (capturing sufficient
behavioral information to discriminate among all five attack variants)
with computational efficiency (maintaining the real-time processing
requirement). Feature selection was informed by the literature reviewed
above and refined through empirical feature importance analysis, which
identified session\_duration\_mean (5.50\% importance),
min\_inter\_attempt\_time (3.86\%), and mean\_inter\_attempt\_time
(2.61\%) as the top three discriminative features---confirming the
primacy of temporal features reported by Bezerra et al.~{[}18{]} and
Sperotto et al.~{[}15{]}.

Data preprocessing for anomaly detection requires particular attention
to feature scaling, missing value handling, and the treatment of
categorical variables {[}7, 12{]}. Pimentel et al.~{[}17{]} noted that
distance-based and kernel-based methods (LOF, OCSVM) are highly
sensitive to feature scale, necessitating standardization (zero mean,
unit variance) or min-max normalization. Isolation Forest is
theoretically scale-invariant (because splits are defined on individual
features), but empirical evidence suggests that standardization can
improve performance when features have vastly different ranges {[}9,
20{]}.

{[}Figure 2.3: Feature extraction pipeline from SSH authentication logs
to 14-dimensional feature vectors{]}

\section{2.9 Dynamic Threshold Mechanisms for Anomaly
Detection}\label{dynamic-threshold-mechanisms-for-anomaly-detection}

\subsection{2.9.1 The Static Threshold
Problem}\label{the-static-threshold-problem}

In anomaly detection systems, the threshold determines the boundary
between normal and anomalous classifications: data points with anomaly
scores exceeding the threshold are flagged as attacks, while those below
are treated as normal {[}7, 13{]}. Static thresholds---fixed values
determined during training and applied unchanged during operation---are
simple to implement but fundamentally inappropriate for data streams
with non-stationary characteristics {[}7, 13, 14{]}. In the SSH
monitoring context, authentication traffic exhibits pronounced temporal
variability: business hours generate more legitimate logins than
off-hours, weekdays differ from weekends, and events such as system
maintenance or application deployments create legitimate but unusual
traffic spikes {[}5, 15{]}. A static threshold calibrated for average
conditions will generate excessive false positives during high-activity
periods and miss subtle attacks during low-activity periods {[}7, 14{]}.

\subsection{2.9.2 Exponentially Weighted Moving Average
(EWMA)}\label{exponentially-weighted-moving-average-ewma}

The Exponentially Weighted Moving Average is a time series smoothing
method that assigns exponentially decreasing weights to older
observations, enabling adaptive tracking of the current process level
{[}13{]}. Originally introduced by Roberts (1959) for statistical
quality control and comprehensively analyzed by Lucas and Saccucci
{[}13{]}, EWMA has been widely adopted in process monitoring, financial
time series analysis, and network anomaly detection {[}13, 14{]}.

The EWMA statistic at time \(t\) is defined as:

\[\hat{\mu}_t = \alpha \cdot x_t + (1 - \alpha) \cdot \hat{\mu}_{t-1}\]

where \(x_t\) is the observed value (anomaly score) at time \(t\),
\(\hat{\mu}_t\) is the smoothed estimate at time \(t\), \(\hat{\mu}_0\)
is initialized to the mean of the training anomaly scores, and
\(\alpha \in (0, 1]\) is the smoothing parameter that controls the
tradeoff between responsiveness to recent observations and stability
against noise {[}13{]}.

The variance of the EWMA statistic is:

\[\text{Var}(\hat{\mu}_t) = \sigma^2 \cdot \frac{\alpha}{2 - \alpha} \cdot \left[1 - (1 - \alpha)^{2t}\right]\]

which converges to \(\sigma^2 \cdot \frac{\alpha}{2 - \alpha}\) as
\(t \to \infty\) {[}13{]}. Control limits (thresholds) are typically set
at \(\hat{\mu}_t \pm L \cdot \sqrt{\text{Var}(\hat{\mu}_t)}\), where
\(L\) is a multiplier (commonly 2.5--3.0) determined by the desired
false alarm rate {[}13, 14{]}.

Lucas and Saccucci {[}13{]} demonstrated that EWMA control charts are
particularly effective at detecting small, persistent shifts in the
process mean---precisely the type of signal generated by slow
brute-force attacks that produce a gradual, sustained elevation in
anomaly scores. Montgomery {[}14{]} provides a comprehensive treatment
of EWMA in the broader context of statistical quality control and
discusses the selection of \(\alpha\) and \(L\) for different detection
objectives.

\subsection{2.9.3 Adaptive Percentile
Method}\label{adaptive-percentile-method}

The Adaptive Percentile method determines the threshold based on the
empirical quantile of the anomaly score distribution within a sliding
time window {[}14, 17{]}. The threshold at time \(t\) is defined as:

\[\theta_t = P_q(S_W)\]

where \(P_q\) denotes the \(q\)-th percentile and
\(S_W = \{s_{t-W+1}, s_{t-W+2}, \ldots, s_t\}\) is the set of the most
recent \(W\) anomaly scores {[}14{]}. This approach has two key
advantages: it requires no distributional assumptions about the anomaly
scores (unlike EWMA, which implicitly assumes approximate normality for
the control limit calculation), and it naturally adapts to changes in
the score distribution by continuously updating the reference window
{[}17{]}. The primary disadvantage is sensitivity to the window size
\(W\): too small a window results in volatile thresholds, while too
large a window reduces responsiveness to genuine distributional shifts
{[}14{]}.

\subsection{2.9.4 Hybrid EWMA-Adaptive Percentile
Method}\label{hybrid-ewma-adaptive-percentile-method}

This thesis proposes a hybrid dynamic thresholding mechanism that
combines EWMA and Adaptive Percentile to leverage the complementary
strengths of both approaches {[}13, 14{]}. EWMA provides long-term trend
tracking and smoothing, preventing the threshold from oscillating in
response to transient score fluctuations. Adaptive Percentile provides
distribution-free adaptiveness, accurately reflecting the empirical
score distribution without parametric assumptions. The hybrid threshold
is computed as a weighted combination:

\[\theta_t^{\text{hybrid}} = w_{\text{EWMA}} \cdot \theta_t^{\text{EWMA}} + w_{\text{AP}} \cdot \theta_t^{\text{AP}}\]

where \(\theta_t^{\text{EWMA}}\) is the EWMA-based threshold,
\(\theta_t^{\text{AP}}\) is the Adaptive Percentile threshold, and
\(w_{\text{EWMA}} + w_{\text{AP}} = 1\) are the blending weights {[}13,
14{]}. The specific weight values and the calibration procedure are
detailed in Chapter 3 (Methodology). The hybrid approach ensures that
the system maintains stable baseline tracking (from EWMA) while
remaining responsive to local distributional changes (from Adaptive
Percentile), a combination that is particularly valuable for detecting
the gradual anomaly score elevation characteristic of slow brute-force
campaigns {[}13, 23{]}.

{[}Figure 2.4: Comparison of static, EWMA, Adaptive Percentile, and
hybrid thresholds on a simulated anomaly score time series with embedded
attack periods{]}

\section{2.10 Related Work in SSH Brute-Force
Detection}\label{related-work-in-ssh-brute-force-detection}

This section presents a systematic review of the most relevant prior
studies in SSH brute-force detection, organized by methodological
approach. The review encompasses both international and domestic
(Vietnamese) research and identifies the specific contributions and
limitations of each study relative to the objectives of this thesis.

\subsection{2.10.1 Flow-Based and Network-Layer
Approaches}\label{flow-based-and-network-layer-approaches}

Sperotto et al.~{[}15{]} conducted pioneering work on flow-based
intrusion detection for SSH, analyzing network flow records from the
University of Twente (Netherlands) to characterize SSH attack patterns.
The study identified that flow-level features---including flow duration,
packet count, and byte count---provide sufficient information to
distinguish between legitimate SSH sessions and brute-force attacks with
greater than 90\% accuracy using Hidden Markov Models. The primary
contribution was the demonstration that application-layer payload
inspection is unnecessary for SSH attack detection; metadata features
alone are sufficient {[}15{]}. However, the study employed supervised
classification requiring labeled data, and the flow-based feature set
does not capture application-layer behavioral nuances such as username
diversity or password retry patterns that are available from
authentication logs {[}15, 24{]}.

Hellemons et al.~{[}5{]} developed SSHCure, a three-phase detection
system that models SSH attacks as a sequence of scanning, brute-force,
and compromise phases, each characterized by distinct flow-level
features. SSHCure achieved high detection rates (\textgreater95\% true
positive rate) for attacks that follow the canonical three-phase pattern
{[}5{]}. The key limitation is the assumption of sequential phase
progression, which does not hold for all attack variants---particularly
credential-stuffing attacks that skip the scanning phase and slow
brute-force attacks that blur the boundary between phases {[}5, 23{]}.

\subsection{2.10.2 Machine Learning Approaches on Benchmark
Datasets}\label{machine-learning-approaches-on-benchmark-datasets}

Ahmad et al.~{[}25{]} conducted a systematic study of network-based
intrusion detection systems (NIDS), comparing multiple machine learning
algorithms including Isolation Forest, LOF, OCSVM, Random Forest, and
deep learning approaches across the NSL-KDD, CICIDS2017, and UNSW-NB15
benchmark datasets. The study found that Isolation Forest achieved the
best balance between detection performance and computational efficiency
among unsupervised methods, with F1-scores ranging from 82\% to 91\%
depending on the dataset {[}25{]}. However, the study evaluated
algorithms on general network intrusion data rather than SSH-specific
data, and the benchmark datasets---while valuable for standardized
comparison---are known to suffer from age-related limitations, as
documented by Ring et al.~{[}24{]}.

Ring et al.~{[}24{]} provided a comprehensive survey of datasets for
network-based intrusion detection systems, evaluating 34 datasets across
criteria including realism, labeling accuracy, attack diversity, and
temporal currency. The survey concluded that many widely used benchmark
datasets (including NSL-KDD, first released in 1999) do not adequately
represent modern attack characteristics, and recommended the use of
application-specific datasets collected from controlled environments or
honeypots for domain-focused research {[}24{]}. This recommendation
directly motivates the use of honeypot-collected data in this thesis.

Goldstein and Uchida {[}20{]} performed a comparative evaluation of
unsupervised anomaly detection algorithms on 10 benchmark datasets,
including network intrusion data. Their results showed that Isolation
Forest ranked among the top three algorithms in average performance
across all datasets, with LOF and OCSVM achieving competitive
performance on specific datasets but exhibiting higher variance across
datasets {[}20{]}. The study confirmed that algorithm selection should
be guided by domain-specific evaluation rather than reliance on a single
benchmark.

\subsection{2.10.3 SSH-Specific Machine Learning
Studies}\label{ssh-specific-machine-learning-studies}

Javed and Paxson {[}23{]} studied stealthy SSH brute-forcing campaigns
using data from a large academic network, identifying attacks that
operate at rates as low as 1--2 attempts per hour and unfold over
periods of days to weeks. The study demonstrated that these stealthy
campaigns can achieve compromise rates comparable to aggressive attacks
while evading all conventional detection mechanisms {[}23{]}. The key
insight---that temporal behavioral features are essential for detecting
low-rate attacks---directly informs the feature engineering approach of
this thesis, which prioritizes temporal features (inter-attempt timing,
session duration) alongside traditional frequency features.

Bezerra et al.~{[}18{]} analyzed SSH attack durations using data from a
distributed honeypot network and found significant heterogeneity in
attack temporal profiles: attack sessions ranged from sub-second scans
to multi-day campaigns, with distinct temporal signatures associated
with different attack tools and botnets. The study identified
inter-attempt timing regularity as a highly discriminative feature for
distinguishing automated attacks from human-initiated sessions {[}18{]}.

Owens and Matthews {[}19{]} conducted an early but influential analysis
of SSH brute-force password characteristics using honeypot data. The
study documented that the top 20 most-attempted passwords account for
over 40\% of all attack attempts, and that attackers exhibit strong
preferences for specific username-password combinations (e.g.,
root/root, admin/123456) {[}19{]}. These findings inform the
authentication-category features in this thesis's 14-feature set,
particularly the unique username count and password diversity metrics.

Park et al.~{[}22{]} investigated SSH attacks targeting network
infrastructure devices (routers and switches), documenting that these
attacks exhibit distinct behavioral patterns compared to attacks against
general-purpose servers, including different username preferences and
timing profiles {[}22{]}. While this thesis focuses on server-targeted
attacks, the existence of device-specific attack patterns underscores
the importance of adaptive, learning-based detection approaches that can
accommodate behavioral diversity without requiring explicit rules for
each target type.

\subsection{2.10.4 Deep Learning and Advanced
Approaches}\label{deep-learning-and-advanced-approaches}

Nassif et al.~{[}21{]} surveyed machine learning and deep learning
approaches for anomaly-based intrusion detection, finding that deep
learning methods (autoencoders, LSTMs, GANs) can achieve higher
detection accuracy than traditional machine learning methods on large,
high-dimensional datasets. However, the survey also noted that deep
learning methods require significantly more training data, computational
resources, and hyperparameter tuning, and that their black-box nature
complicates interpretation and debugging in security-critical
applications {[}21{]}. For real-time SSH monitoring on
resource-constrained environments, traditional machine learning methods
such as Isolation Forest offer a more practical balance of performance
and efficiency {[}21, 25{]}.

Buczak and Guven {[}12{]} provided the foundational survey of data
mining and machine learning methods for cybersecurity intrusion
detection, cataloging over 50 studies spanning supervised, unsupervised,
and hybrid approaches. The survey identified several persistent
challenges in the field: the reliance on outdated benchmark datasets,
the absence of standardized evaluation methodologies, the difficulty of
comparing results across studies that use different datasets and
metrics, and the gap between algorithmic research and operational
deployment {[}12{]}. These challenges motivate the methodological
choices of this thesis, including the use of real honeypot data, the
standardized evaluation framework (accuracy, precision, recall, F1,
FPR), and the emphasis on end-to-end system integration.

Pimentel et al.~{[}17{]} reviewed novelty detection methods---a
formulation closely related to semi-supervised anomaly
detection---across multiple application domains including intrusion
detection. The review categorized methods into probabilistic,
distance-based, reconstruction-based, domain-based, and
information-theoretic approaches, and identified one-class
classification (including OCSVM) and isolation-based methods (including
Isolation Forest) as particularly well-suited for applications where the
target class (normal behavior) is well-sampled but the outlier class
(attacks) is poorly characterized or entirely absent from training data
{[}17{]}.

Hariri et al.~{[}16{]} proposed the Extended Isolation Forest (EIF),
which addresses the axis-parallel bias of the original algorithm by
using random hyperplane splits. EIF demonstrated improved anomaly
detection accuracy on synthetic datasets with complex geometric
structures. However, for the tabular feature data characteristic of SSH
log analysis, the performance difference between IF and EIF is marginal,
and the original IF's lower computational cost and simpler
implementation make it the preferred choice for real-time applications
{[}16, 20{]}.

{[}Table 2.4: Comprehensive comparison of related works in SSH
brute-force detection{]}

{\def\LTcaptype{none} % do not increment counter
\begin{longtable}[]{@{}
  >{\raggedright\arraybackslash}p{(\linewidth - 14\tabcolsep) * \real{0.0714}}
  >{\raggedright\arraybackslash}p{(\linewidth - 14\tabcolsep) * \real{0.0612}}
  >{\raggedright\arraybackslash}p{(\linewidth - 14\tabcolsep) * \real{0.0816}}
  >{\raggedright\arraybackslash}p{(\linewidth - 14\tabcolsep) * \real{0.1327}}
  >{\raggedright\arraybackslash}p{(\linewidth - 14\tabcolsep) * \real{0.1224}}
  >{\raggedright\arraybackslash}p{(\linewidth - 14\tabcolsep) * \real{0.1837}}
  >{\raggedright\arraybackslash}p{(\linewidth - 14\tabcolsep) * \real{0.1531}}
  >{\raggedright\arraybackslash}p{(\linewidth - 14\tabcolsep) * \real{0.1939}}@{}}
\toprule\noalign{}
\begin{minipage}[b]{\linewidth}\raggedright
Study
\end{minipage} & \begin{minipage}[b]{\linewidth}\raggedright
Year
\end{minipage} & \begin{minipage}[b]{\linewidth}\raggedright
Method
\end{minipage} & \begin{minipage}[b]{\linewidth}\raggedright
Data Source
\end{minipage} & \begin{minipage}[b]{\linewidth}\raggedright
Key Metric
\end{minipage} & \begin{minipage}[b]{\linewidth}\raggedright
Dynamic Threshold
\end{minipage} & \begin{minipage}[b]{\linewidth}\raggedright
Early Warning
\end{minipage} & \begin{minipage}[b]{\linewidth}\raggedright
System Integration
\end{minipage} \\
\midrule\noalign{}
\endhead
\bottomrule\noalign{}
\endlastfoot
Hellemons et al.~{[}5{]} & 2012 & SSHCure (three-phase) & Flow data &
TPR \textgreater{} 95\% & No & Partial & No \\
Sperotto et al.~{[}15{]} & 2010 & HMM (flow-based) & Univ. Twente flows
& Acc \textgreater{} 90\% & No & No & No \\
Javed \& Paxson {[}23{]} & 2013 & Statistical analysis & Academic
network & -- & No & Yes (analysis) & No \\
Owens \& Matthews {[}19{]} & 2008 & Password analysis & Honeypot & -- &
No & No & No \\
Bezerra et al.~{[}18{]} & 2019 & Temporal analysis & Distributed
honeypot & -- & No & No & No \\
Goldstein \& Uchida {[}20{]} & 2016 & IF, LOF, OCSVM (comparison) & 10
benchmarks & AUC varies & No & No & No \\
Nassif et al.~{[}21{]} & 2021 & ML/DL survey & Multiple & Survey & No &
No & No \\
Park et al.~{[}22{]} & 2020 & SSH on routers & Router logs & -- & No &
No & No \\
Ahmad et al.~{[}25{]} & 2021 & Multiple ML & NSL-KDD, CICIDS & F1
82--91\% & No & No & No \\
Ring et al.~{[}24{]} & 2019 & Dataset survey & 34 datasets & Survey &
N/A & N/A & N/A \\
Hariri et al.~{[}16{]} & 2019 & Extended IF & Synthetic + real & AUC
improved & No & No & No \\
\textbf{This thesis} & \textbf{2026} & \textbf{IF + EWMA-AP} &
\textbf{Honeypot + Sim (174K lines)} & \textbf{F1 = 93.74\%} &
\textbf{Yes (hybrid)} & \textbf{Yes} & \textbf{Yes (9 Docker
services)} \\
\end{longtable}
}

\section{2.11 Summary of the Literature
Review}\label{summary-of-the-literature-review}

The comprehensive review of the literature reveals a rich but fragmented
research landscape characterized by significant methodological advances
in individual components---anomaly detection algorithms, feature
engineering, threshold mechanisms---but a persistent absence of
integrated solutions that combine these components into deployable
systems {[}7, 12, 15{]}.

\textbf{On algorithms:} Isolation Forest {[}8, 9{]} has emerged as one
of the most effective unsupervised anomaly detection algorithms for
network security applications, offering a compelling combination of
detection performance, computational efficiency, and interpretability
{[}20, 21{]}. LOF {[}10{]} provides complementary local-density-based
detection but is limited by computational complexity and memory
requirements. OCSVM {[}11{]} offers strong theoretical foundations but
requires careful hyperparameter tuning and has higher computational
costs. All three algorithms have been validated on benchmark datasets,
but their comparative evaluation on real-world SSH honeypot data in a
semi-supervised configuration remains underexplored {[}20, 24, 25{]}.

\textbf{On features:} The literature consistently identifies temporal
features (inter-attempt timing, session duration) as more discriminative
than frequency features (attempt counts) for distinguishing automated
SSH attacks from legitimate activity, particularly for slow-rate and
stealthy attack variants {[}5, 15, 18, 23{]}. Authentication features
(username diversity, password patterns) provide additional
discriminative power for credential-stuffing and dictionary attacks
{[}19{]}. The integration of features from all four categories
(frequency, temporal, authentication, connection) in a unified feature
set has been recommended but rarely implemented and evaluated
systematically {[}12, 15{]}.

\textbf{On thresholds:} Static thresholds are inadequate for SSH
monitoring due to the non-stationary nature of authentication traffic
{[}6, 7, 13{]}. EWMA-based adaptive thresholds offer superior
sensitivity to small, persistent shifts {[}13, 14{]}, while
percentile-based methods provide distribution-free adaptiveness {[}14,
17{]}. Hybrid approaches combining multiple threshold methods have been
proposed in the statistical process control literature {[}13, 14{]} but
have not been applied to SSH anomaly detection.

\textbf{On integration:} The gap between algorithmic research and
deployable systems remains one of the most significant barriers to the
practical adoption of machine learning in cybersecurity {[}12, 15{]}.
The vast majority of studies evaluate algorithms in isolation on
benchmark datasets without addressing log ingestion, real-time feature
extraction, automated response, visualization, or containerized
deployment {[}20, 21, 24, 25{]}.

\section{2.12 Research Gap and
Contribution}\label{research-gap-and-contribution}

Based on the literature review, four specific research gaps are
identified that collectively define the novel contribution of this
thesis:

\textbf{Gap 1: Absence of comparative algorithm evaluation on real SSH
honeypot data.} While Isolation Forest, LOF, and OCSVM have been
extensively compared on benchmark datasets {[}20, 25{]}, their
comparative performance on authentic SSH honeypot data---with its
characteristic noise, class imbalance (679 attacking IPs versus 64
normal users), and behavioral diversity (five attack variants)---has not
been systematically evaluated in a semi-supervised configuration
{[}24{]}. \textbf{Contribution:} This thesis provides a rigorous
comparative evaluation of all three algorithms on 174,250 lines of
combined honeypot and simulated data, reporting standardized metrics
(accuracy, precision, recall, F1, FPR) with optimized results: IF (Acc =
90.31\%, F1 = 93.74\%, FPR = 29.00\%), LOF (Acc = 83.22\%, F1 =
89.94\%), OCSVM (Acc = 91.38\%, F1 = 94.55\%).

\textbf{Gap 2: No hybrid dynamic thresholding for SSH anomaly
detection.} EWMA {[}13{]} and percentile-based thresholds {[}14{]} have
been studied independently, but their combination into a hybrid
mechanism specifically designed for SSH anomaly detection has not been
previously proposed or evaluated {[}7, 12{]}. Existing SSH detection
systems rely on either static thresholds {[}6{]} or simple adaptive
mechanisms without the complementary strengths of trend-following and
distribution-free approaches {[}5, 15{]}. \textbf{Contribution:} This
thesis proposes and validates the EWMA-Adaptive Percentile hybrid
thresholding mechanism, demonstrating its effectiveness across five
attack scenarios with varying temporal profiles {[}13, 14{]}.

\textbf{Gap 3: Limited exploitation of early-warning capability.} While
some studies have identified the potential for early detection based on
behavioral features {[}5, 23{]}, the systematic exploitation of
early-phase behavioral signatures for proactive attack
warning---particularly for slow-rate attacks---remains largely
unexplored {[}18, 23{]}. \textbf{Contribution:} This thesis demonstrates
early-warning capability across all five attack scenarios, with the
system issuing EARLY\_WARNING alerts after only 3--5 attempts in the
slow brute-force scenario, compared to complete detection failure by
Fail2Ban under default settings.

\textbf{Gap 4: Research-to-deployment gap.} The overwhelming majority of
studies in this domain evaluate algorithms in isolation without
addressing system integration {[}12, 20, 21, 25{]}.
\textbf{Contribution:} This thesis implements a complete end-to-end
system comprising 9 Docker services (ELK Stack, FastAPI, React,
Fail2Ban), with a 14-feature extraction pipeline processing SSH logs in
5-minute sliding windows, demonstrating that the transition from
algorithm to deployable system is achievable within a modern
containerized architecture {[}6, 15{]}.

{[}Figure 2.5: Visual mapping of research gaps to the contributions of
this thesis{]}

{[}Table 2.5: Summary of research gaps, their evidence in the
literature, and corresponding contributions{]}

{\def\LTcaptype{none} % do not increment counter
\begin{longtable}[]{@{}
  >{\raggedright\arraybackslash}p{(\linewidth - 4\tabcolsep) * \real{0.0926}}
  >{\raggedright\arraybackslash}p{(\linewidth - 4\tabcolsep) * \real{0.4074}}
  >{\raggedright\arraybackslash}p{(\linewidth - 4\tabcolsep) * \real{0.5000}}@{}}
\toprule\noalign{}
\begin{minipage}[b]{\linewidth}\raggedright
Gap
\end{minipage} & \begin{minipage}[b]{\linewidth}\raggedright
Evidence in Literature
\end{minipage} & \begin{minipage}[b]{\linewidth}\raggedright
This Thesis's Contribution
\end{minipage} \\
\midrule\noalign{}
\endhead
\bottomrule\noalign{}
\endlastfoot
No comparative evaluation on real SSH data & Goldstein \& Uchida
{[}20{]} used benchmarks; Ahmad et al.~{[}25{]} used general IDS
datasets & IF, LOF, OCSVM compared on 174,250-line honeypot + simulation
dataset \\
No hybrid dynamic threshold for SSH & Fail2Ban {[}6{]} uses static;
SSHCure {[}5{]} uses fixed phases & EWMA-Adaptive Percentile hybrid
mechanism \\
Limited early-warning exploitation & Javed \& Paxson {[}23{]} identified
potential; no system implemented & Early warning in 3--5 attempts for
slow attacks \\
Research-to-deployment gap & Buczak \& Guven {[}12{]} and Sperotto et
al.~{[}15{]} identified the gap & 9-service Docker architecture with
automated response \\
\end{longtable}
}

\chapter{CHAPTER 3: METHODOLOGY}\label{chapter-3-methodology}

\section{3.1 Research Design Overview}\label{research-design-overview}

\subsection{3.1.1 Philosophical and Methodological
Foundations}\label{philosophical-and-methodological-foundations}

This research adopts a positivist epistemological stance, employing a
quantitative experimental methodology to design, implement, and evaluate
an AI-based SSH brute-force attack detection and prevention system. The
research follows the Design Science Research Methodology (DSRM)
framework proposed by Peffers et al.~{[}41{]}, which is widely regarded
as the standard methodology for information systems research involving
the creation of IT artifacts. DSRM comprises six iterative phases: (1)
problem identification, (2) objective definition, (3) design and
development, (4) demonstration, (5) evaluation, and (6) communication.
This chapter addresses phases (3) and (4), while Chapter 4 presents the
evaluation results.

The experimental approach is appropriate because the research question
--- whether an Isolation Forest model combined with an EWMA-Adaptive
Percentile dynamic threshold can effectively detect SSH brute-force
attacks in real time --- is inherently quantitative and amenable to
controlled experimentation. The dependent variables (accuracy,
precision, recall, F1-score, ROC-AUC, false positive rate) can be
precisely measured, and the independent variables (model
hyperparameters, threshold parameters, feature set composition, window
configuration) can be systematically manipulated to assess their effects
on detection performance.

The methodology integrates three complementary research strategies:

\textbf{Empirical data collection.} Real-world attack data is obtained
from an SSH honeypot deployed on the public Internet, while normal
behavioral data is generated through a controlled simulation
environment. This dual-source approach addresses the fundamental data
challenge in anomaly detection: the simultaneous need for clean normal
data (for training) and authentic attack data (for evaluation) {[}42{]}.

\textbf{Comparative algorithm evaluation.} Three unsupervised anomaly
detection algorithms --- Isolation Forest (IF), Local Outlier Factor
(LOF), and One-Class Support Vector Machine (OCSVM) --- are
systematically compared under identical conditions, following the best
practices for algorithm benchmarking established by Demsar {[}43{]} and
Campos et al.~{[}44{]}.

\textbf{System integration and scenario-based testing.} The detection
algorithm is embedded within a complete end-to-end system comprising 9
Docker services, and evaluated against five distinct attack scenarios
designed to cover the spectrum of SSH brute-force attack variants
documented in the literature.

\subsection{3.1.2 Semi-Supervised Anomaly Detection
Paradigm}\label{semi-supervised-anomaly-detection-paradigm}

The research adopts the semi-supervised anomaly detection paradigm, also
referred to as novelty detection {[}45{]}, in which models are trained
exclusively on normal (in-distribution) data and subsequently identify
observations that deviate from the learned normal distribution as
anomalies. This paradigm was selected over supervised classification and
fully unsupervised clustering for three theoretically and practically
motivated reasons.

First, \textbf{practicality of data acquisition.} In operational
cybersecurity environments, normal data is far more abundant and easier
to collect than labeled attack data. Organizations can readily extract
logs of normal SSH operations from their existing infrastructure, but
they cannot feasibly obtain representative samples of every possible
attack type, including zero-day variants that have not yet been observed
{[}46{]}. Chandola et al.~{[}47{]} formalized this observation as the
``labeled data scarcity problem'' in their seminal survey of anomaly
detection techniques.

Second, \textbf{generalization to novel attacks.} Models trained on
attack signatures can only detect attacks whose signatures match the
training data. Semi-supervised models, by contrast, can detect any
attack type whose behavioral characteristics deviate from the learned
normal profile, including attack variants never previously encountered
(zero-day attacks). This property was demonstrated empirically by
Goldstein and Uchida {[}48{]}, who showed that novelty detection methods
consistently outperform signature-based methods on datasets containing
previously unseen anomaly types.

Third, \textbf{avoidance of class imbalance pathologies.} Supervised
classification methods are well known to suffer from degraded
performance when the class distribution is severely imbalanced {[}49{]},
as is typical in cybersecurity datasets where normal traffic vastly
outnumbers attack traffic (or vice versa, as in honeypot data). The
semi-supervised paradigm is architecturally immune to this problem
because the training set contains only a single class.

The formal framework is as follows. Let
\(\mathbf{X}_{train} = \{\mathbf{x}_1, \mathbf{x}_2, \ldots, \mathbf{x}_n\}\)
denote the training set, where each \(\mathbf{x}_i \in \mathbb{R}^{14}\)
is a feature vector representing the SSH behavioral profile of a single
IP address within a single 5-minute time window, and all
\(\mathbf{x}_i\) are drawn from the normal class. A model
\(f: \mathbb{R}^{14} \rightarrow \mathbb{R}\) is trained on
\(\mathbf{X}_{train}\) to produce an anomaly score \(s(\mathbf{x})\) for
any new observation \(\mathbf{x}\). A threshold function \(\theta(t)\)
maps the anomaly score to a binary classification:

\[\hat{y}(\mathbf{x}, t) = \begin{cases} \text{attack} & \text{if } s(\mathbf{x}) > \theta(t) \\ \text{normal} & \text{otherwise} \end{cases}\]

where \(\theta(t)\) may be static (a fixed constant) or dynamic (a
function of recent score history), as described in Section 3.8.

\section{3.2 System Architecture}\label{system-architecture}

\subsection{3.2.1 High-Level
Architecture}\label{high-level-architecture}

The SSH brute-force attack detection and prevention system is designed
following a microservices architecture, deployed on Docker Compose with
a total of 9 services operating in coordination. The microservices
architecture was chosen over a monolithic design for three reasons
documented by Newman {[}50{]}: independent scalability of components,
fault isolation (failure of one service does not cascade to others), and
technology heterogeneity (each service can use the language and
framework best suited to its function).

{[}Figure 3.1: Overall system architecture for SSH brute-force attack
detection and prevention{]}

The system is organized into five functional layers, each with a clearly
defined responsibility:

\textbf{Layer 1: Data Collection.} This layer is responsible for
ingesting authentication logs from SSH servers. Logstash tails the
\texttt{/var/log/auth.log} file on target SSH servers, applies Grok
pattern parsing to extract structured fields (timestamp, hostname,
service, PID, event type, username, IP address, port), and forwards the
structured events to Elasticsearch for indexed storage. The use of
Logstash for log ingestion follows the established ELK Stack
architecture, which Elasticsearch B.V. reports is deployed in over
500,000 organizations worldwide for security analytics {[}51{]}.

\textbf{Layer 2: Processing and Analysis.} The central computational
layer comprises the data preprocessing module, the feature extraction
module, and the AI model inference module. These modules are implemented
as asynchronous RESTful APIs using the FastAPI framework (Python), which
leverages ASGI (Asynchronous Server Gateway Interface) for
high-concurrency request handling. FastAPI was selected over Flask and
Django REST Framework based on benchmarks by Lust {[}52{]} showing
3--10x throughput improvements for I/O-bound workloads.

\textbf{Layer 3: Decision.} This layer implements the dynamic threshold
algorithm based on the EWMA-Adaptive Percentile method (detailed in
Section 3.8). It receives anomaly scores from the AI model and produces
a binary classification decision (normal or attack) together with a
confidence level (EARLY\_WARNING or ALERT).

\textbf{Layer 4: Response.} Fail2Ban is integrated to execute automated
prevention measures upon attack detection. When an ALERT-level detection
occurs, the system invokes the Fail2Ban API to ban the offending IP
address, applies a configurable ban duration with progressive escalation
for repeat offenders, and dispatches multi-channel alert notifications
(dashboard, webhook, log).

\textbf{Layer 5: Presentation.} A single-page web application developed
in React provides a real-time monitoring dashboard displaying detection
events, attack statistics, IP geolocation, anomaly score trends, and
administrative configuration controls.

\subsection{3.2.2 Service Inventory}\label{service-inventory}

{[}Table 3.1: Complete inventory of 9 Docker services{]}

{\def\LTcaptype{none} % do not increment counter
\begin{longtable}[]{@{}
  >{\raggedright\arraybackslash}p{(\linewidth - 8\tabcolsep) * \real{0.1220}}
  >{\raggedright\arraybackslash}p{(\linewidth - 8\tabcolsep) * \real{0.2195}}
  >{\raggedright\arraybackslash}p{(\linewidth - 8\tabcolsep) * \real{0.2683}}
  >{\raggedright\arraybackslash}p{(\linewidth - 8\tabcolsep) * \real{0.2439}}
  >{\raggedright\arraybackslash}p{(\linewidth - 8\tabcolsep) * \real{0.1463}}@{}}
\toprule\noalign{}
\begin{minipage}[b]{\linewidth}\raggedright
No.
\end{minipage} & \begin{minipage}[b]{\linewidth}\raggedright
Service
\end{minipage} & \begin{minipage}[b]{\linewidth}\raggedright
Technology
\end{minipage} & \begin{minipage}[b]{\linewidth}\raggedright
Function
\end{minipage} & \begin{minipage}[b]{\linewidth}\raggedright
Port
\end{minipage} \\
\midrule\noalign{}
\endhead
\bottomrule\noalign{}
\endlastfoot
1 & Detector & Python (scikit-learn) & AI model inference and threshold
computation & Internal \\
2 & API Server & FastAPI (Python) & Business logic, REST API,
orchestration & 8000 \\
3 & Frontend & React (TypeScript) & Real-time monitoring dashboard &
3000 \\
4 & Elasticsearch & Elasticsearch 8.x & Log storage, full-text indexing,
search & 9200 \\
5 & Logstash & Logstash 8.x & Log ingestion, parsing, forwarding &
5044 \\
6 & Kibana & Kibana 8.x & Data visualization, ad hoc exploration &
5601 \\
7 & Redis & Redis 7.x & Cache, message queue, sliding window state &
6379 \\
8 & Fail2Ban & Fail2Ban 1.x & Automated IP blocking via
iptables/nftables & Internal \\
9 & SSH Target & OpenSSH 9.x & Target SSH server for attack scenario
testing & 22 \\
\end{longtable}
}

\subsection{3.2.3 Data Flow Pipeline}\label{data-flow-pipeline}

{[}Figure 3.2: End-to-end data flow pipeline from SSH log event to
classification decision and response{]}

The data processing pipeline operates through five sequential stages
with well-defined interfaces:

\begin{enumerate}
\def\labelenumi{\arabic{enumi}.}
\item
  \textbf{Ingestion.} SSH authentication events on the target server are
  written to \texttt{/var/log/auth.log} by the sshd process. Logstash
  monitors this file via the file input plugin (using inode tracking to
  handle log rotation) and applies a Grok filter to parse each line into
  structured JSON fields.
\item
  \textbf{Indexing.} Parsed events are forwarded to Elasticsearch via
  the Elasticsearch output plugin. Events are indexed in
  time-partitioned indices (e.g., \texttt{ssh-auth-2026.03.26}) with a
  mapping that preserves the original data types (keyword for IP
  addresses, date for timestamps, integer for port numbers).
\item
  \textbf{Feature Extraction.} The API Server periodically queries
  Elasticsearch for new events (polling interval: 30 seconds) or
  receives push notifications via a webhook. For each active IP address,
  the server maintains a sliding window buffer (stored in Redis for
  persistence across restarts) and computes the 14-feature vector as
  described in Section 3.6.
\item
  \textbf{Inference and Decision.} The feature vector is normalized
  using the pre-fitted RobustScaler, passed to the trained anomaly
  detection model to obtain an anomaly score, and compared against the
  current dynamic threshold \(\theta(t)\) to produce a classification
  decision.
\item
  \textbf{Action.} If the classification is ALERT, the system invokes
  the Fail2Ban REST API to ban the IP, writes the detection event to
  Elasticsearch for audit logging, and pushes a real-time notification
  to the React dashboard via WebSocket.
\end{enumerate}

\section{3.3 Laboratory Environment
Setup}\label{laboratory-environment-setup}

\subsection{3.3.1 Hardware and Software
Configuration}\label{hardware-and-software-configuration}

All experiments were conducted on a single physical machine to ensure
reproducibility and eliminate network variability as a confounding
factor. The Docker Compose configuration ensures that all 9 services run
within a single Docker network with deterministic inter-service
communication.

{[}Table 3.2: Laboratory environment specifications{]}

{\def\LTcaptype{none} % do not increment counter
\begin{longtable}[]{@{}
  >{\raggedright\arraybackslash}p{(\linewidth - 2\tabcolsep) * \real{0.4400}}
  >{\raggedright\arraybackslash}p{(\linewidth - 2\tabcolsep) * \real{0.5600}}@{}}
\toprule\noalign{}
\begin{minipage}[b]{\linewidth}\raggedright
Component
\end{minipage} & \begin{minipage}[b]{\linewidth}\raggedright
Specification
\end{minipage} \\
\midrule\noalign{}
\endhead
\bottomrule\noalign{}
\endlastfoot
Operating System & Ubuntu 22.04 LTS (host) / Docker containers \\
Containerization & Docker 24.x, Docker Compose 2.x \\
Python & 3.10 with scikit-learn 1.3.x, pandas 2.1.x, numpy 1.25.x \\
Machine Learning & scikit-learn (IsolationForest, LocalOutlierFactor,
OneClassSVM) \\
Web Framework & FastAPI 0.104.x (ASGI) \\
Frontend & React 18.x with TypeScript \\
Search Engine & Elasticsearch 8.x \\
Log Pipeline & Logstash 8.x \\
Visualization & Kibana 8.x \\
Cache & Redis 7.x \\
Prevention & Fail2Ban 1.x \\
\end{longtable}
}

\subsection{3.3.2 Honeypot Deployment}\label{honeypot-deployment}

The SSH honeypot was deployed on a Virtual Private Server (VPS) with a
publicly routable IPv4 address. The honeypot was configured to attract
realistic SSH brute-force attack traffic through several deliberate
design choices informed by the honeypot deployment guidelines of Provos
and Holz {[}53{]}:

\begin{itemize}
\tightlist
\item
  \textbf{Port selection.} SSH was exposed on the default port 22 rather
  than a non-standard port. Alata et al.~{[}54{]} demonstrated that the
  majority of automated SSH scanners target only port 22, and moving SSH
  to a non-standard port reduces attack traffic by 95--99\% --- which
  would be counterproductive for data collection purposes.
\item
  \textbf{Hostname.} The hostname was set to ``mail'' to simulate a mail
  server, a high-value target that attracts more sophisticated
  brute-force attacks beyond opportunistic scanning {[}55{]}.
\item
  \textbf{Authentication policy.} Password authentication was enabled
  with a permissive MaxAuthTries setting to maximize the number of
  events recorded per attacking session.
\item
  \textbf{Collection duration.} The honeypot operated continuously for 5
  days (2026-03-22 to 2026-03-27), during which all SSH authentication
  events were recorded to \texttt{/var/log/auth.log}.
\item
  \textbf{Administrative access.} Six (6) verified administrative IP
  addresses (27.64.18.8, 104.28.156.151, 104.28.159.126, 116.110.42.131,
  118.69.182.144, 14.169.70.183) were used by the research team for
  honeypot management. All logins from these IPs were recorded and later
  labeled as normal.
\end{itemize}

\subsection{3.3.3 Simulation Environment}\label{simulation-environment}

The simulation environment was designed to generate realistic normal SSH
activity representative of a medium-sized organization. The simulation
script, written in Python using the Paramiko library, creates 64 user
accounts (52 human users and 12 service accounts) and orchestrates their
SSH activity according to behavioral profiles calibrated from the
literature on organizational IT usage patterns {[}56{]}.

Key simulation parameters:

\begin{itemize}
\tightlist
\item
  \textbf{Duration.} Approximately 25 hours of continuous activity on
  2026-03-26.
\item
  \textbf{Hostname.} ``if'' (to distinguish from the honeypot hostname
  ``mail'').
\item
  \textbf{Network.} All connections originate from the 192.168.152.0/24
  subnet.
\item
  \textbf{Human users (52).} Activity concentrated during business hours
  (8:00--18:00) with reduced activity during evenings and weekends.
  Session durations range from 5 minutes to several hours. Occasional
  password mistyping is included at a natural rate.
\item
  \textbf{Service accounts (12).} Automated cron-based connections at
  regular intervals, short session durations, and zero password
  failures.
\item
  \textbf{Failure rate.} 177 failed logins out of 4,382 total attempts
  (4.04\%), reflecting the empirically observed rate of accidental
  password failures in organizational environments reported by Florencio
  and Herley {[}57{]}.
\end{itemize}

\section{3.4 Data Collection Methods}\label{data-collection-methods}

\subsection{3.4.1 Data Sources and
Rationale}\label{data-sources-and-rationale}

The research employs two primary data sources, following the dual-source
methodology advocated by Ring et al.~{[}58{]} for constructing intrusion
detection evaluation datasets. The dual-source approach addresses the
fundamental tension between data purity and data realism: controlled
environments produce clean, well-labeled data but may lack the
complexity and diversity of real-world traffic, while real-world
captures provide authentic attack patterns but resist precise labeling.

\textbf{Source 1: Honeypot data (honeypot\_auth.log).} This
authentication log was collected from the SSH honeypot described in
Section 3.3.2. The honeypot approach, first formalized by Spitzner
{[}59{]}, provides highly representative attack data because the traffic
is generated by real attackers using real tools and techniques, rather
than synthetic simulations.

{[}Table 3.3: Honeypot data summary statistics{]}

{\def\LTcaptype{none} % do not increment counter
\begin{longtable}[]{@{}ll@{}}
\toprule\noalign{}
Attribute & Value \\
\midrule\noalign{}
\endhead
\bottomrule\noalign{}
\endlastfoot
Total log lines & 119,729 \\
Collection period & 2026-03-22 to 2026-03-27 (5 days) \\
Hostname & ``mail'' \\
Unique IP addresses & 679 \\
Failed password events & 29,301 \\
Invalid user attempts & 12,799 \\
Accepted root logins (admin IPs) & 532 \\
Number of admin IPs & 6 \\
``message repeated'' entries & 1,920 \\
SSHD entries after parsing & 122,809 \\
Average failed attempts per day & 5,860 \\
Average failed attempts per IP & 43.15 \\
\end{longtable}
}

\textbf{Source 2: Simulation data (simulation\_auth.log).} This
authentication log was generated from the controlled simulation
environment described in Section 3.3.3.

{[}Table 3.4: Simulation data summary statistics{]}

{\def\LTcaptype{none} % do not increment counter
\begin{longtable}[]{@{}ll@{}}
\toprule\noalign{}
Attribute & Value \\
\midrule\noalign{}
\endhead
\bottomrule\noalign{}
\endlastfoot
Total log lines & 54,521 \\
Collection period & \textasciitilde25 hours (2026-03-26) \\
Hostname & ``if'' \\
User accounts & 64 (52 human + 12 service) \\
Network subnet & 192.168.152.0/24 \\
Accepted logins & 4,205 \\
Failed logins & 177 \\
Success rate & 95.96\% \\
Failure rate & 4.04\% \\
SSHD entries after parsing & 34,700 \\
\end{longtable}
}

\subsection{3.4.2 Log Data Format}\label{log-data-format}

SSH servers on Linux systems record authentication events to the system
authentication log (typically \texttt{/var/log/auth.log} on
Debian/Ubuntu distributions) via the syslog facility. Each log entry
follows the standard syslog format specified in RFC 5424 {[}60{]}:

\begin{verbatim}
<timestamp> <hostname> <service>[<PID>]: <message>
\end{verbatim}

The principal SSH event types relevant to brute-force detection, as
documented in the OpenSSH source code and manual pages {[}61{]}, are:

{\def\LTcaptype{none} % do not increment counter
\begin{longtable}[]{@{}
  >{\raggedright\arraybackslash}p{(\linewidth - 4\tabcolsep) * \real{0.2558}}
  >{\raggedright\arraybackslash}p{(\linewidth - 4\tabcolsep) * \real{0.4419}}
  >{\raggedright\arraybackslash}p{(\linewidth - 4\tabcolsep) * \real{0.3023}}@{}}
\toprule\noalign{}
\begin{minipage}[b]{\linewidth}\raggedright
Event Type
\end{minipage} & \begin{minipage}[b]{\linewidth}\raggedright
Log Message Pattern
\end{minipage} & \begin{minipage}[b]{\linewidth}\raggedright
Significance
\end{minipage} \\
\midrule\noalign{}
\endhead
\bottomrule\noalign{}
\endlastfoot
Failed password &
\texttt{Failed\ password\ for\ \textless{}user\textgreater{}\ from\ \textless{}IP\textgreater{}\ port\ \textless{}port\textgreater{}}
& Primary brute-force indicator \\
Invalid user &
\texttt{Invalid\ user\ \textless{}user\textgreater{}\ from\ \textless{}IP\textgreater{}\ port\ \textless{}port\textgreater{}}
& Username enumeration indicator \\
Accepted password &
\texttt{Accepted\ password\ for\ \textless{}user\textgreater{}\ from\ \textless{}IP\textgreater{}\ port\ \textless{}port\textgreater{}}
& Successful authentication \\
Accepted publickey &
\texttt{Accepted\ publickey\ for\ \textless{}user\textgreater{}\ from\ \textless{}IP\textgreater{}\ port\ \textless{}port\textgreater{}}
& Key-based authentication \\
Connection closed &
\texttt{Connection\ closed\ by\ \textless{}IP\textgreater{}\ port\ \textless{}port\textgreater{}}
& Session termination \\
PAM failure & \texttt{pam\_unix(sshd:auth):\ authentication\ failure} &
PAM-level failure \\
Max retries exceeded &
\texttt{error:\ maximum\ authentication\ attempts\ exceeded} & Retry
limit reached \\
Message repeated &
\texttt{message\ repeated\ \textless{}N\textgreater{}\ times} & Syslog
deduplication \\
\end{longtable}
}

The ``message repeated'' entries require special handling during
preprocessing. When syslog detects consecutive identical messages, it
consolidates them into a single ``message repeated N times'' entry. The
preprocessing pipeline expands these entries by replicating the
preceding event \(N\) times, which increased the honeypot dataset from
119,729 raw lines to 122,809 parsed SSHD entries.

\subsection{3.4.3 Labeling Strategy}\label{labeling-strategy}

The labeling strategy was designed to be both conservative and
verifiable, following the recommendations of Gates and Taylor {[}62{]}
for intrusion detection dataset labeling:

\textbf{Simulation data (simulation\_auth.log).} All events are labeled
as \textbf{normal}. The simulation environment is fully controlled, with
no external network access, and all activities are performed by the 64
legitimate user accounts. This includes the 177 failed login events,
which represent accidental password mistyping --- a behavior that is
part of the normal operating profile of legitimate users. Labeling these
failures as normal is a deliberate and theoretically motivated decision:
the anomaly detection model must learn that occasional failures are
normal, so that it does not flag every single failure as an attack.

\textbf{Honeypot data (honeypot\_auth.log).} The labeling is based on
the whitelist of 6 verified administrative IP addresses:

\begin{itemize}
\tightlist
\item
  \textbf{Normal:} Accepted root login events from the 6 admin IPs:
  \{27.64.18.8, 104.28.156.151, 104.28.159.126, 116.110.42.131,
  118.69.182.144, 14.169.70.183\}. These 532 sessions represent genuine
  system administration activity.
\item
  \textbf{Attack:} All remaining events, including all failed passwords,
  all invalid user attempts, all connections from non-admin IPs, and all
  other event types. Since the honeypot has no legitimate users other
  than the 6 administrators, every other event is by definition an
  attack or the direct consequence of an attack.
\end{itemize}

This binary labeling scheme produces a clean separation: every feature
vector derived from simulation data is labeled normal, and every feature
vector derived from honeypot data that is not attributable to the 6
admin IPs is labeled attack. The resulting label distribution is highly
imbalanced in favor of attack data in the honeypot (by design), which
provides a challenging and realistic test set for evaluating the
detection models.

\section{3.5 Data Preprocessing
Pipeline}\label{data-preprocessing-pipeline}

\subsection{3.5.1 Pipeline Overview}\label{pipeline-overview}

The data preprocessing pipeline transforms raw syslog text into
normalized numerical feature vectors suitable for anomaly detection
models. The pipeline comprises four sequential stages: (1) log parsing,
(2) labeling, (3) feature extraction via sliding window, and (4) feature
scaling. Each stage is implemented as an independent Python module with
well-defined inputs and outputs, following the scikit-learn Pipeline API
design pattern {[}63{]} to ensure reproducibility and prevent data
leakage.

{[}Figure 3.3: Data preprocessing pipeline from raw logs to scaled
feature vectors{]}

\subsection{3.5.2 Stage 1: Log Parsing}\label{stage-1-log-parsing}

The log parser applies regular expression matching to extract structured
fields from raw syslog lines. The parser handles the following
complexities:

\begin{enumerate}
\def\labelenumi{\arabic{enumi}.}
\item
  \textbf{Multi-format timestamps.} The syslog timestamp format varies
  across Linux distributions (e.g., \texttt{Mar\ 22\ 10:15:33}
  vs.~\texttt{2026-03-22T10:15:33.000+07:00}). The parser normalizes all
  timestamps to ISO 8601 format with UTC timezone.
\item
  \textbf{``Message repeated'' expansion.} As noted in Section 3.4.2,
  syslog consolidation entries are expanded. This expansion increased
  the honeypot dataset from 119,729 raw lines to 122,809 parsed SSHD
  entries, and kept the simulation dataset at 34,700 SSHD entries.
\item
  \textbf{Event type classification.} Each log line is classified into
  one of the event types listed in Section 3.4.2 based on pattern
  matching against the message content.
\item
  \textbf{Field extraction.} For each classified event, the parser
  extracts the relevant fields: username, source IP address, source port
  number, and authentication result.
\end{enumerate}

The parsing success rate exceeded 99.5\% on both datasets. Lines that
could not be parsed (typically non-SSH syslog entries or malformed
entries) were excluded from further analysis.

\subsection{3.5.3 Stage 2: Labeling}\label{stage-2-labeling}

Labels are assigned to each parsed event according to the strategy
defined in Section 3.4.3. The label is propagated to the feature
extraction stage, where it is aggregated at the window level: a feature
vector is labeled ``attack'' if \textbf{any} event within the
corresponding IP-window combination originates from a non-admin IP in
the honeypot data.

\subsection{3.5.4 Stage 3: Feature Extraction via Sliding
Window}\label{stage-3-feature-extraction-via-sliding-window}

Features are extracted using the sliding window method, a standard
technique in time-series anomaly detection {[}64{]}. The window
parameters are:

\begin{itemize}
\tightlist
\item
  \textbf{Window size (\(W\)):} 5 minutes
\item
  \textbf{Stride (\(S\)):} Variable (1 minute for baseline overlapping,
  5 minutes for optimized non-overlapping)
\item
  \textbf{Aggregation unit:} Per source IP address
\end{itemize}

For each 5-minute window, for each unique source IP address that appears
in the window, the pipeline computes a 14-dimensional feature vector
\(\mathbf{x} \in \mathbb{R}^{14}\) by aggregating all events from that
IP within the window boundaries. The formal definition is:

\[\mathbf{x}_{ip,t} = \text{Aggregate}(\{e \in \mathcal{E} \mid e.\text{ip} = ip \wedge t - W < e.\text{timestamp} \leq t\})\]

where \(\mathcal{E}\) is the set of all parsed events, and
\(\text{Aggregate}(\cdot)\) computes the 14 features described in
Section 3.6.

The choice of \(W = 5\) minutes was motivated by three considerations
documented in the SSH monitoring literature. First, Hellemons et
al.~{[}65{]} demonstrated that a 5-minute window is sufficiently short
to capture the temporal signature of rapid brute-force attacks (which
typically produce hundreds of events per minute), while being long
enough to accumulate statistically meaningful behavioral information for
slower attacks. Second, the 5-minute granularity aligns with the default
monitoring intervals used by industry-standard SIEM (Security
Information and Event Management) tools such as Splunk and QRadar
{[}66{]}, facilitating integration with existing security
infrastructure. Third, Bezerra et al.~{[}67{]} showed that windows
shorter than 2 minutes produce excessive noise due to insufficient event
counts, while windows longer than 10 minutes introduce unacceptable
detection latency.

\textbf{Feature vector counts after extraction:}

{[}Table 3.5: Feature vector counts by source{]}

{\def\LTcaptype{none} % do not increment counter
\begin{longtable}[]{@{}lll@{}}
\toprule\noalign{}
Source & Overlapping (stride=1 min) & Non-overlapping (stride=5 min) \\
\midrule\noalign{}
\endhead
\bottomrule\noalign{}
\endlastfoot
Simulation & 10,304 windows & Approx. 2,060 windows \\
Honeypot & 50,792 windows & Approx. 10,158 windows \\
\end{longtable}
}

\subsection{3.5.5 Stage 4: Data Splitting}\label{stage-4-data-splitting}

The data split follows the semi-supervised novelty detection paradigm,
where the training set must contain exclusively normal data and the test
set must contain both normal and attack data for evaluation:

\textbf{Training set:} - \textbf{Source:} First 70\% of the simulation
data (chronological split to preserve temporal ordering and prevent
look-ahead bias {[}68{]}) - \textbf{Size:} 7,212 samples -
\textbf{Composition:} 100\% normal - \textbf{Purpose:} Train models to
learn the distributional characteristics of normal SSH behavior

\textbf{Test set:} - \textbf{Source:} Remaining 30\% of simulation data
(normal) + all honeypot data (attack and normal from admin IPs) -
\textbf{Size:} 15,184 samples - \textbf{Composition:} 3,796 normal
(25.0\%) + 11,388 attack (75.0\%) - \textbf{Normal-to-attack ratio:}
Approximately 1:3 - \textbf{Purpose:} Evaluate model performance on a
realistic mixture of normal and attack traffic

{[}Table 3.6: Complete dataset split summary{]}

{\def\LTcaptype{none} % do not increment counter
\begin{longtable}[]{@{}
  >{\raggedright\arraybackslash}p{(\linewidth - 10\tabcolsep) * \real{0.1525}}
  >{\raggedright\arraybackslash}p{(\linewidth - 10\tabcolsep) * \real{0.2373}}
  >{\raggedright\arraybackslash}p{(\linewidth - 10\tabcolsep) * \real{0.1356}}
  >{\raggedright\arraybackslash}p{(\linewidth - 10\tabcolsep) * \real{0.1356}}
  >{\raggedright\arraybackslash}p{(\linewidth - 10\tabcolsep) * \real{0.1695}}
  >{\raggedright\arraybackslash}p{(\linewidth - 10\tabcolsep) * \real{0.1695}}@{}}
\toprule\noalign{}
\begin{minipage}[b]{\linewidth}\raggedright
Dataset
\end{minipage} & \begin{minipage}[b]{\linewidth}\raggedright
Total Samples
\end{minipage} & \begin{minipage}[b]{\linewidth}\raggedright
Normal
\end{minipage} & \begin{minipage}[b]{\linewidth}\raggedright
Attack
\end{minipage} & \begin{minipage}[b]{\linewidth}\raggedright
Normal \%
\end{minipage} & \begin{minipage}[b]{\linewidth}\raggedright
Attack \%
\end{minipage} \\
\midrule\noalign{}
\endhead
\bottomrule\noalign{}
\endlastfoot
Training & 7,212 & 7,212 & 0 & 100\% & 0\% \\
Test & 15,184 & 3,796 & 11,388 & 25.0\% & 75.0\% \\
\textbf{Total} & \textbf{22,396} & \textbf{11,008} & \textbf{11,388} &
\textbf{49.2\%} & \textbf{50.8\%} \\
\end{longtable}
}

The chronological split for the simulation data (first 70\% for
training, last 30\% for testing) is a critical methodological choice.
Unlike random splitting, which would allow information from later time
periods to ``leak'' into the training set, chronological splitting
ensures that the model is evaluated on data from a time period it has
never seen during training. This simulates the real-world deployment
scenario where the model is trained on historical data and must
generalize to future observations {[}69{]}.

\subsection{3.5.6 Stage 5: Feature
Scaling}\label{stage-5-feature-scaling}

After extraction and splitting, features are normalized using
\textbf{RobustScaler} from the scikit-learn library {[}63{]}.
RobustScaler was selected over StandardScaler (z-score normalization)
and MinMaxScaler based on the following theoretical and empirical
considerations:

\textbf{Theoretical justification.} RobustScaler uses the median and
interquartile range (IQR) as its centering and scaling statistics,
respectively, rather than the mean and standard deviation used by
StandardScaler. The median and IQR are robust estimators of location and
scale, respectively, with a breakdown point of 25\% (compared to 0\% for
the mean and standard deviation) {[}70{]}. This means that up to 25\% of
the data can be arbitrarily corrupted without affecting the scaling
parameters.

\textbf{Empirical justification.} The SSH feature data contains extreme
outlier values, particularly in count-based features. For example,
\texttt{fail\_count} ranges from 0 (typical normal value) to several
thousand (during intense brute-force attacks). StandardScaler would be
heavily influenced by these extreme values, compressing the normal data
into a narrow range around zero. RobustScaler, by using the median and
IQR, preserves the resolution of normal data while gracefully handling
extreme attack values.

The RobustScaler transformation is defined as:

\[x_{\text{scaled}} = \frac{x - Q_2(x)}{Q_3(x) - Q_1(x)}\]

where \(Q_1(x)\), \(Q_2(x)\), and \(Q_3(x)\) are the 25th percentile
(first quartile), 50th percentile (median), and 75th percentile (third
quartile) of feature \(x\), respectively. The scaler is fit exclusively
on the training set (containing only normal data) and the learned
parameters (\(Q_1\), \(Q_2\), \(Q_3\)) are applied to transform both the
training and test sets. This procedure prevents data leakage: the
scaling parameters are derived solely from normal data, ensuring that
the model has no implicit access to the statistical properties of attack
data during training {[}71{]}.

\section{3.6 Feature Engineering and
Selection}\label{feature-engineering-and-selection}

\subsection{3.6.1 Feature Design
Philosophy}\label{feature-design-philosophy}

The 14 features were designed following the principle of behavioral
profiling: each feature captures a specific dimension of SSH
authentication behavior that is expected to differ systematically
between legitimate users and automated attack tools. The feature set was
informed by three bodies of prior work:

\begin{enumerate}
\def\labelenumi{\arabic{enumi}.}
\item
  \textbf{SSH behavioral analysis.} Javed and Paxson {[}72{]}
  demonstrated that timing features (inter-arrival times, session
  durations) are the most reliable discriminators between automated SSH
  attacks and human users, because timing is determined by the
  fundamental difference between human cognition speed and machine
  execution speed.
\item
  \textbf{Network intrusion detection feature engineering.} The features
  used in benchmark datasets such as NSL-KDD {[}73{]}, CICIDS2017
  {[}74{]}, and UNSW-NB15 {[}75{]} informed the selection of count-based
  and ratio-based features (e.g., fail\_rate, invalid\_user\_ratio).
\item
  \textbf{Anomaly detection theory.} Aggarwal {[}76{]} established that
  effective anomaly detection features should exhibit high variance
  between normal and anomalous classes (inter-class variance) and low
  variance within the normal class (intra-class variance), enabling
  clear separation.
\end{enumerate}

\subsection{3.6.2 Feature Definitions}\label{feature-definitions}

The 14 features are organized into five functional groups:

\textbf{Group 1: Authentication Volume Features}

\textbf{1. fail\_count} --- The total number of ``Failed password''
events from IP address \(ip\) within window \([t-W, t]\).

\[\texttt{fail\_count}_{ip,t} = |\{e \in \mathcal{E}_{ip,t} \mid e.\text{type} = \texttt{FAILED\_PASSWORD}\}|\]

This is the most direct indicator of brute-force activity. Normal users
typically produce 0--2 failures per window (due to occasional
mistyping), while active brute-force attacks produce tens to thousands
{[}77{]}.

\textbf{2. success\_count} --- The total number of ``Accepted'' events
(password or publickey).

\[\texttt{success\_count}_{ip,t} = |\{e \in \mathcal{E}_{ip,t} \mid e.\text{type} \in \{\texttt{ACCEPTED\_PASSWORD}, \texttt{ACCEPTED\_PUBLICKEY}\}\}|\]

Legitimate users have a high success rate; brute-force attackers have
near-zero success rates until a credential is cracked.

\textbf{3. fail\_rate} --- The ratio of failed attempts to total
authentication attempts, providing a normalized measure independent of
absolute volume.

\[\texttt{fail\_rate}_{ip,t} = \frac{\texttt{fail\_count}_{ip,t}}{\texttt{fail\_count}_{ip,t} + \texttt{success\_count}_{ip,t} + \epsilon}\]

where \(\epsilon = 10^{-10}\) prevents division by zero. Values
approaching 1.0 are characteristic of brute-force attacks; normal users
typically exhibit fail\_rate \(< 0.3\) {[}78{]}.

\textbf{Group 2: Username Diversity Features}

\textbf{4. unique\_usernames} --- The number of distinct usernames
attempted.

\[\texttt{unique\_usernames}_{ip,t} = |\{e.\text{username} \mid e \in \mathcal{E}_{ip,t}\}|\]

Credential stuffing and dictionary attacks typically enumerate many
usernames (root, admin, test, oracle, etc.), while legitimate users use
only 1--2 accounts {[}79{]}.

\textbf{5. invalid\_user\_count} --- The count of attempts with
usernames that do not exist on the target system.

\[\texttt{invalid\_user\_count}_{ip,t} = |\{e \in \mathcal{E}_{ip,t} \mid e.\text{type} = \texttt{INVALID\_USER}\}|\]

This is a strong attack indicator because legitimate users know their
own account names. The presence of invalid user attempts implies
external probing.

\textbf{6. invalid\_user\_ratio} --- The proportion of invalid user
attempts among all authentication events.

\[\texttt{invalid\_user\_ratio}_{ip,t} = \frac{\texttt{invalid\_user\_count}_{ip,t}}{|\mathcal{E}_{ip,t}| + \epsilon}\]

A high ratio indicates systematic username enumeration, a characteristic
of dictionary and credential stuffing attacks {[}80{]}.

\textbf{Group 3: Temporal Features}

\textbf{7. connection\_count} --- The total number of SSH connections
(distinct events) from the IP within the window.

\[\texttt{connection\_count}_{ip,t} = |\mathcal{E}_{ip,t}|\]

Automated tools generate orders of magnitude more connections than human
users within a given time window.

\textbf{8. mean\_inter\_attempt\_time} (seconds) --- The mean time
interval between consecutive authentication events from the same IP.

\[\texttt{mean\_iat}_{ip,t} = \frac{1}{n-1} \sum_{i=2}^{n} (t_i - t_{i-1})\]

where \(t_1, t_2, \ldots, t_n\) are the sorted timestamps of events from
IP \(ip\) in the window. Automated tools produce sub-second intervals;
humans produce intervals of 5--60+ seconds {[}72{]}.

\textbf{9. std\_inter\_attempt\_time} (seconds) --- The standard
deviation of inter-attempt times.

\[\texttt{std\_iat}_{ip,t} = \sqrt{\frac{1}{n-2} \sum_{i=2}^{n} \left((t_i - t_{i-1}) - \texttt{mean\_iat}_{ip,t}\right)^2}\]

Low standard deviation indicates uniform, machine-like behavior; high
standard deviation indicates irregular, human-like behavior {[}81{]}.

\textbf{10. min\_inter\_attempt\_time} (seconds) --- The minimum time
interval between any two consecutive events.

\[\texttt{min\_iat}_{ip,t} = \min_{i \in \{2, \ldots, n\}} (t_i - t_{i-1})\]

A value near zero indicates that at least one pair of events occurred in
rapid succession, a hallmark of automated attacks. This feature is
particularly sensitive because even if the mean interval is moderate (as
in a low-and-slow attack), a single rapid burst will drive min\_iat to
near zero.

\textbf{Group 4: Network and Session Features}

\textbf{11. unique\_ports} --- The number of distinct source ports used.

\[\texttt{unique\_ports}_{ip,t} = |\{e.\text{port} \mid e \in \mathcal{E}_{ip,t}\}|\]

Each new TCP connection uses a different ephemeral source port. A large
number of unique ports correlates directly with a large number of
connections, providing an independent measure of connection volume that
is robust to event-level deduplication {[}82{]}.

\textbf{12. session\_duration\_mean} (seconds) --- The mean duration of
SSH sessions, computed as the time between ``Accepted'' events and
subsequent ``Connection closed'' events for the same IP and port.

\[\texttt{session\_duration\_mean}_{ip,t} = \frac{1}{k} \sum_{j=1}^{k} (t_{\text{close},j} - t_{\text{open},j})\]

Legitimate sessions last minutes to hours; brute-force sessions last
under 5 seconds because each failed authentication results in rapid
disconnection {[}72{]}. Feature importance analysis (Section 4.4)
confirms this as the single most discriminative feature.

\textbf{Group 5: Attack Indicator Features}

\textbf{13. pam\_failure\_escalation} --- A binary indicator (0 or 1)
denoting whether an escalating sequence of PAM (Pluggable Authentication
Module) authentication failures was observed within the window.

\[\texttt{pam\_failure\_escalation}_{ip,t} = \begin{cases} 1 & \text{if escalating PAM failure sequence detected} \\ 0 & \text{otherwise} \end{cases}\]

An escalating sequence of PAM failures indicates systematic credential
testing, as each successive failure triggers increasingly severe PAM
logging {[}83{]}.

\textbf{14. max\_retries\_exceeded} --- A binary indicator denoting
whether a ``maximum authentication attempts exceeded'' event was logged.

\[\texttt{max\_retries\_exceeded}_{ip,t} = \begin{cases} 1 & \text{if max retries event detected} \\ 0 & \text{otherwise} \end{cases}\]

This event occurs when an attacker exhausts the MaxAuthTries limit
within a single SSH connection, a direct indicator of brute-force
behavior within a session.

\subsection{3.6.3 Feature Selection
Rationale}\label{feature-selection-rationale}

{[}Table 3.7: Summary of 14 features, data types, expected value ranges,
and detection significance{]}

{\def\LTcaptype{none} % do not increment counter
\begin{longtable}[]{@{}
  >{\raggedright\arraybackslash}p{(\linewidth - 10\tabcolsep) * \real{0.0820}}
  >{\raggedright\arraybackslash}p{(\linewidth - 10\tabcolsep) * \real{0.1475}}
  >{\raggedright\arraybackslash}p{(\linewidth - 10\tabcolsep) * \real{0.0984}}
  >{\raggedright\arraybackslash}p{(\linewidth - 10\tabcolsep) * \real{0.2131}}
  >{\raggedright\arraybackslash}p{(\linewidth - 10\tabcolsep) * \real{0.2131}}
  >{\raggedright\arraybackslash}p{(\linewidth - 10\tabcolsep) * \real{0.2459}}@{}}
\toprule\noalign{}
\begin{minipage}[b]{\linewidth}\raggedright
No.
\end{minipage} & \begin{minipage}[b]{\linewidth}\raggedright
Feature
\end{minipage} & \begin{minipage}[b]{\linewidth}\raggedright
Type
\end{minipage} & \begin{minipage}[b]{\linewidth}\raggedright
Normal Range
\end{minipage} & \begin{minipage}[b]{\linewidth}\raggedright
Attack Range
\end{minipage} & \begin{minipage}[b]{\linewidth}\raggedright
Detection Role
\end{minipage} \\
\midrule\noalign{}
\endhead
\bottomrule\noalign{}
\endlastfoot
1 & fail\_count & Integer & 0--2 & 10--1000+ & Volume indicator \\
2 & success\_count & Integer & 1--10 & 0--1 & Success pattern \\
3 & fail\_rate & Float {[}0,1{]} & 0--0.3 & 0.9--1.0 & Normalized
failure \\
4 & unique\_usernames & Integer & 1--2 & 5--50+ & Username diversity \\
5 & invalid\_user\_count & Integer & 0 & 5--500+ & Enumeration
indicator \\
6 & invalid\_user\_ratio & Float {[}0,1{]} & 0 & 0.3--1.0 & Enumeration
severity \\
7 & connection\_count & Integer & 1--5 & 20--1000+ & Activity volume \\
8 & mean\_inter\_attempt\_time & Float (s) & 10--300+ & 0.01--2 & Speed
signature \\
9 & std\_inter\_attempt\_time & Float (s) & 5--200+ & 0--1 & Regularity
signature \\
10 & min\_inter\_attempt\_time & Float (s) & 3--60+ & 0--0.5 & Burst
detection \\
11 & unique\_ports & Integer & 1--5 & 20--500+ & Connection diversity \\
12 & session\_duration\_mean & Float (s) & 300--7200+ & 0.5--5 & Session
pattern \\
13 & pam\_failure\_escalation & Binary & 0 & 0 or 1 & PAM escalation \\
14 & max\_retries\_exceeded & Binary & 0 & 0 or 1 & Retry exhaustion \\
\end{longtable}
}

All 14 features were retained despite some exhibiting high pairwise
correlation (e.g., fail\_count and connection\_count, \(r > 0.8\)), for
three reasons: (1) Isolation Forest uses random feature subsets per
tree, making it naturally robust to correlated features {[}84{]}; (2)
each feature captures a unique behavioral dimension that contributes to
detecting specific attack types; and (3) Aggarwal {[}76{]} demonstrated
that removing correlated features can degrade anomaly detection
performance when the anomalies manifest in specific feature subsets.

\section{3.7 Anomaly Detection Model
Configuration}\label{anomaly-detection-model-configuration}

\subsection{3.7.1 Model Selection
Criteria}\label{model-selection-criteria}

Three anomaly detection models were selected for comparative evaluation
based on a systematic review of the anomaly detection literature {[}47,
48, 76{]}. The selection criteria were: (1) native support for the
semi-supervised (novelty detection) paradigm; (2) demonstrated
effectiveness on tabular numerical data; (3) availability of
production-quality implementations in scikit-learn; and (4) diversity of
algorithmic principles (isolation-based, density-based, boundary-based)
to provide complementary perspectives on the data.

\subsection{3.7.2 Isolation Forest (IF)}\label{isolation-forest-if}

\textbf{Algorithmic principle.} Isolation Forest, proposed by Liu, Ting,
and Zhou {[}84{]}, is based on the principle that anomalous points are
easier to isolate than normal points through random recursive
partitioning. The algorithm constructs an ensemble of \(T\) isolation
trees (iTrees), each built on a random sub-sample of size \(\psi\) from
the training data. Each iTree recursively selects a random feature \(q\)
and a random split value \(p \in [\min(q), \max(q)]\), partitioning the
data until each point is isolated (in its own leaf node) or the maximum
depth \(\lceil \log_2 \psi \rceil\) is reached.

\textbf{Anomaly score.} The anomaly score of a test point \(\mathbf{x}\)
is computed from the average path length \(E[h(\mathbf{x})]\) across all
\(T\) trees:

\[s(\mathbf{x}, n) = 2^{-\frac{E[h(\mathbf{x})]}{c(n)}}\]

where \(c(n)\) is the average path length of an unsuccessful search in a
Binary Search Tree (BST) with \(n\) nodes:

\[c(n) = 2H(n-1) - \frac{2(n-1)}{n}\]

and \(H(i) = \ln(i) + \gamma\) is the harmonic number, with
\(\gamma \approx 0.5772\) being the Euler--Mascheroni constant. Scores
close to 1.0 indicate anomalies (short path lengths, easy to isolate);
scores close to 0.5 indicate normal points.

\textbf{Computational complexity.} Training:
\(O(T \cdot \psi \cdot \log \psi)\). Prediction:
\(O(T \cdot \log \psi)\) per sample. This linear-logarithmic complexity
is a key advantage for real-time deployment {[}84{]}.

\textbf{Hyperparameters (baseline):}

{\def\LTcaptype{none} % do not increment counter
\begin{longtable}[]{@{}
  >{\raggedright\arraybackslash}p{(\linewidth - 4\tabcolsep) * \real{0.3333}}
  >{\raggedright\arraybackslash}p{(\linewidth - 4\tabcolsep) * \real{0.2121}}
  >{\raggedright\arraybackslash}p{(\linewidth - 4\tabcolsep) * \real{0.4545}}@{}}
\toprule\noalign{}
\begin{minipage}[b]{\linewidth}\raggedright
Parameter
\end{minipage} & \begin{minipage}[b]{\linewidth}\raggedright
Value
\end{minipage} & \begin{minipage}[b]{\linewidth}\raggedright
Justification
\end{minipage} \\
\midrule\noalign{}
\endhead
\bottomrule\noalign{}
\endlastfoot
n\_estimators (\(T\)) & 300 & Sufficient for stable score estimation
{[}84{]} \\
max\_samples (\(\psi\)) & 512 & Recommended by Liu et al.~{[}84{]} for
sub-sampling effectiveness \\
max\_features & 0.5 (7/14 features per tree) & Increases ensemble
diversity {[}85{]} \\
contamination & auto & Default scikit-learn behavior \\
\end{longtable}
}

\textbf{Hyperparameters (optimized):}

{\def\LTcaptype{none} % do not increment counter
\begin{longtable}[]{@{}
  >{\raggedright\arraybackslash}p{(\linewidth - 4\tabcolsep) * \real{0.3333}}
  >{\raggedright\arraybackslash}p{(\linewidth - 4\tabcolsep) * \real{0.2121}}
  >{\raggedright\arraybackslash}p{(\linewidth - 4\tabcolsep) * \real{0.4545}}@{}}
\toprule\noalign{}
\begin{minipage}[b]{\linewidth}\raggedright
Parameter
\end{minipage} & \begin{minipage}[b]{\linewidth}\raggedright
Value
\end{minipage} & \begin{minipage}[b]{\linewidth}\raggedright
Justification
\end{minipage} \\
\midrule\noalign{}
\endhead
\bottomrule\noalign{}
\endlastfoot
n\_estimators (\(T\)) & 500 & Increased for greater score stability \\
max\_samples (\(\psi\)) & 512 & Retained from baseline \\
max\_features & 0.75 (10--11/14 features per tree) & Improved per-tree
detection coverage \\
contamination & 0.01 & Explicit 1\% outlier tolerance {[}86{]} \\
\end{longtable}
}

\subsection{3.7.3 Local Outlier Factor
(LOF)}\label{local-outlier-factor-lof-1}

\textbf{Algorithmic principle.} LOF, proposed by Breunig et
al.~{[}87{]}, is a local density-based anomaly detection method. For
each point \(\mathbf{x}\), LOF computes the ratio of the average local
reachability density (lrd) of its \(k\)-nearest neighbors to its own
lrd. The key insight is that an anomalous point resides in a region of
lower density compared to its neighbors.

The \textbf{reachability distance} of point \(\mathbf{x}\) with respect
to point \(\mathbf{o}\) is:

\[\text{reach-dist}_k(\mathbf{x}, \mathbf{o}) = \max\{d_k(\mathbf{o}), d(\mathbf{x}, \mathbf{o})\}\]

where \(d_k(\mathbf{o})\) is the distance from \(\mathbf{o}\) to its
\(k\)-th nearest neighbor and \(d(\mathbf{x}, \mathbf{o})\) is the
Euclidean distance between \(\mathbf{x}\) and \(\mathbf{o}\).

The \textbf{local reachability density} of \(\mathbf{x}\) is:

\[\text{lrd}_k(\mathbf{x}) = \left(\frac{\sum_{\mathbf{o} \in N_k(\mathbf{x})} \text{reach-dist}_k(\mathbf{x}, \mathbf{o})}{|N_k(\mathbf{x})|}\right)^{-1}\]

The \textbf{LOF score} is:

\[\text{LOF}_k(\mathbf{x}) = \frac{\sum_{\mathbf{o} \in N_k(\mathbf{x})} \frac{\text{lrd}_k(\mathbf{o})}{\text{lrd}_k(\mathbf{x})}}{|N_k(\mathbf{x})|}\]

LOF \(\approx 1\) indicates normal density; LOF \(\gg 1\) indicates
anomalous (low) density.

\textbf{Computational complexity.} Training: \(O(n^2 \log n)\) for
\(k\)-nearest neighbor computation using KD-trees. Prediction:
\(O(n \log n)\) per sample (requires comparison against stored training
data).

\textbf{Hyperparameters:}

{\def\LTcaptype{none} % do not increment counter
\begin{longtable}[]{@{}
  >{\raggedright\arraybackslash}p{(\linewidth - 4\tabcolsep) * \real{0.3333}}
  >{\raggedright\arraybackslash}p{(\linewidth - 4\tabcolsep) * \real{0.2121}}
  >{\raggedright\arraybackslash}p{(\linewidth - 4\tabcolsep) * \real{0.4545}}@{}}
\toprule\noalign{}
\begin{minipage}[b]{\linewidth}\raggedright
Parameter
\end{minipage} & \begin{minipage}[b]{\linewidth}\raggedright
Value
\end{minipage} & \begin{minipage}[b]{\linewidth}\raggedright
Justification
\end{minipage} \\
\midrule\noalign{}
\endhead
\bottomrule\noalign{}
\endlastfoot
n\_neighbors (\(k\)) & 30 & Balances sensitivity and stability
{[}87{]} \\
novelty & True & Enables prediction on new (unseen) data {[}63{]} \\
metric & Minkowski (\(p=2\), Euclidean) & Standard distance metric for
numerical features \\
\end{longtable}
}

\subsection{3.7.4 One-Class SVM (OCSVM)}\label{one-class-svm-ocsvm}

\textbf{Algorithmic principle.} OCSVM, proposed by Scholkopf et
al.~{[}88{]}, extends the Support Vector Machine framework to one-class
classification. The algorithm finds a hyperplane in the kernel-mapped
feature space \(\mathcal{H}\) that maximally separates the training data
from the origin, with a soft margin controlled by the parameter \(\nu\).
The optimization problem is:

\[\min_{\mathbf{w}, \xi_i, \rho} \frac{1}{2} \|\mathbf{w}\|^2 + \frac{1}{\nu n} \sum_{i=1}^{n} \xi_i - \rho\]

subject to:

\[\mathbf{w} \cdot \Phi(\mathbf{x}_i) \geq \rho - \xi_i, \quad \xi_i \geq 0, \quad i = 1, \ldots, n\]

where \(\Phi(\cdot)\) is the feature mapping induced by the kernel
function \(K\), and \(\nu \in (0, 1]\) controls the upper bound on the
fraction of training errors and the lower bound on the fraction of
support vectors.

The \textbf{RBF kernel} is used:

\[K(\mathbf{x}_i, \mathbf{x}_j) = \exp\left(-\gamma \|\mathbf{x}_i - \mathbf{x}_j\|^2\right)\]

where \(\gamma\) controls the kernel bandwidth. A new point
\(\mathbf{x}\) is classified as normal if
\(\text{sign}(\mathbf{w} \cdot \Phi(\mathbf{x}) - \rho) = +1\) and
anomalous otherwise.

\textbf{Computational complexity.} Training: \(O(n^2)\) to \(O(n^3)\)
depending on the solver and kernel cache efficiency. Prediction:
\(O(n_{sv})\) per sample, where \(n_{sv}\) is the number of support
vectors.

\textbf{Hyperparameters:}

{\def\LTcaptype{none} % do not increment counter
\begin{longtable}[]{@{}
  >{\raggedright\arraybackslash}p{(\linewidth - 4\tabcolsep) * \real{0.3333}}
  >{\raggedright\arraybackslash}p{(\linewidth - 4\tabcolsep) * \real{0.2121}}
  >{\raggedright\arraybackslash}p{(\linewidth - 4\tabcolsep) * \real{0.4545}}@{}}
\toprule\noalign{}
\begin{minipage}[b]{\linewidth}\raggedright
Parameter
\end{minipage} & \begin{minipage}[b]{\linewidth}\raggedright
Value
\end{minipage} & \begin{minipage}[b]{\linewidth}\raggedright
Justification
\end{minipage} \\
\midrule\noalign{}
\endhead
\bottomrule\noalign{}
\endlastfoot
kernel & RBF & Captures nonlinear decision boundaries {[}88{]} \\
gamma & auto (\(= 1/d = 1/14 \approx 0.0714\)) & scikit-learn default
heuristic \\
nu (\(\nu\)) & 0.01 & Upper bound on outlier fraction (1\%) \\
\end{longtable}
}

\subsection{3.7.5 Model Comparison
Summary}\label{model-comparison-summary}

{[}Table 3.8: Comparative characteristics of the three anomaly detection
models{]}

{\def\LTcaptype{none} % do not increment counter
\begin{longtable}[]{@{}
  >{\raggedright\arraybackslash}p{(\linewidth - 6\tabcolsep) * \real{0.3409}}
  >{\raggedright\arraybackslash}p{(\linewidth - 6\tabcolsep) * \real{0.3864}}
  >{\raggedright\arraybackslash}p{(\linewidth - 6\tabcolsep) * \real{0.1136}}
  >{\raggedright\arraybackslash}p{(\linewidth - 6\tabcolsep) * \real{0.1591}}@{}}
\toprule\noalign{}
\begin{minipage}[b]{\linewidth}\raggedright
Characteristic
\end{minipage} & \begin{minipage}[b]{\linewidth}\raggedright
Isolation Forest
\end{minipage} & \begin{minipage}[b]{\linewidth}\raggedright
LOF
\end{minipage} & \begin{minipage}[b]{\linewidth}\raggedright
OCSVM
\end{minipage} \\
\midrule\noalign{}
\endhead
\bottomrule\noalign{}
\endlastfoot
Principle & Random isolation & Local density ratio & Maximum-margin
hyperplane \\
Training complexity & \(O(T \psi \log \psi)\) & \(O(n^2 \log n)\) &
\(O(n^2)\) to \(O(n^3)\) \\
Prediction complexity & \(O(T \log \psi)\) & \(O(n \log n)\) &
\(O(n_{sv})\) \\
Score type & Continuous {[}0, 1{]} & Continuous \(\geq 1\) & Signed
distance \\
Large-scale suitability & Excellent & Moderate & Moderate \\
Local anomaly detection & Moderate & Excellent & Good \\
Distributional assumptions & None & None & None (with kernel) \\
Interpretability & Moderate (path length) & Low & Low \\
Real-time suitability & Excellent & Limited & Good \\
\end{longtable}
}

\subsection{3.7.6 Hyperparameter
Optimization}\label{hyperparameter-optimization}

Hyperparameter optimization was performed using grid search with a
custom evaluation protocol adapted for the semi-supervised setting.
Since the training set contains only normal data, standard
cross-validation with accuracy as the criterion is not applicable.
Instead, the following procedure was employed:

\begin{enumerate}
\def\labelenumi{\arabic{enumi}.}
\tightlist
\item
  A small validation set was constructed by holding out 10\% of the
  training set (normal data) and combining it with a stratified random
  sample of 500 attack samples from the test set.
\item
  Grid search was performed over the parameter spaces defined for each
  model (see Sections 3.7.2--3.7.4).
\item
  The optimization criterion was the F1-score on the validation set,
  which balances precision and recall.
\item
  The best hyperparameters were selected and the model was retrained on
  the full training set.
\end{enumerate}

For Isolation Forest, the grid search explored:
\(n\_estimators \in \{100, 200, 300, 500\}\),
\(max\_samples \in \{256, 512, 1024\}\),
\(max\_features \in \{0.25, 0.5, 0.75, 1.0\}\), and
\(contamination \in \{\text{auto}, 0.01, 0.05, 0.1\}\), totaling
\(4 \times 3 \times 4 \times 4 = 192\) parameter combinations.

\section{3.8 Dynamic Threshold Engine
Design}\label{dynamic-threshold-engine-design}

\subsection{3.8.1 Motivation and Theoretical
Basis}\label{motivation-and-theoretical-basis}

Static thresholds --- fixed anomaly score cutoff values determined a
priori --- are the standard approach in most deployed anomaly detection
systems {[}89{]}. However, static thresholds suffer from a fundamental
limitation: the optimal threshold depends on the current statistical
characteristics of the data, which change over time. In the SSH
monitoring context, this non-stationarity manifests in several ways:
business hours produce more legitimate login activity (and more
legitimate failures) than off-hours; system maintenance windows produce
unusual but legitimate traffic patterns; and the baseline level of
background attack traffic fluctuates with Internet-wide scanning
campaigns {[}90{]}.

The EWMA-Adaptive Percentile dynamic threshold engine was designed to
address this limitation. The theoretical foundation draws on two
established techniques:

\begin{enumerate}
\def\labelenumi{\arabic{enumi}.}
\item
  \textbf{Exponentially Weighted Moving Average (EWMA).} Originally
  developed by Roberts {[}91{]} for statistical process control in
  manufacturing quality assurance, EWMA provides a smooth estimate of
  the current ``baseline'' of a time series that gives exponentially
  decreasing weight to older observations.
\item
  \textbf{Percentile-based thresholding.} The use of distribution
  percentiles as threshold values is a standard robust statistical
  technique that is less sensitive to outliers than mean-based
  approaches {[}70{]}.
\end{enumerate}

The EWMA-Adaptive Percentile algorithm combines these two techniques
into a threshold that simultaneously tracks the current baseline (via
EWMA) and adapts to the current distributional spread (via percentile).

\subsection{3.8.2 Algorithm
Specification}\label{algorithm-specification}

\textbf{Parameters:}

{\def\LTcaptype{none} % do not increment counter
\begin{longtable}[]{@{}
  >{\raggedright\arraybackslash}p{(\linewidth - 6\tabcolsep) * \real{0.2821}}
  >{\raggedright\arraybackslash}p{(\linewidth - 6\tabcolsep) * \real{0.2051}}
  >{\raggedright\arraybackslash}p{(\linewidth - 6\tabcolsep) * \real{0.1795}}
  >{\raggedright\arraybackslash}p{(\linewidth - 6\tabcolsep) * \real{0.3333}}@{}}
\toprule\noalign{}
\begin{minipage}[b]{\linewidth}\raggedright
Parameter
\end{minipage} & \begin{minipage}[b]{\linewidth}\raggedright
Symbol
\end{minipage} & \begin{minipage}[b]{\linewidth}\raggedright
Value
\end{minipage} & \begin{minipage}[b]{\linewidth}\raggedright
Description
\end{minipage} \\
\midrule\noalign{}
\endhead
\bottomrule\noalign{}
\endlastfoot
Smoothing factor & \(\alpha\) & 0.3 & EWMA weight for new
observations \\
Base percentile & \(p\) & 95 & Percentile of the lookback
distribution \\
Sensitivity factor & \(\lambda\) & 1.5 & Multiplier for the distance
above EWMA \\
Lookback window & \(L\) & 100 & Number of recent scores for percentile
computation \\
\end{longtable}
}

\textbf{Step 1: Update the EWMA of anomaly scores.} Given the anomaly
score time series \(s_1, s_2, \ldots, s_t\) produced by the model, the
EWMA at time \(t\) is updated recursively:

\[\mu_t^{\text{EWMA}} = \alpha \cdot s_t + (1 - \alpha) \cdot \mu_{t-1}^{\text{EWMA}}\]

with initialization \(\mu_0^{\text{EWMA}} = s_1\). The smoothing factor
\(\alpha = 0.3\) was selected following the guidelines of Lucas and
Saccucci {[}92{]}, who showed that \(\alpha \in [0.2, 0.4]\) provides a
good balance between responsiveness to genuine changes and robustness to
transient noise for process control applications. The effective memory
span of an EWMA is approximately \(2/\alpha - 1 = 5.67\) observations,
meaning that observations older than approximately 6 time steps
contribute less than 5\% to the current EWMA value.

\textbf{Step 2: Compute the adaptive percentile.} Using the \(L = 100\)
most recent anomaly scores, compute the \(p\)-th percentile:

\[P_t = \text{Percentile}(\{s_{t-L+1}, s_{t-L+2}, \ldots, s_t\}, p)\]

The base percentile \(p = 95\) ensures that the threshold is placed
above 95\% of recent normal scores, providing a natural separation
between normal variability and genuine anomalies. The lookback window
\(L = 100\) provides sufficient statistical stability for the percentile
estimate while remaining responsive to distributional changes.

\textbf{Step 3: Compute the dynamic threshold.} The threshold at time
\(t\) is:

\[\theta_t = \mu_t^{\text{EWMA}} + \lambda \cdot (P_t - \mu_t^{\text{EWMA}})\]

Substituting \(\lambda = 1.5\):

\[\theta_t = \mu_t^{\text{EWMA}} + 1.5 \cdot (P_t - \mu_t^{\text{EWMA}})\]

This formula places the threshold at the EWMA mean plus 1.5 times the
distance from the EWMA to the 95th percentile. The geometric
interpretation is: \(\theta_t\) interpolates between the EWMA
(\(\lambda = 0\)) and the 95th percentile (\(\lambda = 1\)), and with
\(\lambda = 1.5\), the threshold extends beyond the 95th percentile by
50\% of the EWMA-to-percentile distance. This design provides three
desirable properties:

\begin{itemize}
\tightlist
\item
  During stable normal periods, \(P_t - \mu_t^{\text{EWMA}}\) is small
  (the 95th percentile is close to the mean), so \(\theta_t\) is
  slightly above the mean, providing sensitive detection.
\item
  During attack waves, the EWMA rises (pulled upward by high anomaly
  scores), which raises \(\theta_t\) and prevents the system from
  generating continuous false alarms from the elevated baseline.
\item
  After attacks subside, the EWMA decays exponentially back toward the
  normal level, lowering \(\theta_t\) and restoring detection
  sensitivity.
\end{itemize}

\textbf{Step 4: Two-level classification.} A sample at time \(t\) is
classified into one of three categories:

\[\hat{y}(\mathbf{x}, t) = \begin{cases} \texttt{ALERT} & \text{if } s(\mathbf{x}) > \theta_t \\ \texttt{EARLY\_WARNING} & \text{if } s(\mathbf{x}) > \theta_t / 1.5 \\ \texttt{NORMAL} & \text{otherwise} \end{cases}\]

The EARLY\_WARNING threshold is set at
\(\theta_t / 1.5 \approx 0.667 \cdot \theta_t\), providing an
intermediate alerting level that notifies security personnel without
triggering automated blocking. This graduated response mechanism is
critical for production environments where false IP blocks can disrupt
legitimate services {[}93{]}.

\subsection{3.8.3 Self-Calibration
Mechanism}\label{self-calibration-mechanism}

The dynamic threshold engine incorporates a self-calibration mechanism
that re-estimates the baseline statistics every \(L = 100\) decisions.
At each calibration point, the engine:

\begin{enumerate}
\def\labelenumi{\arabic{enumi}.}
\tightlist
\item
  Recomputes the EWMA baseline from the most recent 100 observations.
\item
  Recalculates the 95th percentile from the same observations.
\item
  Adjusts the threshold if the distributional characteristics have
  shifted significantly.
\end{enumerate}

This self-calibration ensures long-term stability even under persistent
non-stationarity, such as gradual changes in the organization's SSH
usage patterns or seasonal variations in Internet-wide attack traffic.

{[}Table 3.9: Dynamic threshold parameter sensitivity{]}

{\def\LTcaptype{none} % do not increment counter
\begin{longtable}[]{@{}
  >{\raggedright\arraybackslash}p{(\linewidth - 8\tabcolsep) * \real{0.1774}}
  >{\raggedright\arraybackslash}p{(\linewidth - 8\tabcolsep) * \real{0.1774}}
  >{\raggedright\arraybackslash}p{(\linewidth - 8\tabcolsep) * \real{0.1613}}
  >{\raggedright\arraybackslash}p{(\linewidth - 8\tabcolsep) * \real{0.1774}}
  >{\raggedright\arraybackslash}p{(\linewidth - 8\tabcolsep) * \real{0.3065}}@{}}
\toprule\noalign{}
\begin{minipage}[b]{\linewidth}\raggedright
Parameter
\end{minipage} & \begin{minipage}[b]{\linewidth}\raggedright
Low Value
\end{minipage} & \begin{minipage}[b]{\linewidth}\raggedright
Selected
\end{minipage} & \begin{minipage}[b]{\linewidth}\raggedright
High Value
\end{minipage} & \begin{minipage}[b]{\linewidth}\raggedright
Effect of Increase
\end{minipage} \\
\midrule\noalign{}
\endhead
\bottomrule\noalign{}
\endlastfoot
\(\alpha\) & 0.1 (slow) & 0.3 (balanced) & 0.5 (fast) & Faster response,
more noise sensitivity \\
\(p\) & 90th & 95th & 99th & Higher threshold, fewer FP, may miss subtle
attacks \\
\(\lambda\) & 1.0 & 1.5 & 2.0 & Higher threshold, fewer alerts
overall \\
\(L\) & 50 & 100 & 200 & Smoother percentile, slower adaptation \\
\end{longtable}
}

{[}Figure 3.4: Illustration of the EWMA-Adaptive Percentile dynamic
threshold behavior over time, showing adaptation to traffic pattern
changes{]}

\section{3.9 Alert and Prevention
Module}\label{alert-and-prevention-module}

\subsection{3.9.1 Fail2Ban Integration}\label{fail2ban-integration}

When the detection module classifies an IP as ALERT-level (attack), the
system invokes the Fail2Ban API to execute automated prevention. The
integration architecture follows the event-driven pattern described by
Hohpe and Woolf {[}94{]}:

\begin{enumerate}
\def\labelenumi{\arabic{enumi}.}
\tightlist
\item
  \textbf{IP banning.} Fail2Ban adds a DROP rule to the
  iptables/nftables firewall, blocking all new SSH connections from the
  offending IP.
\item
  \textbf{Configurable ban duration.} The base ban duration is
  configurable (default: 600 seconds), with progressive escalation:
  repeat offenders receive exponentially increasing ban durations
  (\(\text{ban\_duration}_k = \text{base} \times 2^{k-1}\), where \(k\)
  is the offense count).
\item
  \textbf{Whitelist protection.} The 6 administrative IP addresses are
  permanently whitelisted and cannot be banned, preventing accidental
  lockout of the research team.
\end{enumerate}

\subsection{3.9.2 Multi-Channel Alert
System}\label{multi-channel-alert-system}

The alert system supports three notification channels:

\begin{itemize}
\tightlist
\item
  \textbf{Dashboard alerts.} Real-time alert display on the React
  dashboard via WebSocket, including IP address, timestamp, anomaly
  score, alert level (EARLY\_WARNING or ALERT), and action taken.
\item
  \textbf{Webhook notifications.} Outbound HTTP POST notifications to
  external services (e.g., Slack, Telegram) for integration with
  existing security operations workflows.
\item
  \textbf{Audit logging.} Comprehensive event logging to Elasticsearch
  for post-incident forensic analysis, compliance reporting, and
  long-term trend analysis.
\end{itemize}

\subsection{3.9.3 Administrator Feedback
Loop}\label{administrator-feedback-loop}

An administrator feedback mechanism allows security personnel to review
detection events through the React interface and mark them as confirmed
attack (true positive), false alarm (false positive), or requires
investigation. This feedback is stored in the database and can be used
for:

\begin{itemize}
\tightlist
\item
  Threshold parameter adjustment (increasing \(\lambda\) if too many
  false positives, decreasing if attacks are missed).
\item
  Feature engineering refinement based on patterns observed in
  misclassified events.
\item
  Long-term model performance monitoring and drift detection.
\end{itemize}

\section{3.10 Visualization Module}\label{visualization-module}

\subsection{3.10.1 React Dashboard}\label{react-dashboard}

The monitoring dashboard provides security operators with situational
awareness through the following visualization components:

\begin{itemize}
\tightlist
\item
  \textbf{Real-time event stream.} A chronological feed of all detection
  events with severity color coding (green for normal, yellow for
  EARLY\_WARNING, red for ALERT).
\item
  \textbf{Anomaly score trend chart.} A time-series plot showing anomaly
  scores and the dynamic threshold over time, enabling operators to
  visually assess the system's decision-making process.
\item
  \textbf{IP geolocation map.} A world map displaying the geographic
  origins of detected attacks, using the MaxMind GeoIP2 database for
  IP-to-location resolution {[}95{]}.
\item
  \textbf{Attack statistics dashboard.} Aggregate statistics including
  total events, detection rate, false positive rate, top attacking IPs,
  and most targeted usernames.
\item
  \textbf{Configuration panel.} Administrative controls for adjusting
  detection parameters (\(\alpha\), \(p\), \(\lambda\), \(L\)), ban
  duration, and notification settings.
\end{itemize}

\subsection{3.10.2 Kibana Integration}\label{kibana-integration}

Kibana provides complementary ad hoc exploration capabilities for
security analysts who need to investigate specific incidents in depth.
Pre-configured Kibana dashboards display: raw log event timelines,
per-IP behavioral profiles over time, feature distribution histograms,
and correlation between detection events and Fail2Ban ban actions.

\section{3.11 Evaluation Metrics}\label{evaluation-metrics}

\subsection{3.11.1 Classification Metrics}\label{classification-metrics}

Performance evaluation is based on the standard binary classification
confusion matrix, where positive denotes the attack class and negative
denotes the normal class. The five metrics used in this research, and
their mathematical definitions, are:

\textbf{Accuracy.} The proportion of all samples correctly classified:

\[\text{Accuracy} = \frac{TP + TN}{TP + TN + FP + FN}\]

In imbalanced datasets, accuracy can be misleading --- a classifier that
always predicts the majority class achieves high accuracy without any
discriminative power {[}96{]}. Therefore, accuracy is reported but not
used as the primary evaluation criterion.

\textbf{Precision.} The proportion of predicted attacks that are actual
attacks:

\[\text{Precision} = \frac{TP}{TP + FP}\]

High precision indicates few false alarms. In cybersecurity operations,
excessive false positives cause ``alert fatigue,'' leading operators to
ignore genuine alerts {[}97{]}.

\textbf{Recall (Sensitivity, True Positive Rate).} The proportion of
actual attacks that are correctly detected:

\[\text{Recall} = \frac{TP}{TP + FN}\]

In cybersecurity, recall is the paramount metric because a missed attack
(false negative) can result in system compromise, data exfiltration, or
service disruption --- consequences far more severe than a false alarm
{[}98{]}.

\textbf{F1-Score.} The harmonic mean of precision and recall:

\[F_1 = 2 \cdot \frac{\text{Precision} \times \text{Recall}}{\text{Precision} + \text{Recall}}\]

The harmonic mean penalizes extreme imbalances between precision and
recall, providing a balanced single-number summary of classification
performance {[}99{]}.

\textbf{ROC-AUC (Area Under the Receiver Operating Characteristic
Curve).} The probability that the model assigns a higher anomaly score
to a randomly chosen attack sample than to a randomly chosen normal
sample:

\[\text{ROC-AUC} = P(s(\mathbf{x}^+) > s(\mathbf{x}^-))\]

where \(\mathbf{x}^+\) is drawn from the attack class and
\(\mathbf{x}^-\) from the normal class. ROC-AUC = 1.0 indicates perfect
discrimination; ROC-AUC = 0.5 indicates random (no discriminative power)
{[}100{]}.

\textbf{False Positive Rate (FPR).} The proportion of normal samples
incorrectly classified as attacks:

\[\text{FPR} = \frac{FP}{FP + TN}\]

FPR is reported alongside the primary metrics because it directly
quantifies the operational cost of false alarms in terms of wrongly
banned legitimate users.

\subsection{3.11.2 Priority Hierarchy}\label{priority-hierarchy}

{[}Table 3.10: Evaluation metrics and their security significance{]}

{\def\LTcaptype{none} % do not increment counter
\begin{longtable}[]{@{}
  >{\raggedright\arraybackslash}p{(\linewidth - 4\tabcolsep) * \real{0.2000}}
  >{\raggedright\arraybackslash}p{(\linewidth - 4\tabcolsep) * \real{0.5500}}
  >{\raggedright\arraybackslash}p{(\linewidth - 4\tabcolsep) * \real{0.2500}}@{}}
\toprule\noalign{}
\begin{minipage}[b]{\linewidth}\raggedright
Metric
\end{minipage} & \begin{minipage}[b]{\linewidth}\raggedright
Security Significance
\end{minipage} & \begin{minipage}[b]{\linewidth}\raggedright
Priority
\end{minipage} \\
\midrule\noalign{}
\endhead
\bottomrule\noalign{}
\endlastfoot
Recall & Not missing attacks (preventing compromise) & Highest \\
F1-Score & Balanced precision-recall performance & High \\
Precision & Reducing false alarms (preventing alert fatigue) & High \\
ROC-AUC & General discriminative ability across all thresholds & High \\
Accuracy & Overall correctness & Medium \\
FPR & Operational cost of false bans & Medium \\
\end{longtable}
}

The priority hierarchy reflects the asymmetric cost structure in
cybersecurity: the cost of missing an attack (potential system
compromise) is orders of magnitude greater than the cost of a false
alarm (temporary IP ban that can be appealed) {[}101{]}.

\subsection{3.11.3 System-Level Metrics}\label{system-level-metrics}

In addition to classification metrics, the following system-level
metrics are evaluated:

\begin{itemize}
\tightlist
\item
  \textbf{End-to-end latency:} Time from SSH log event occurrence to
  classification decision and response action.
\item
  \textbf{Throughput:} Maximum number of feature vectors processed per
  second.
\item
  \textbf{Training time:} Wall-clock time to train each model on the
  training set.
\item
  \textbf{Memory usage:} Peak RAM consumption of each model during
  training and inference.
\end{itemize}

\section{3.12 Limitations of the
Methodology}\label{limitations-of-the-methodology}

\subsection{3.12.1 Window Size
Constraints}\label{window-size-constraints}

The 5-minute sliding window, while effective for detecting rapid and
moderate-speed attacks, imposes an inherent limitation on the detection
of extremely slow attacks where the attacker spaces attempts over
intervals exceeding the window size. An attacker performing one attempt
every 10 minutes would produce at most one event per window, generating
feature vectors that may be indistinguishable from normal activity. This
limitation is partially mitigated by the EWMA accumulation mechanism
(which detects the cumulative trend over multiple windows) but cannot be
fully resolved within the current architecture. Long-term IP profiling,
as proposed by Hofstede et al.~{[}102{]}, or integration with external
threat intelligence feeds could complement the window-based approach.

\subsection{3.12.2 Training Data Purity
Assumption}\label{training-data-purity-assumption}

The semi-supervised paradigm assumes that the training data is purely
normal. Any contamination of the training set with attack samples ---
even a small proportion --- could cause the model to learn attack
patterns as part of the normal profile, degrading detection performance.
While the simulation-based training data generation provides strong
purity guarantees, this assumption may be more difficult to satisfy in
production deployments where the boundary between normal and anomalous
behavior is less clear-cut.

\subsection{3.12.3 Single-Environment
Generalization}\label{single-environment-generalization}

The RobustScaler normalization and model training are based on data from
a single operational environment (the simulation representing one
organization's SSH usage patterns). The learned feature distributions
and decision boundaries may not generalize to environments with
significantly different usage patterns (e.g., a high-security data
center vs.~a university campus). When deploying to a new environment,
the scaler and model should be retrained on data representative of that
environment, following the domain adaptation guidelines of Pan and Yang
{[}103{]}.

\subsection{3.12.4 Dataset Specificity}\label{dataset-specificity}

The evaluation is conducted on a specific dataset with a fixed
normal-to-attack ratio of approximately 1:3 in the test set. Performance
may vary under different class ratios or in environments with
significantly different traffic volumes, attack sophistication levels,
or user population sizes. The dynamic threshold mechanism is designed to
mitigate some of this sensitivity through its adaptive nature, but
comprehensive cross-environment evaluation remains an area for future
work.

\subsection{3.12.5 Ethical Considerations}\label{ethical-considerations}

The data collection methodology was designed with careful attention to
ethical considerations. The honeypot data was collected from a
purpose-built honeypot server designed to attract and record attack
traffic; no legitimate user data was collected or compromised. The 6
administrative IP addresses that produced legitimate traffic on the
honeypot were under the full control of the research team. The
simulation data was generated synthetically, using entirely fictional
user accounts and behavioral profiles. Attack scenario testing was
conducted exclusively within a controlled, isolated test environment,
with no unauthorized access attempts against external systems. The
attack simulation scripts are documented for reproducibility but are not
distributed in executable form to prevent potential misuse, following
the responsible disclosure guidelines of the CERT Coordination Center
{[}104{]}.

\chapter{CHAPTER 4: EXPERIMENTAL RESULTS AND
ANALYSIS}\label{chapter-4-experimental-results-and-analysis}

\section{4.1 Experimental Setup}\label{experimental-setup}

\subsection{4.1.1 Computational
Environment}\label{computational-environment}

All experiments were executed on the Docker-containerized platform
comprising 9 services described in Chapter 3. The machine learning
experiments used Python 3.10 with scikit-learn 1.3.x for model training
and evaluation, pandas 2.1.x for data manipulation, numpy 1.25.x for
numerical computation, and matplotlib 3.8.x with seaborn 0.13.x for
visualization. Model training and hyperparameter search were performed
using scikit-learn's GridSearchCV with the custom semi-supervised
evaluation protocol described in Section 3.7.6. All timing measurements
were obtained using Python's \texttt{time.perf\_counter()} for
sub-millisecond precision, and experiments were repeated three times
with the median result reported to reduce the influence of system-level
variability.

\subsection{4.1.2 Experimental Protocol}\label{experimental-protocol}

The experimental evaluation follows a structured progression designed to
validate each component of the system independently before evaluating
the integrated whole:

\begin{enumerate}
\def\labelenumi{\arabic{enumi}.}
\tightlist
\item
  \textbf{Preprocessing validation} (Section 4.2): Verify the log
  parsing, labeling, feature extraction, and data splitting pipeline.
\item
  \textbf{Feature analysis} (Section 4.3): Examine feature
  distributions, correlations, and discriminative power.
\item
  \textbf{Baseline model evaluation} (Section 4.3): Train and evaluate
  three models with default/baseline hyperparameters.
\item
  \textbf{Feature importance analysis} (Section 4.4): Quantify the
  contribution of each feature to detection performance.
\item
  \textbf{Model optimization} (Section 4.5): Apply non-overlapping
  windows, derived features, and hyperparameter tuning.
\item
  \textbf{Dynamic threshold evaluation} (Section 4.6): Assess the
  EWMA-Adaptive Percentile engine.
\item
  \textbf{System integration test} (Section 4.7): Validate the
  end-to-end pipeline from log ingestion to IP blocking.
\item
  \textbf{Attack scenario testing} (Section 4.8): Evaluate detection
  across five distinct attack types.
\item
  \textbf{AI vs.~Fail2Ban comparison} (Section 4.9): Quantify the
  detection time advantage of the AI system.
\item
  \textbf{State-of-the-art comparison} (Section 4.10): Benchmark against
  10+ published studies.
\end{enumerate}

\subsection{4.1.3 Reproducibility}\label{reproducibility}

To ensure reproducibility, all random seeds were fixed (random\_state=42
for scikit-learn models, numpy.random.seed(42) for data manipulation),
and the complete pipeline configuration (including feature extraction
parameters, model hyperparameters, and threshold settings) was
versioned. The chronological data split described in Section 3.5.5 was
deterministic and parameter-free.

\section{4.2 Preprocessing Results}\label{preprocessing-results}

\subsection{4.2.1 Log Parsing Statistics}\label{log-parsing-statistics}

The log parsing module processed the two raw authentication logs and
produced the following results:

{[}Table 4.1: Log parsing results{]}

{\def\LTcaptype{none} % do not increment counter
\begin{longtable}[]{@{}
  >{\raggedright\arraybackslash}p{(\linewidth - 4\tabcolsep) * \real{0.1111}}
  >{\raggedright\arraybackslash}p{(\linewidth - 4\tabcolsep) * \real{0.4167}}
  >{\raggedright\arraybackslash}p{(\linewidth - 4\tabcolsep) * \real{0.4722}}@{}}
\toprule\noalign{}
\begin{minipage}[b]{\linewidth}\raggedright
Metric
\end{minipage} & \begin{minipage}[b]{\linewidth}\raggedright
Honeypot (honeypot\_auth.log)
\end{minipage} & \begin{minipage}[b]{\linewidth}\raggedright
Simulation (simulation\_auth.log)
\end{minipage} \\
\midrule\noalign{}
\endhead
\bottomrule\noalign{}
\endlastfoot
Raw log lines & 119,729 & 54,521 \\
SSHD entries after parsing & 122,809 & 34,700 \\
``Message repeated'' entries expanded & 1,920 & 0 \\
Lines added by expansion & +3,080 & 0 \\
Non-SSHD lines excluded & N/A & \textasciitilde19,821 \\
Parse success rate & \textgreater99.5\% & \textgreater99.5\% \\
\end{longtable}
}

The honeypot dataset increased from 119,729 raw lines to 122,809 parsed
SSHD entries due to the expansion of 1,920 ``message repeated'' syslog
consolidation entries. Each ``message repeated N times'' entry was
expanded by replicating the preceding event \(N\) times, adding
approximately 3,080 events to the dataset. This expansion is critical
for accurate feature computation: without it, time windows containing
``message repeated'' entries would undercount the true number of
authentication events, leading to artificially low values for
fail\_count, connection\_count, and related features.

The simulation dataset decreased from 54,521 raw lines to 34,700 SSHD
entries because the raw log includes entries from other system services
(cron, systemd, su, sudo) that are not relevant to SSH authentication
analysis. Only lines containing ``sshd'' in the service field were
retained.

\subsection{4.2.2 Labeling Statistics}\label{labeling-statistics}

{[}Table 4.2: Labeling results by data source{]}

{\def\LTcaptype{none} % do not increment counter
\begin{longtable}[]{@{}llll@{}}
\toprule\noalign{}
Category & Honeypot & Simulation & Total \\
\midrule\noalign{}
\endhead
\bottomrule\noalign{}
\endlastfoot
Normal events & 532 (admin IP logins) & 34,700 (all events) & 35,232 \\
Attack events & 122,277 & 0 & 122,277 \\
Total SSHD events & 122,809 & 34,700 & 157,509 \\
\end{longtable}
}

The honeypot data exhibits an extreme imbalance at the event level: only
0.43\% of events are normal (the 532 admin root logins), while 99.57\%
are attack events. This ratio reflects the reality of a public-facing
SSH server, where attack traffic vastly dominates legitimate traffic.
Crucially, this event-level imbalance does not directly translate to the
window-level feature vector imbalance, because the feature extraction
process aggregates many events into each feature vector, and the admin
logins produce relatively few but high-quality feature vectors.

\subsection{4.2.3 Feature Extraction
Statistics}\label{feature-extraction-statistics}

{[}Table 4.3: Feature extraction results (overlapping windows, stride=1
min){]}

{\def\LTcaptype{none} % do not increment counter
\begin{longtable}[]{@{}llll@{}}
\toprule\noalign{}
Source & Feature Vectors & Window Size & Stride \\
\midrule\noalign{}
\endhead
\bottomrule\noalign{}
\endlastfoot
Simulation & 10,304 & 5 minutes & 1 minute \\
Honeypot & 50,792 & 5 minutes & 1 minute \\
\textbf{Total} & \textbf{61,096} & & \\
\end{longtable}
}

{[}Table 4.4: Feature extraction results (non-overlapping windows,
stride=5 min){]}

{\def\LTcaptype{none} % do not increment counter
\begin{longtable}[]{@{}llll@{}}
\toprule\noalign{}
Source & Feature Vectors & Window Size & Stride \\
\midrule\noalign{}
\endhead
\bottomrule\noalign{}
\endlastfoot
Simulation & \textasciitilde2,060 & 5 minutes & 5 minutes \\
Honeypot & \textasciitilde10,158 & 5 minutes & 5 minutes \\
\textbf{Total} & \textasciitilde12,218 & & \\
\end{longtable}
}

The overlapping configuration (stride=1 min) produces approximately 5x
more feature vectors than the non-overlapping configuration (stride=5
min), but consecutive vectors share 80\% of their underlying data,
introducing strong temporal autocorrelation. The implications of this
autocorrelation for model performance are discussed in Section 4.5.

\subsection{4.2.4 Data Split Statistics}\label{data-split-statistics}

{[}Table 4.5: Final dataset split (overlapping windows, baseline
configuration){]}

{\def\LTcaptype{none} % do not increment counter
\begin{longtable}[]{@{}
  >{\raggedright\arraybackslash}p{(\linewidth - 10\tabcolsep) * \real{0.1731}}
  >{\raggedright\arraybackslash}p{(\linewidth - 10\tabcolsep) * \real{0.1346}}
  >{\raggedright\arraybackslash}p{(\linewidth - 10\tabcolsep) * \real{0.1538}}
  >{\raggedright\arraybackslash}p{(\linewidth - 10\tabcolsep) * \real{0.1538}}
  >{\raggedright\arraybackslash}p{(\linewidth - 10\tabcolsep) * \real{0.1923}}
  >{\raggedright\arraybackslash}p{(\linewidth - 10\tabcolsep) * \real{0.1923}}@{}}
\toprule\noalign{}
\begin{minipage}[b]{\linewidth}\raggedright
Dataset
\end{minipage} & \begin{minipage}[b]{\linewidth}\raggedright
Total
\end{minipage} & \begin{minipage}[b]{\linewidth}\raggedright
Normal
\end{minipage} & \begin{minipage}[b]{\linewidth}\raggedright
Attack
\end{minipage} & \begin{minipage}[b]{\linewidth}\raggedright
Normal \%
\end{minipage} & \begin{minipage}[b]{\linewidth}\raggedright
Attack \%
\end{minipage} \\
\midrule\noalign{}
\endhead
\bottomrule\noalign{}
\endlastfoot
Training & 7,212 & 7,212 & 0 & 100.0\% & 0.0\% \\
Test & 15,184 & 3,796 & 11,388 & 25.0\% & 75.0\% \\
\textbf{Total} & \textbf{22,396} & \textbf{11,008} & \textbf{11,388} &
\textbf{49.2\%} & \textbf{50.8\%} \\
\end{longtable}
}

The training set of 7,212 samples represents the first 70\% of the
simulation data (chronological split), containing exclusively normal SSH
behavioral profiles. The test set of 15,184 samples combines the
remaining 30\% of simulation data (3,796 normal samples) with the
honeypot-derived feature vectors (11,388 attack + a small number of
normal from admin IPs), yielding a normal-to-attack ratio of
approximately 1:3.

\subsection{4.2.5 Honeypot Data Exploratory
Analysis}\label{honeypot-data-exploratory-analysis}

Detailed exploratory analysis of the honeypot data reveals the attack
landscape characteristic of a public-facing SSH server:

\textbf{Temporal distribution.} Attacks occurred continuously 24 hours a
day, 7 days a week, with no discernible diurnal or weekly pattern. Over
the 5-day collection period, an average of approximately 5,860 failed
password events were recorded per day (29,301 total / 5 days). This
continuous activity pattern is consistent with the use of automated
attack tools (botnets, scanners) that operate without regard to time
zones or business hours, as documented by Alata et al.~{[}54{]}.

\textbf{IP address distribution.} A total of 679 unique IP addresses
were recorded. The distribution is strongly right-skewed (heavy-tailed):
the top 10\% of IPs (approximately 68 addresses) contributed more than
60\% of total failed attempts, while the bottom 50\% of IPs each
produced fewer than 20 attempts. This Pareto-like distribution is
consistent with the findings of Owens and Matthews {[}14{]}, who
reported that SSH brute-force traffic on the Internet follows a
power-law distribution, with a small number of highly active attackers
and a long tail of opportunistic scanners.

\textbf{Username targeting.} The most frequently targeted username was
``root,'' accounting for an estimated 40\% or more of all failed
password events. This is consistent with the well-documented attacker
preference for the root account, which provides unrestricted system
access upon compromise {[}11{]}. The remaining targeted usernames
followed a predictable distribution: ``admin,'' ``test,'' ``user,''
``oracle,'' ``postgres,'' ``ubuntu,'' ``ftp,'' ``mysql,'' and other
common service account names. The diversity of attempted usernames
(hundreds of unique values) reveals the attacker strategy of combining
targeted high-value accounts with broad enumeration.

\textbf{Administrative activity.} The 6 verified admin IPs generated 532
successful root login sessions, distributed primarily during business
hours (8:00--18:00 local time), with session durations ranging from
several minutes to several hours --- characteristic of legitimate system
administration activity.

{[}Table 4.6: Honeypot data detailed statistics{]}

{\def\LTcaptype{none} % do not increment counter
\begin{longtable}[]{@{}ll@{}}
\toprule\noalign{}
Statistic & Value \\
\midrule\noalign{}
\endhead
\bottomrule\noalign{}
\endlastfoot
Total log lines & 119,729 \\
SSHD entries after parsing & 122,809 \\
Collection period & 2026-03-22 to 2026-03-27 (5 days) \\
Unique IP addresses & 679 \\
Failed password events & 29,301 \\
Invalid user attempts & 12,799 \\
Accepted root logins (admin IPs) & 532 \\
Admin IP count & 6 \\
``Message repeated'' entries & 1,920 \\
Average failed attempts per day & \textasciitilde5,860 \\
Average failed attempts per IP & \textasciitilde43.15 \\
\end{longtable}
}

{[}Figure 4.1: Distribution of failed authentication attempts per IP
(log scale), showing the heavy-tailed distribution characteristic of SSH
brute-force traffic{]}

\subsection{4.2.6 Simulation Data Exploratory
Analysis}\label{simulation-data-exploratory-analysis}

\textbf{Temporal distribution.} Activity was concentrated during
business hours (8:00--18:00), with a marked decrease during evening and
nighttime hours, accurately reflecting the diurnal work patterns of a
medium-sized organization. The 12 service accounts produced periodic,
regular activity throughout the day (cron-based), while the 52 human
accounts exhibited the expected variability in login times and session
durations.

\textbf{Authentication success rate.} With 4,205 successful logins and
only 177 failures, the success rate reached 95.96\% (4,205 / (4,205 +
177)). The failure rate of 4.04\% reflects accidental password mistyping
--- entirely normal behavior that the anomaly detection model must learn
to tolerate. This failure rate is consistent with the empirically
observed rates of 3--5\% reported by Florencio and Herley {[}57{]} in
their large-scale study of user authentication behavior.

\textbf{User distribution.} The 64 user accounts exhibited varying SSH
usage frequencies. Developer and system administrator accounts had
significantly higher login frequencies (5--20 sessions per day) than
ordinary users (1--3 sessions per day). The 12 service accounts produced
highly regular, predictable patterns.

{[}Table 4.7: Simulation data detailed statistics{]}

{\def\LTcaptype{none} % do not increment counter
\begin{longtable}[]{@{}ll@{}}
\toprule\noalign{}
Statistic & Value \\
\midrule\noalign{}
\endhead
\bottomrule\noalign{}
\endlastfoot
Total log lines & 54,521 \\
SSHD entries after parsing & 34,700 \\
Collection period & \textasciitilde25 hours (2026-03-26) \\
User accounts & 64 (52 human + 12 service) \\
Network subnet & 192.168.152.0/24 \\
Accepted logins & 4,205 \\
Failed logins & 177 \\
Success rate & 95.96\% \\
Failure rate & 4.04\% \\
\end{longtable}
}

{[}Figure 4.2: Hourly activity distribution comparison between honeypot
(attack) data and simulation (normal) data, illustrating the diurnal
pattern of normal activity versus the 24/7 pattern of attack activity{]}

\subsection{4.2.7 Feature Distribution
Comparison}\label{feature-distribution-comparison}

Exploratory analysis of the 14 features across normal and attack classes
reveals clear distributional separation:

{[}Table 4.8: Descriptive statistics of key features by class{]}

{\def\LTcaptype{none} % do not increment counter
\begin{longtable}[]{@{}
  >{\raggedright\arraybackslash}p{(\linewidth - 6\tabcolsep) * \real{0.1364}}
  >{\raggedright\arraybackslash}p{(\linewidth - 6\tabcolsep) * \real{0.3333}}
  >{\raggedright\arraybackslash}p{(\linewidth - 6\tabcolsep) * \real{0.3333}}
  >{\raggedright\arraybackslash}p{(\linewidth - 6\tabcolsep) * \real{0.1970}}@{}}
\toprule\noalign{}
\begin{minipage}[b]{\linewidth}\raggedright
Feature
\end{minipage} & \begin{minipage}[b]{\linewidth}\raggedright
Normal (mean +/- std)
\end{minipage} & \begin{minipage}[b]{\linewidth}\raggedright
Attack (mean +/- std)
\end{minipage} & \begin{minipage}[b]{\linewidth}\raggedright
Separability
\end{minipage} \\
\midrule\noalign{}
\endhead
\bottomrule\noalign{}
\endlastfoot
fail\_count & 0.1 +/- 0.4 & 47.3 +/- 128.6 & Very High \\
success\_count & 2.8 +/- 1.9 & 0.02 +/- 0.15 & Very High \\
fail\_rate & 0.04 +/- 0.12 & 0.98 +/- 0.07 & Very High \\
unique\_usernames & 1.2 +/- 0.5 & 8.7 +/- 12.3 & High \\
invalid\_user\_count & 0.0 +/- 0.02 & 18.9 +/- 52.4 & Very High \\
invalid\_user\_ratio & 0.0 +/- 0.01 & 0.41 +/- 0.35 & High \\
connection\_count & 3.1 +/- 2.4 & 53.6 +/- 141.2 & Very High \\
mean\_inter\_attempt\_time & 87.4 +/- 62.1 & 0.8 +/- 2.3 & Very High \\
std\_inter\_attempt\_time & 43.2 +/- 35.7 & 0.5 +/- 1.8 & Very High \\
min\_inter\_attempt\_time & 12.6 +/- 18.3 & 0.02 +/- 0.1 & Very High \\
unique\_ports & 2.8 +/- 2.1 & 42.1 +/- 112.7 & High \\
session\_duration\_mean & 847.3 +/- 1204.5 & 2.1 +/- 4.7 & Very High \\
pam\_failure\_escalation & 0.0 +/- 0.03 & 0.28 +/- 0.45 & Moderate \\
max\_retries\_exceeded & 0.0 +/- 0.01 & 0.19 +/- 0.39 & Moderate \\
\end{longtable}
}

The temporal features (mean\_inter\_attempt\_time,
min\_inter\_attempt\_time, session\_duration\_mean) and the ratio
features (fail\_rate, invalid\_user\_ratio) exhibit the greatest
distributional separation between classes, with nearly non-overlapping
distributions. The count features (fail\_count, connection\_count) also
show large separation but with heavier-tailed distributions in the
attack class. The binary features (pam\_failure\_escalation,
max\_retries\_exceeded) show moderate separation --- they are present in
a significant fraction of attack windows but absent in nearly all normal
windows.

{[}Figure 4.3: Boxplot comparison of the 14 feature distributions
between normal (blue) and attack (red) classes, with logarithmic y-axis
for count features{]}

\section{4.3 Model Training and Evaluation
(Baseline)}\label{model-training-and-evaluation-baseline}

\subsection{4.3.1 Training Configuration}\label{training-configuration}

All three models were trained on the training set of 7,212 exclusively
normal samples, preprocessed with RobustScaler as described in Section
3.5.6. The baseline hyperparameters were set as specified in Section
3.7:

\begin{itemize}
\tightlist
\item
  \textbf{Isolation Forest:} n\_estimators=300, max\_samples=512,
  max\_features=0.5, contamination=auto
\item
  \textbf{LOF:} n\_neighbors=30, novelty=True, metric=Minkowski
  (Euclidean)
\item
  \textbf{OCSVM:} kernel=RBF, gamma=auto (1/14 \(\approx\) 0.0714),
  nu=0.01
\end{itemize}

\subsection{4.3.2 Baseline Results}\label{baseline-results}

{[}Table 4.9: Performance comparison of three models on the test set
(baseline, 15,184 samples){]}

{\def\LTcaptype{none} % do not increment counter
\begin{longtable}[]{@{}llll@{}}
\toprule\noalign{}
Metric & Isolation Forest & LOF & OCSVM \\
\midrule\noalign{}
\endhead
\bottomrule\noalign{}
\endlastfoot
Accuracy & 0.8076 & 0.8415 & \textbf{0.8573} \\
Precision & 0.7959 & 0.8256 & \textbf{0.8401} \\
Recall & 0.9999 & \textbf{1.0000} & \textbf{1.0000} \\
F1-Score & 0.8863 & 0.9045 & \textbf{0.9131} \\
ROC-AUC & 0.8316 & \textbf{0.9759} & 0.9003 \\
FPR & 0.7692 & 0.6338 & \textbf{0.5709} \\
Training Time & 0.4005s & 0.083s & \textbf{0.0132s} \\
\end{longtable}
}

All three models achieve near-perfect recall (\(\geq 0.9999\)),
confirming the fundamental viability of the semi-supervised anomaly
detection approach for SSH brute-force detection: the models trained
exclusively on normal data can identify virtually all attack samples in
the test set. This result aligns with the theoretical expectations
established by Goldstein and Uchida {[}48{]}, who showed that novelty
detection methods are highly effective when the normal class is
well-defined and the anomalies are sufficiently different.

\subsection{4.3.3 Confusion Matrices}\label{confusion-matrices}

{[}Table 4.10: Confusion matrices for three models (baseline){]}

{\def\LTcaptype{none} % do not increment counter
\begin{longtable}[]{@{}
  >{\raggedright\arraybackslash}p{(\linewidth - 8\tabcolsep) * \real{0.2188}}
  >{\raggedright\arraybackslash}p{(\linewidth - 8\tabcolsep) * \real{0.2188}}
  >{\raggedright\arraybackslash}p{(\linewidth - 8\tabcolsep) * \real{0.1875}}
  >{\raggedright\arraybackslash}p{(\linewidth - 8\tabcolsep) * \real{0.1875}}
  >{\raggedright\arraybackslash}p{(\linewidth - 8\tabcolsep) * \real{0.1875}}@{}}
\toprule\noalign{}
\begin{minipage}[b]{\linewidth}\raggedright
Model
\end{minipage} & \begin{minipage}[b]{\linewidth}\raggedright
True Negative (TN)
\end{minipage} & \begin{minipage}[b]{\linewidth}\raggedright
False Positive (FP)
\end{minipage} & \begin{minipage}[b]{\linewidth}\raggedright
False Negative (FN)
\end{minipage} & \begin{minipage}[b]{\linewidth}\raggedright
True Positive (TP)
\end{minipage} \\
\midrule\noalign{}
\endhead
\bottomrule\noalign{}
\endlastfoot
Isolation Forest & 876 & 2,920 & 1 & 11,387 \\
LOF & 1,390 & 2,406 & 0 & 11,388 \\
OCSVM & 1,629 & 2,167 & 0 & 11,388 \\
\end{longtable}
}

The confusion matrices reveal the critical performance difference among
the three models: while all achieve near-perfect TP counts, they differ
substantially in the TN/FP trade-off:

\textbf{Isolation Forest} correctly identifies only 876 of 3,796 normal
samples as normal (TN=876), misclassifying 2,920 normal samples as
attacks (FP=2,920), yielding an FPR of 76.92\%. The single false
negative (FN=1) is an attack sample whose feature values happened to
fall within the model's learned normal boundary --- likely a very
low-activity attack window with only 1--2 events.

\textbf{LOF} improves the normal identification to TN=1,390 (FP=2,406,
FPR=63.38\%), with zero false negatives. The improvement over IF in
false positive rate is attributable to LOF's local density-based
approach, which creates a tighter boundary around the normal data
clusters.

\textbf{OCSVM} achieves the best normal identification at TN=1,629
(FP=2,167, FPR=57.09\%), with zero false negatives. The RBF kernel
enables OCSVM to capture the nonlinear boundary of the normal region
more precisely than the axis-aligned partitions of Isolation Forest.

{[}Figure 4.4: Normalized confusion matrices for three models
(baseline), visualized as heatmaps{]}

\subsection{4.3.4 Baseline Analysis}\label{baseline-analysis}

The baseline results establish several critical findings:

\textbf{Finding 1: High recall is achievable.} All three models achieve
recall \(\geq 0.9999\), meaning that virtually no attacks are missed. In
the context of SSH security, where a single missed attack can lead to
system compromise, this is the paramount requirement. The
semi-supervised approach --- training only on normal data --- is
sufficient to achieve this recall level.

\textbf{Finding 2: False positive rate is the primary weakness.} The FPR
ranges from 57.09\% (OCSVM) to 76.92\% (IF). This means that 57--77\% of
normal samples are incorrectly flagged as attacks. While this is
acceptable as a baseline (given the asymmetric cost of FPs vs.~FNs in
cybersecurity), it motivates the optimization process described in
Section 4.5 and the dynamic threshold mechanism described in Section
4.6.

\textbf{Finding 3: OCSVM achieves the best baseline performance.} OCSVM
leads in accuracy (0.8573), precision (0.8401), and F1-score (0.9131).
The nonlinear RBF kernel decision boundary provides a tighter fit to the
normal data distribution than IF's axis-aligned random partitions or
LOF's density-based approach.

\textbf{Finding 4: LOF achieves the best ROC-AUC.} LOF's ROC-AUC of
0.9759 significantly exceeds both IF (0.8316) and OCSVM (0.9003),
indicating that LOF produces the best overall separation between normal
and attack score distributions across all threshold values, even though
its performance at the specific operating point (fixed threshold) is not
the best.

\textbf{Finding 5: Training times are negligible.} All three models
train in under 0.5 seconds on the 7,212-sample training set. OCSVM is
the fastest (0.0132s), benefiting from the relatively small training set
size. IF takes the longest (0.4005s) due to the construction of 300
isolation trees, but this is still well within practical requirements.

{[}Figure 4.5: ROC curves for three models (baseline), with AUC values
annotated{]}

\section{4.4 Feature Importance
Experiment}\label{feature-importance-experiment}

\subsection{4.4.1 Methodology}\label{methodology}

Feature importance was evaluated using the \textbf{permutation
importance} method {[}105{]} applied to the Isolation Forest model on
the test set. Permutation importance measures the decrease in model
performance (F1-score) when a single feature's values are randomly
shuffled, destroying the information content of that feature while
preserving the marginal distributions of all other features. The
procedure was repeated 10 times for each feature to obtain stable
importance estimates with confidence intervals.

The permutation importance of feature \(j\) is defined as:

\[\text{PI}_j = \frac{1}{K} \sum_{k=1}^{K} \left[ F_1(\mathbf{X}_{\text{test}}) - F_1(\mathbf{X}_{\text{test}}^{(j,k)}) \right]\]

where \(\mathbf{X}_{\text{test}}^{(j,k)}\) is the test set with the
\(j\)-th feature randomly permuted in the \(k\)-th repetition, and
\(K = 10\) is the number of repetitions.

\subsection{4.4.2 Results}\label{results}

{[}Table 4.11: Feature importance ranking (all 14 features){]}

{\def\LTcaptype{none} % do not increment counter
\begin{longtable}[]{@{}lllll@{}}
\toprule\noalign{}
Rank & Feature & Importance (\%) & Std (\%) & Category \\
\midrule\noalign{}
\endhead
\bottomrule\noalign{}
\endlastfoot
1 & session\_duration\_mean & 5.50 & 0.32 & Session \\
2 & min\_inter\_attempt\_time & 3.86 & 0.28 & Temporal \\
3 & mean\_inter\_attempt\_time & 2.61 & 0.24 & Temporal \\
4 & std\_inter\_attempt\_time & 1.64 & 0.19 & Temporal \\
5 & unique\_ports & 1.42 & 0.17 & Network \\
6 & connection\_count & 1.16 & 0.15 & Temporal \\
7 & success\_count & 0.44 & 0.09 & Volume \\
8 & fail\_rate & \textless0.10 & -- & Volume \\
9 & fail\_count & \textless0.10 & -- & Volume \\
10 & unique\_usernames & \textless0.10 & -- & Username \\
11 & invalid\_user\_count & \textless0.10 & -- & Username \\
12 & invalid\_user\_ratio & \textless0.10 & -- & Username \\
13 & pam\_failure\_escalation & \textless0.10 & -- & Indicator \\
14 & max\_retries\_exceeded & \textless0.10 & -- & Indicator \\
\end{longtable}
}

{[}Figure 4.6: Bar chart of permutation importance for all 14 features,
with error bars showing standard deviation across 10 repetitions{]}

\subsection{4.4.3 Interpretation}\label{interpretation}

The feature importance results reveal a clear and theoretically
meaningful hierarchy:

\textbf{The dominance of temporal and session features.} The top 6
features are all temporal, session, or network features. Collectively,
the top 5 features account for 15.03\% of total importance, while the
remaining 9 features collectively account for less than 1\%. This
dramatic skew demonstrates that the model's discrimination between
normal and attack traffic is overwhelmingly driven by the timing and
session characteristics of SSH connections, rather than by count-based
or username-based features.

\textbf{session\_duration\_mean (5.50\%)} is the single most important
feature. This reflects the most fundamental behavioral difference
between attacks and normal usage: brute-force attacks produce extremely
short SSH sessions (typically under 5 seconds, because each failed
authentication results in rapid disconnection by the server), while
legitimate users maintain sessions lasting minutes to hours. The bimodal
gap between these distributions (attack sessions: 0.5--5 seconds; normal
sessions: 300--7,200+ seconds) provides an almost noise-free
classification signal. This finding is consistent with Javed and Paxson
{[}72{]}, who identified session duration as one of the most reliable
indicators of automated SSH activity.

\textbf{min\_inter\_attempt\_time (3.86\%)} is the second most important
feature. Automated attack tools dispatch authentication attempts at
machine speed, producing inter-attempt intervals of milliseconds to low
seconds. Even in a low-and-slow attack where the mean interval is
deliberately increased, the minimum interval often reveals the
attacker's true timing signature --- for example, when the tool
dispatches a rapid burst before entering a deliberate delay phase.

\textbf{mean\_inter\_attempt\_time (2.61\%)} and
\textbf{std\_inter\_attempt\_time (1.64\%)} together capture the overall
speed and regularity of authentication behavior. Automated tools produce
low mean and low standard deviation (fast and uniform), while humans
produce high mean and high standard deviation (slow and irregular).

\textbf{unique\_ports (1.42\%)} serves as a proxy for the total number
of distinct TCP connections. Each new SSH connection uses a different
ephemeral source port, so a large number of unique ports directly
indicates a large number of separate connection attempts --- a hallmark
of brute-force behavior.

\textbf{The low importance of count-based features.} Surprisingly,
fail\_count, fail\_rate, unique\_usernames, and invalid\_user\_count ---
the features that correspond most directly to the traditional rule-based
detection approach (e.g., Fail2Ban's maxretry threshold) --- have
negligible permutation importance (\textless0.10\%). This does not mean
these features are uninformative in an absolute sense, but rather that
the Isolation Forest model's decision boundary is primarily shaped by
temporal and session features, and the count-based features provide
redundant information once the temporal features are available. This
finding has important practical implications: it suggests that
traditional count-based detection methods (which rely exclusively on
features like fail\_count) are using the least discriminative dimensions
of the feature space, which explains their documented vulnerability to
rate-limiting evasion techniques {[}13{]}.

\section{4.5 Model Optimization
Results}\label{model-optimization-results}

\subsection{4.5.1 Optimization Strategy}\label{optimization-strategy}

The optimization process targeted three sources of performance
degradation identified through analysis of the baseline results:

\textbf{Source 1: Temporal autocorrelation from overlapping windows.}
The baseline configuration used overlapping sliding windows (window=5
min, stride=1 min), causing consecutive feature vectors to share 80\% of
their underlying events. This temporal correlation inflates the apparent
training set size without proportionally increasing its information
content, and causes the model to learn correlated patterns that do not
generalize well to non-overlapping test conditions. Switching to
non-overlapping windows (stride=5 min) eliminates this correlation,
producing independent feature vectors that each represent a distinct
5-minute behavioral snapshot. Cerqueira et al.~{[}106{]} demonstrated
that temporal correlation between training samples is a major source of
overfitting in time-series anomaly detection.

\textbf{Source 2: Limited feature expressiveness.} The baseline used 14
raw features. The optimized configuration adds 9 derived features that
capture second-order behavioral patterns, including: the ratio of failed
attempts to unique usernames (indicating whether the attacker uses many
passwords per username or many usernames per password), the coefficient
of variation of inter-attempt times (capturing the regularity of timing
patterns independent of absolute speed), and interaction terms between
temporal and count features. These derived features increase the
effective dimensionality from 14 to 23, providing additional
discriminative power for edge cases where the raw features alone are
ambiguous {[}107{]}.

\textbf{Source 3: Suboptimal hyperparameters.} The baseline
hyperparameters were set conservatively. Grid search over the expanded
parameter space identified improved values for contamination,
max\_features, and n\_estimators.

\subsection{4.5.2 Optimized Results}\label{optimized-results}

{[}Table 4.12: Optimized model performance comparison{]}

{\def\LTcaptype{none} % do not increment counter
\begin{longtable}[]{@{}llll@{}}
\toprule\noalign{}
Metric & Isolation Forest & LOF & OCSVM \\
\midrule\noalign{}
\endhead
\bottomrule\noalign{}
\endlastfoot
Accuracy & 0.9031 & 0.8322 & \textbf{0.9138} \\
F1-Score & 0.9374 & 0.8994 & \textbf{0.9455} \\
Recall & 0.9675 & \textbf{1.0000} & 0.9965 \\
FPR & 0.2900 & 0.6710 & \textbf{0.3342} \\
Training Time & 0.65s & 0.09s & 0.02s \\
\end{longtable}
}

\textbf{Optimized Isolation Forest hyperparameters:} contamination=0.01,
max\_features=0.75, max\_samples=512, n\_estimators=500.

{[}Table 4.13: Confusion matrices for three models (optimized){]}

{\def\LTcaptype{none} % do not increment counter
\begin{longtable}[]{@{}lllll@{}}
\toprule\noalign{}
Model & TN & FP & FN & TP \\
\midrule\noalign{}
\endhead
\bottomrule\noalign{}
\endlastfoot
Isolation Forest & 546 & 223 & 75 & 2,232 \\
LOF & 253 & 516 & 0 & 2,307 \\
OCSVM & 512 & 257 & 8 & 2,299 \\
\end{longtable}
}

Note: The optimized evaluation uses non-overlapping windows, resulting
in a smaller test set (3,076 samples vs.~15,184 in the baseline) with a
different class distribution. The total samples per model in the
confusion matrix reflect this reduced test set size.

\subsection{4.5.3 Improvement Analysis}\label{improvement-analysis}

{[}Table 4.14: Isolation Forest baseline vs.~optimized comparison{]}

{\def\LTcaptype{none} % do not increment counter
\begin{longtable}[]{@{}
  >{\raggedright\arraybackslash}p{(\linewidth - 8\tabcolsep) * \real{0.1538}}
  >{\raggedright\arraybackslash}p{(\linewidth - 8\tabcolsep) * \real{0.1923}}
  >{\raggedright\arraybackslash}p{(\linewidth - 8\tabcolsep) * \real{0.2115}}
  >{\raggedright\arraybackslash}p{(\linewidth - 8\tabcolsep) * \real{0.1538}}
  >{\raggedright\arraybackslash}p{(\linewidth - 8\tabcolsep) * \real{0.2885}}@{}}
\toprule\noalign{}
\begin{minipage}[b]{\linewidth}\raggedright
Metric
\end{minipage} & \begin{minipage}[b]{\linewidth}\raggedright
Baseline
\end{minipage} & \begin{minipage}[b]{\linewidth}\raggedright
Optimized
\end{minipage} & \begin{minipage}[b]{\linewidth}\raggedright
Change
\end{minipage} & \begin{minipage}[b]{\linewidth}\raggedright
Interpretation
\end{minipage} \\
\midrule\noalign{}
\endhead
\bottomrule\noalign{}
\endlastfoot
Accuracy & 0.8076 & 0.9031 & +0.0955 (+9.55 pp) & Substantial
improvement \\
F1-Score & 0.8863 & 0.9374 & +0.0511 (+5.11 pp) & Significant
improvement \\
Recall & 0.9999 & 0.9675 & -0.0324 (-3.24 pp) & Minor trade-off \\
FPR & 0.7692 & 0.2900 & -0.4792 (-47.92 pp) & Dramatic improvement \\
Precision & 0.7959 & \textasciitilde0.91 & +0.11 (+11 pp) & Major
improvement \\
\end{longtable}
}

The optimization process produced dramatic improvements across all
primary metrics except recall, which decreased by 3.24 percentage points
(from 0.9999 to 0.9675). This slight recall reduction represents a
conscious and desirable trade-off: the baseline model achieved
near-perfect recall at the cost of an unacceptably high false positive
rate (76.92\%), which would make the system unusable in practice due to
excessive false alarms. The optimized model sacrifices a small amount of
recall (75 missed attacks out of 2,307 total) in exchange for a
48-percentage-point reduction in FPR (from 76.92\% to 29.00\%),
representing a far more practical operating point for production
deployment.

\textbf{Attribution of improvements.} The three optimization changes
contributed to the improvement as follows:

\begin{enumerate}
\def\labelenumi{\arabic{enumi}.}
\item
  \textbf{Non-overlapping windows} (primary contributor): Estimated to
  account for approximately 20 percentage points of FPR reduction. By
  eliminating the 80\% temporal overlap between consecutive feature
  vectors, the model learns from genuinely independent observations,
  producing a more accurate representation of the normal behavioral
  boundary.
\item
  \textbf{Derived features} (secondary contributor): Estimated to
  account for approximately 15 percentage points of FPR reduction. The
  second-order features provide critical discriminative information for
  edge cases --- particularly normal windows with slightly elevated
  fail\_count or slightly reduced mean\_inter\_attempt\_time that the
  raw features alone cannot distinguish from mild attack activity.
\item
  \textbf{Hyperparameter tuning} (tertiary contributor): Estimated to
  account for approximately 13 percentage points of FPR reduction.
  Setting contamination=0.01 explicitly constrains the model to treat at
  most 1\% of training samples as outliers, creating a tighter normal
  boundary. Increasing max\_features from 0.5 to 0.75 allows each tree
  to consider more features simultaneously, improving detection of
  attacks that manifest in specific feature subsets. Increasing
  n\_estimators from 300 to 500 provides more stable score estimates
  through greater ensemble averaging.
\end{enumerate}

\subsection{4.5.4 Model Selection
Justification}\label{model-selection-justification}

Although OCSVM achieves the highest F1-Score (0.9455) and accuracy
(0.9138) in the optimized configuration, \textbf{Isolation Forest was
selected as the primary model} for the deployed system. This decision
was based on three operational considerations:

\textbf{First, score distribution suitability.} IF produces continuous
anomaly scores in the range {[}0, 1{]} based on the average path length
in random decision trees. The EWMA-Adaptive Percentile dynamic threshold
(Section 3.8) requires a continuous, smoothly varying score sequence to
compute meaningful EWMA and percentile values. IF's score distribution
is smoother and less sensitive to extreme outliers than LOF's density
ratios or OCSVM's signed distances.

\textbf{Second, computational efficiency.} IF has a training complexity
of \(O(T \psi \log \psi)\), where \(T = 500\) and \(\psi = 512\), making
retraining feasible in under 1 second even on modest hardware. OCSVM's
\(O(n^2)\) to \(O(n^3)\) training complexity becomes prohibitive as the
training set grows, and LOF's \(O(n^2 \log n)\) complexity similarly
limits scalability.

\textbf{Third, interpretability.} IF's path-length-based anomaly score
has a natural interpretation: shorter path lengths indicate points that
are easier to isolate, which corresponds intuitively to behavioral
anomalies. This interpretability facilitates communication with security
operators who need to understand why a specific IP was flagged.

\section{4.6 Dynamic Threshold Engine
Evaluation}\label{dynamic-threshold-engine-evaluation}

\subsection{4.6.1 Experimental Design}\label{experimental-design}

The EWMA-Adaptive Percentile dynamic threshold engine was evaluated with
the parameters specified in Section 3.8.2: \(\alpha = 0.3\), \(p = 95\),
\(\lambda = 1.5\), \(L = 100\). The evaluation assessed three aspects:
(1) adaptation to distribution shifts, (2) reduction of burst false
positives, and (3) early detection capability.

\subsection{4.6.2 Distribution Shift
Adaptation}\label{distribution-shift-adaptation}

When the anomaly score distribution changes due to shifts in traffic
patterns (e.g., transition from business hours to off-hours, or onset of
a scanning campaign), the dynamic threshold adjusts automatically.
During normal periods, the EWMA tracks the baseline score level, and the
threshold remains close to the 95th percentile of recent scores,
providing sensitive detection. When a sustained increase in mildly
anomalous activity occurs (e.g., many users simultaneously mistyping
passwords after a policy change), the EWMA rises, pulling the threshold
upward and preventing a cascade of false alarms.

{[}Figure 4.7: Anomaly score and dynamic threshold evolution over time
on test data, showing automatic threshold adjustment during distribution
shifts{]}

\subsection{4.6.3 Burst False Positive
Reduction}\label{burst-false-positive-reduction}

With a static threshold (set at the optimal operating point for the test
set), bursts of mildly anomalous legitimate activity generate clusters
of false positives. The dynamic threshold detects the general increase
in baseline scores through the EWMA component and elevates the threshold
proportionally, reducing burst false positive rates by an estimated
10--15\% compared to the static threshold.

{[}Table 4.15: Dynamic threshold vs.~static threshold comparison
(Isolation Forest, optimized){]}

{\def\LTcaptype{none} % do not increment counter
\begin{longtable}[]{@{}lll@{}}
\toprule\noalign{}
Metric & Static Threshold (Optimal) & Dynamic Threshold (EWMA) \\
\midrule\noalign{}
\endhead
\bottomrule\noalign{}
\endlastfoot
Precision & \textasciitilde0.91 & \textasciitilde0.93 \\
Recall & 0.9675 & \textasciitilde0.96 \\
F1-Score & 0.9374 & \textasciitilde0.94 \\
FPR & 0.2900 & \textasciitilde0.25 \\
Adaptation to shifts & None & Automatic \\
Early warning capability & None & Yes (EARLY\_WARNING level) \\
\end{longtable}
}

The dynamic threshold does not necessarily outperform the static
threshold on a static, i.i.d. test set (since the static threshold can
be optimized retrospectively for that specific set). However, the
dynamic threshold's advantage is operational: it performs well on
non-stationary streaming data where the optimal threshold cannot be
known in advance, and it provides the EARLY\_WARNING capability that is
absent from static approaches.

\subsection{4.6.4 Two-Level Detection
Behavior}\label{two-level-detection-behavior}

The two-level detection mechanism (EARLY\_WARNING at \(\theta_t / 1.5\)
and ALERT at \(\theta_t\)) provides graduated response:

\begin{itemize}
\item
  \textbf{EARLY\_WARNING} triggers when the anomaly score exceeds
  approximately 67\% of the current dynamic threshold. At this level,
  the system logs the event and notifies the security administrator but
  does not block the IP. This is appropriate for borderline cases (e.g.,
  a legitimate user with slightly anomalous behavior, or the early
  stages of a slow attack) where premature blocking could disrupt
  service.
\item
  \textbf{ALERT} triggers when the anomaly score exceeds the full
  dynamic threshold. At this level, the system automatically invokes
  Fail2Ban to ban the IP address, in addition to logging and
  notification.
\end{itemize}

The two-level design is particularly valuable for the low-and-slow
attack scenario (discussed in detail in Section 4.8), where individual
anomaly scores are only slightly elevated and may fluctuate around the
EARLY\_WARNING threshold for several windows before reaching the ALERT
level.

\subsection{4.6.5 Parameter Sensitivity
Analysis}\label{parameter-sensitivity-analysis}

The sensitivity of detection performance to each threshold parameter was
investigated by varying one parameter at a time while holding the others
at their default values:

\textbf{Alpha (\(\alpha\)):} At \(\alpha = 0.1\), the EWMA responds
slowly to changes, suitable for very stable environments but slow to
detect sudden onset attacks. At \(\alpha = 0.3\) (selected), a balanced
trade-off between responsiveness and noise filtering is achieved. At
\(\alpha = 0.5\), the EWMA responds rapidly but is more sensitive to
transient score fluctuations, causing the threshold to oscillate.

\textbf{Base percentile (\(p\)):} At \(p = 90\), the threshold is lower,
detecting more anomalies but generating approximately 20\% more false
positives. At \(p = 95\) (selected), an effective balance between
detection sensitivity and false alarm rate is achieved. At \(p = 99\),
the threshold is high, producing very few false positives but
potentially missing subtle attacks.

\textbf{Sensitivity factor (\(\lambda\)):} At \(\lambda = 1.0\), the
threshold equals the percentile, providing maximum sensitivity. At
\(\lambda = 1.5\) (selected), the threshold extends 50\% above the
EWMA-to-percentile gap, effectively reducing false positives. At
\(\lambda = 2.0\), the threshold is very conservative, suitable only for
environments with extremely low tolerance for alert volume.

{[}Figure 4.8: Effect of \(\alpha\), \(p\), and \(\lambda\) on F1-Score
and FPR, showing the optimal region around the selected parameter
values{]}

\subsection{4.6.6 Self-Calibration
Results}\label{self-calibration-results}

The self-calibration mechanism (recalibration every 100 decisions) was
evaluated by simulating a deployment scenario with drifting traffic
patterns. Over a simulated 24-hour period with varying attack intensity,
the self-calibrating threshold maintained a consistent detection rate
(F1-Score within \(\pm 0.02\) of the static-optimal value) despite the
distribution drift, while the static threshold showed F1-Score
degradation of up to 0.08 during high-drift periods.

\section{4.7 System Integration Test
Results}\label{system-integration-test-results}

\subsection{4.7.1 End-to-End Pipeline
Validation}\label{end-to-end-pipeline-validation}

The complete system pipeline was validated by injecting known attack and
normal SSH traffic into the target SSH server and verifying the
end-to-end flow from log event to detection decision and response
action.

{[}Table 4.16: Latency analysis by pipeline stage{]}

{\def\LTcaptype{none} % do not increment counter
\begin{longtable}[]{@{}
  >{\raggedright\arraybackslash}p{(\linewidth - 4\tabcolsep) * \real{0.2333}}
  >{\raggedright\arraybackslash}p{(\linewidth - 4\tabcolsep) * \real{0.5333}}
  >{\raggedright\arraybackslash}p{(\linewidth - 4\tabcolsep) * \real{0.2333}}@{}}
\toprule\noalign{}
\begin{minipage}[b]{\linewidth}\raggedright
Stage
\end{minipage} & \begin{minipage}[b]{\linewidth}\raggedright
Average Latency
\end{minipage} & \begin{minipage}[b]{\linewidth}\raggedright
Notes
\end{minipage} \\
\midrule\noalign{}
\endhead
\bottomrule\noalign{}
\endlastfoot
Log ingestion (Logstash) & \textless{} 1 second & Batch-dependent \\
Feature extraction (API Server) & \textless{} 100 ms & 14 features from
Redis buffer \\
Model inference (IF) & \textless{} 10 ms & Fastest (path length
computation) \\
Model inference (LOF) & \textless{} 50 ms & k-NN distance computation \\
Model inference (OCSVM) & \textless{} 30 ms & Kernel evaluation against
support vectors \\
Dynamic threshold update & \textless{} 5 ms & EWMA arithmetic \\
Decision + Fail2Ban API call & \textless{} 50 ms & HTTP API
invocation \\
Dashboard notification (WebSocket) & \textless{} 20 ms & Push
notification \\
\textbf{Total end-to-end} & \textbf{\textless{} 1--2 seconds} &
\textbf{Meets real-time requirements} \\
\end{longtable}
}

The total end-to-end latency from the moment an SSH authentication event
occurs to the moment a detection decision is rendered and an IP ban is
executed is under 2 seconds. The vast majority of this latency resides
in the log ingestion stage (Logstash file tailing and parsing); all
AI-related processing stages (feature extraction, model inference,
threshold computation) complete in under 150 ms combined.

\subsection{4.7.2 Throughput Analysis}\label{throughput-analysis}

The system was stress-tested to determine its maximum throughput
capacity. The API Server, using FastAPI's ASGI-based asynchronous
architecture, can process several hundred feature vectors per second on
a single worker, which is sufficient for monitoring dozens of SSH
servers simultaneously. At peak throughput, the bottleneck shifts from
the AI processing pipeline to the Elasticsearch write throughput,
confirming that the AI components do not limit system performance.

\subsection{4.7.3 Resource Utilization}\label{resource-utilization}

{[}Table 4.17: Resource utilization of Docker services under typical
load{]}

{\def\LTcaptype{none} % do not increment counter
\begin{longtable}[]{@{}llll@{}}
\toprule\noalign{}
Service & RAM (Approx.) & CPU (Avg.) & Notes \\
\midrule\noalign{}
\endhead
\bottomrule\noalign{}
\endlastfoot
API Server (FastAPI) & \textasciitilde200 MB & 5--15\% & Scales with
request volume \\
Detector (scikit-learn) & \textasciitilde150 MB & 2--5\% & Loaded model
+ scaler \\
Elasticsearch & \textasciitilde1--2 GB & 10--20\% & Index management,
continuous \\
Logstash & \textasciitilde500 MB & 5--10\% & JVM-based, log rate
dependent \\
Kibana & \textasciitilde400 MB & 3--8\% & JVM-based, increases with
dashboard use \\
React Frontend & \textasciitilde50 MB (client) & \textless1\% & Runs in
browser \\
Redis & \textasciitilde50 MB & \textless1\% & Lightweight in-memory
store \\
Fail2Ban & \textasciitilde30 MB & \textless1\% & Event-driven, minimal
footprint \\
SSH Target & \textasciitilde20 MB & \textless1\% & Standard OpenSSH \\
\end{longtable}
}

The total system resource footprint is approximately 2.5--4.5 GB of RAM,
well within the capacity of a modest server or cloud instance. The
Elasticsearch service is the most resource-intensive component,
consistent with its role as the primary data store and search engine.

\section{4.8 Attack Scenario Testing}\label{attack-scenario-testing}

\subsection{4.8.1 Scenario Design}\label{scenario-design}

Five attack scenarios were designed and executed against the test
environment, each simulating a distinct SSH brute-force attack variant
documented in the literature {[}11, 14, 15{]}. The attack scripts were
implemented in Python using the Paramiko SSH library, which provides
programmatic control over SSH connection parameters including timing,
username selection, and password selection.

{[}Table 4.18: Description of 5 attack scenarios{]}

{\def\LTcaptype{none} % do not increment counter
\begin{longtable}[]{@{}
  >{\raggedright\arraybackslash}p{(\linewidth - 10\tabcolsep) * \real{0.1020}}
  >{\raggedright\arraybackslash}p{(\linewidth - 10\tabcolsep) * \real{0.2041}}
  >{\raggedright\arraybackslash}p{(\linewidth - 10\tabcolsep) * \real{0.1020}}
  >{\raggedright\arraybackslash}p{(\linewidth - 10\tabcolsep) * \real{0.1429}}
  >{\raggedright\arraybackslash}p{(\linewidth - 10\tabcolsep) * \real{0.2245}}
  >{\raggedright\arraybackslash}p{(\linewidth - 10\tabcolsep) * \real{0.2245}}@{}}
\toprule\noalign{}
\begin{minipage}[b]{\linewidth}\raggedright
No.
\end{minipage} & \begin{minipage}[b]{\linewidth}\raggedright
Scenario
\end{minipage} & \begin{minipage}[b]{\linewidth}\raggedright
IPs
\end{minipage} & \begin{minipage}[b]{\linewidth}\raggedright
Speed
\end{minipage} & \begin{minipage}[b]{\linewidth}\raggedright
Usernames
\end{minipage} & \begin{minipage}[b]{\linewidth}\raggedright
Difficulty
\end{minipage} \\
\midrule\noalign{}
\endhead
\bottomrule\noalign{}
\endlastfoot
1 & Basic Brute-Force & 1 & Maximum & root & Easy \\
2 & Distributed Attack & 10+ & Moderate & Various & Medium \\
3 & Low-and-Slow & 1 & 30--120s delay & Various & Hard \\
4 & Credential Stuffing & 1--3 & Moderate-High & Leaked list & Medium \\
5 & Dictionary Attack & 1 & High & root & Easy \\
\end{longtable}
}

\subsection{4.8.2 Scenario 1: Basic
Brute-Force}\label{scenario-1-basic-brute-force}

\textbf{Configuration.} A single IP address attempts to authenticate as
root at maximum speed, dispatching attempts as fast as the SSH server
and network permit (typically 5--20 attempts per second). The password
list is systematically generated.

\textbf{Results.} All three models detect 100\% of attack windows.
Anomaly scores are far above the dynamic threshold (typically 3--5x the
threshold value). Detection occurs within the first 1-minute stride ---
a maximum latency of 60 seconds from attack onset to detection. The
feature profile is unmistakable: fail\_count in the hundreds, fail\_rate
\(\approx 1.0\), mean\_inter\_attempt\_time \(< 0.5\) seconds,
session\_duration\_mean \(< 3\) seconds, unique\_usernames \(= 1\).

\textbf{Key features activated.} fail\_count, fail\_rate,
mean\_inter\_attempt\_time, min\_inter\_attempt\_time,
session\_duration\_mean, connection\_count.

\subsection{4.8.3 Scenario 2: Distributed
Attack}\label{scenario-2-distributed-attack}

\textbf{Configuration.} 10+ IP addresses (simulating a small botnet)
each attempt 2--5 password combinations before cycling to the next
target. The total attack volume is comparable to Scenario 1, but
distributed across many sources to evade per-IP rate limiting.

\textbf{Results.} The system detects more than 90\% of attacking IPs.
The IPs with the highest per-IP attempt counts (3--5 attempts) are
detected reliably. However, IPs performing only 1--2 attempts within a
single 5-minute window produce anomaly scores in the borderline zone
between EARLY\_WARNING and ALERT thresholds. The temporal features
(session\_duration\_mean, min\_inter\_attempt\_time) remain anomalous
even for low-volume IPs because automated tools still produce
characteristically short sessions and machine-speed timing patterns.

\textbf{Key features activated.} session\_duration\_mean,
min\_inter\_attempt\_time, unique\_ports.

{[}Table 4.19: Distributed attack detection rates by per-IP attempt
count{]}

{\def\LTcaptype{none} % do not increment counter
\begin{longtable}[]{@{}
  >{\raggedright\arraybackslash}p{(\linewidth - 4\tabcolsep) * \real{0.3857}}
  >{\raggedright\arraybackslash}p{(\linewidth - 4\tabcolsep) * \real{0.2857}}
  >{\raggedright\arraybackslash}p{(\linewidth - 4\tabcolsep) * \real{0.3286}}@{}}
\toprule\noalign{}
\begin{minipage}[b]{\linewidth}\raggedright
Attempts per IP per Window
\end{minipage} & \begin{minipage}[b]{\linewidth}\raggedright
Detection Rate (IF)
\end{minipage} & \begin{minipage}[b]{\linewidth}\raggedright
Detection Rate (OCSVM)
\end{minipage} \\
\midrule\noalign{}
\endhead
\bottomrule\noalign{}
\endlastfoot
1--2 & \textasciitilde75\% & \textasciitilde82\% \\
3--5 & \textasciitilde95\% & \textasciitilde97\% \\
6+ & 100\% & 100\% \\
\end{longtable}
}

\subsection{4.8.4 Scenario 3:
Low-and-Slow}\label{scenario-3-low-and-slow}

\textbf{Configuration.} A single IP address spaces authentication
attempts over intervals of 30--120 seconds, deliberately staying below
traditional rate-limiting thresholds (e.g., Fail2Ban's default
maxretry=5/findtime=600s). The attacker targets various usernames with a
curated password list.

\textbf{Results.} This is the most challenging scenario and represents
the boundary of the system's detection capabilities. Individual 5-minute
windows contain only 1--3 events per window, producing feature vectors
with modestly elevated anomaly scores that may fluctuate around the
EARLY\_WARNING threshold. The key detection mechanism is the
\textbf{EWMA accumulation}: although each individual window's anomaly
score may not exceed the ALERT threshold, the cumulative effect of 3--5
consecutive windows with slightly elevated scores causes the EWMA value
to rise, eventually crossing the EARLY\_WARNING threshold (at
approximately 2--3 minutes after attack onset) and potentially the ALERT
threshold (at approximately 5--8 minutes).

The detection relies primarily on two features that remain anomalous
even at low attempt rates: \textbf{session\_duration\_mean} (attack
sessions are still extremely short --- under 5 seconds --- because even
though the attacker waits between sessions, each individual session
consists of a single failed attempt followed by disconnection) and
\textbf{invalid\_user\_ratio} (if the attacker uses non-existent
usernames, this ratio is elevated regardless of attempt frequency).

\textbf{Comparison with Fail2Ban.} Under Fail2Ban's default
configuration (maxretry=5, findtime=600 seconds), a low-and-slow attack
spacing attempts at 2-minute intervals would produce only 3 failures
within the 10-minute findtime window, falling below the maxretry=5
threshold and evading detection entirely. The AI system's EWMA
accumulation mechanism detects the same attack pattern within 2--3
minutes through cumulative score elevation.

{[}Table 4.20: Low-and-Slow detection results by model{]}

{\def\LTcaptype{none} % do not increment counter
\begin{longtable}[]{@{}
  >{\raggedright\arraybackslash}p{(\linewidth - 6\tabcolsep) * \real{0.1000}}
  >{\raggedright\arraybackslash}p{(\linewidth - 6\tabcolsep) * \real{0.2143}}
  >{\raggedright\arraybackslash}p{(\linewidth - 6\tabcolsep) * \real{0.4000}}
  >{\raggedright\arraybackslash}p{(\linewidth - 6\tabcolsep) * \real{0.2857}}@{}}
\toprule\noalign{}
\begin{minipage}[b]{\linewidth}\raggedright
Model
\end{minipage} & \begin{minipage}[b]{\linewidth}\raggedright
Detection Rate
\end{minipage} & \begin{minipage}[b]{\linewidth}\raggedright
Avg Windows to EARLY\_WARNING
\end{minipage} & \begin{minipage}[b]{\linewidth}\raggedright
Avg Windows to ALERT
\end{minipage} \\
\midrule\noalign{}
\endhead
\bottomrule\noalign{}
\endlastfoot
IF & \textasciitilde70\% & 3--4 windows & 5--7 windows \\
LOF & \textasciitilde80\% & 3 windows & 4--6 windows \\
OCSVM & \textasciitilde85\% & 2--3 windows & 4--5 windows \\
\end{longtable}
}

\subsection{4.8.5 Scenario 4: Credential
Stuffing}\label{scenario-4-credential-stuffing}

\textbf{Configuration.} 1--3 IP addresses use a database of
username-password pairs (simulating leaked credentials) to test each
pair against the SSH server. The distinguishing characteristic is a
large number of unique usernames (many of which do not exist on the
target system) with only 1--2 attempts per username.

\textbf{Results.} Detection rate exceeds 95\% for all models. The
feature profile is distinctive: unique\_usernames is very high (tens of
different usernames per window), invalid\_user\_ratio is elevated (since
many usernames from leaked lists do not exist on the target system), and
the temporal features still indicate automated behavior (short sessions,
sub-second inter-attempt times).

\textbf{Key features activated.} unique\_usernames,
invalid\_user\_count, invalid\_user\_ratio, session\_duration\_mean.

\subsection{4.8.6 Scenario 5: Dictionary
Attack}\label{scenario-5-dictionary-attack}

\textbf{Configuration.} A single IP address targets root using a curated
dictionary of common passwords (e.g., RockYou, SecLists top-1000). The
attack proceeds at high speed with a single target username.

\textbf{Results.} Detection rate is 100\% for all three models. The
feature profile is similar to Scenario 1 (basic brute-force): very high
fail\_count, fail\_rate \(\approx 1.0\), unique\_usernames \(= 1\), very
low mean\_inter\_attempt\_time and session\_duration\_mean. The primary
difference from Scenario 1 is that the password list is curated rather
than systematic, but this difference does not affect the behavioral
features that the model uses for detection.

\subsection{4.8.7 Consolidated Scenario
Results}\label{consolidated-scenario-results}

{[}Table 4.21: Summary of detection results across 5 attack scenarios{]}

{\def\LTcaptype{none} % do not increment counter
\begin{longtable}[]{@{}
  >{\raggedright\arraybackslash}p{(\linewidth - 10\tabcolsep) * \real{0.1613}}
  >{\raggedright\arraybackslash}p{(\linewidth - 10\tabcolsep) * \real{0.1774}}
  >{\raggedright\arraybackslash}p{(\linewidth - 10\tabcolsep) * \real{0.0968}}
  >{\raggedright\arraybackslash}p{(\linewidth - 10\tabcolsep) * \real{0.0968}}
  >{\raggedright\arraybackslash}p{(\linewidth - 10\tabcolsep) * \real{0.1129}}
  >{\raggedright\arraybackslash}p{(\linewidth - 10\tabcolsep) * \real{0.3548}}@{}}
\toprule\noalign{}
\begin{minipage}[b]{\linewidth}\raggedright
Scenario
\end{minipage} & \begin{minipage}[b]{\linewidth}\raggedright
Difficulty
\end{minipage} & \begin{minipage}[b]{\linewidth}\raggedright
IF
\end{minipage} & \begin{minipage}[b]{\linewidth}\raggedright
LOF
\end{minipage} & \begin{minipage}[b]{\linewidth}\raggedright
OCSVM
\end{minipage} & \begin{minipage}[b]{\linewidth}\raggedright
Primary Discriminators
\end{minipage} \\
\midrule\noalign{}
\endhead
\bottomrule\noalign{}
\endlastfoot
1. Basic Brute-Force & Easy & 100\% & 100\% & 100\% & fail\_count,
fail\_rate, timing \\
2. Distributed & Medium & \textgreater90\% & \textgreater92\% &
\textgreater93\% & session\_duration, timing, ports \\
3. Low-and-Slow & Hard & \textasciitilde70\% & \textasciitilde80\% &
\textasciitilde85\% & session\_duration, EWMA accum. \\
4. Credential Stuffing & Medium & \textgreater95\% & \textgreater96\% &
\textgreater97\% & unique\_usernames, invalid\_user \\
5. Dictionary Attack & Easy & 100\% & 100\% & 100\% & fail\_count,
timing, session \\
\end{longtable}
}

{[}Figure 4.9: Heatmap of detection rates by attack scenario and
model{]}

The consolidated results yield several important findings:

\begin{enumerate}
\def\labelenumi{\arabic{enumi}.}
\item
  \textbf{High-speed attacks (Scenarios 1, 5) are trivially detected.}
  All models achieve 100\% detection, with anomaly scores far exceeding
  any reasonable threshold. These scenarios validate the basic
  functionality of the system but do not test its limits.
\item
  \textbf{Distributed attacks (Scenario 2) require multidimensional
  analysis.} Per-IP count features alone are insufficient when each IP
  makes only a few attempts. The temporal and session features (which
  remain anomalous even at low per-IP volumes because automated tools
  still produce machine-speed timing and short sessions) are essential
  for detection.
\item
  \textbf{Low-and-slow attacks (Scenario 3) are the critical challenge.}
  Detection rates of 70--85\% represent the system's performance
  frontier. The EWMA accumulation mechanism provides the key capability:
  detecting the persistent, cumulative pattern of slightly anomalous
  activity that no single-window analysis can reliably identify. This is
  the scenario where the dynamic threshold provides the greatest
  marginal value over static approaches.
\item
  \textbf{OCSVM consistently achieves the highest per-scenario detection
  rates.} The nonlinear RBF kernel decision boundary is more sensitive
  to subtle deviations from the normal profile than IF's axis-aligned
  random partitions. However, IF's computational efficiency and score
  distribution properties make it the preferred choice for the real-time
  system.
\end{enumerate}

\section{4.9 Detection Time Comparison: AI
vs.~Fail2Ban}\label{detection-time-comparison-ai-vs.-fail2ban}

\subsection{4.9.1 Comparative Framework}\label{comparative-framework}

To quantify the detection time advantage of the AI-based system over
traditional rate-limiting approaches, both systems were evaluated on the
same attack traffic under the same conditions. Fail2Ban was configured
with its default SSH parameters: maxretry=5, findtime=600 seconds,
bantime=600 seconds.

\subsection{4.9.2 Detection Time Results}\label{detection-time-results}

{[}Table 4.22: Detection time comparison: AI system vs.~Fail2Ban{]}

{\def\LTcaptype{none} % do not increment counter
\begin{longtable}[]{@{}
  >{\raggedright\arraybackslash}p{(\linewidth - 6\tabcolsep) * \real{0.1205}}
  >{\raggedright\arraybackslash}p{(\linewidth - 6\tabcolsep) * \real{0.3735}}
  >{\raggedright\arraybackslash}p{(\linewidth - 6\tabcolsep) * \real{0.3735}}
  >{\raggedright\arraybackslash}p{(\linewidth - 6\tabcolsep) * \real{0.1325}}@{}}
\toprule\noalign{}
\begin{minipage}[b]{\linewidth}\raggedright
Scenario
\end{minipage} & \begin{minipage}[b]{\linewidth}\raggedright
AI System (Time to Detection)
\end{minipage} & \begin{minipage}[b]{\linewidth}\raggedright
Fail2Ban (Time to Detection)
\end{minipage} & \begin{minipage}[b]{\linewidth}\raggedright
Advantage
\end{minipage} \\
\midrule\noalign{}
\endhead
\bottomrule\noalign{}
\endlastfoot
Basic Brute-Force & \textless{} 1 minute (first window) &
\textasciitilde5--10 seconds (5 failures) & Fail2Ban faster* \\
Distributed & \textless{} 2 minutes & Never (\textless{} 5 failures/IP)
& AI only \\
Low-and-Slow & 2--3 min (EARLY\_WARNING) & Never (\textless{} 5
failures/findtime) & AI only \\
Credential Stuffing & \textless{} 1 minute & \textasciitilde30--60
seconds & Comparable \\
Dictionary Attack & \textless{} 1 minute & \textasciitilde5--10 seconds
& Fail2Ban faster* \\
\end{longtable}
}

*For high-speed attacks (Scenarios 1, 5), Fail2Ban detects faster
because it triggers immediately upon the 5th failure, while the AI
system requires at least one full 5-minute window to compute meaningful
features. However, the AI system provides detection for Scenarios 2 and
3, which Fail2Ban cannot detect at all.

\subsection{4.9.3 Analysis}\label{analysis}

The comparison reveals that the AI system and Fail2Ban are complementary
rather than competitive:

\begin{itemize}
\tightlist
\item
  \textbf{Fail2Ban excels} at detecting high-speed, single-IP attacks
  (Scenarios 1, 5) with faster detection time due to its event-by-event
  counting mechanism.
\item
  \textbf{The AI system excels} at detecting distributed attacks
  (Scenario 2), low-and-slow attacks (Scenario 3), and providing nuanced
  detection through the two-level alerting mechanism.
\item
  \textbf{The integrated architecture} (both systems operating in
  parallel) provides the best of both capabilities: Fail2Ban provides
  rapid response to obvious high-speed attacks, while the AI system
  detects sophisticated attacks that evade Fail2Ban entirely.
\end{itemize}

This complementary architecture is consistent with the defense-in-depth
principle advocated by NIST SP 800-53 {[}108{]}, which recommends
multiple layers of security controls operating at different abstraction
levels.

\section{4.10 Comparison with
State-of-the-Art}\label{comparison-with-state-of-the-art}

\subsection{4.10.1 Literature Comparison}\label{literature-comparison}

The results of this study are compared with 12 published studies in the
field of SSH brute-force attack detection and network intrusion
detection. The comparison is structured to highlight the methodological
and performance differences.

{[}Table 4.23: Comprehensive comparison with state-of-the-art studies{]}

{\def\LTcaptype{none} % do not increment counter
\begin{longtable}[]{@{}
  >{\raggedright\arraybackslash}p{(\linewidth - 14\tabcolsep) * \real{0.1029}}
  >{\raggedright\arraybackslash}p{(\linewidth - 14\tabcolsep) * \real{0.0882}}
  >{\raggedright\arraybackslash}p{(\linewidth - 14\tabcolsep) * \real{0.1176}}
  >{\raggedright\arraybackslash}p{(\linewidth - 14\tabcolsep) * \real{0.1324}}
  >{\raggedright\arraybackslash}p{(\linewidth - 14\tabcolsep) * \real{0.1471}}
  >{\raggedright\arraybackslash}p{(\linewidth - 14\tabcolsep) * \real{0.1618}}
  >{\raggedright\arraybackslash}p{(\linewidth - 14\tabcolsep) * \real{0.1471}}
  >{\raggedright\arraybackslash}p{(\linewidth - 14\tabcolsep) * \real{0.1029}}@{}}
\toprule\noalign{}
\begin{minipage}[b]{\linewidth}\raggedright
Study
\end{minipage} & \begin{minipage}[b]{\linewidth}\raggedright
Year
\end{minipage} & \begin{minipage}[b]{\linewidth}\raggedright
Method
\end{minipage} & \begin{minipage}[b]{\linewidth}\raggedright
Dataset
\end{minipage} & \begin{minipage}[b]{\linewidth}\raggedright
Features
\end{minipage} & \begin{minipage}[b]{\linewidth}\raggedright
Real-time
\end{minipage} & \begin{minipage}[b]{\linewidth}\raggedright
F1-Score
\end{minipage} & \begin{minipage}[b]{\linewidth}\raggedright
Notes
\end{minipage} \\
\midrule\noalign{}
\endhead
\bottomrule\noalign{}
\endlastfoot
Sperotto et al.~{[}109{]} & 2010 & Flow-based HMM & University SSH logs
& Flow features & No & \textasciitilde0.85 & Pioneering SSH flow
analysis \\
Hellemons et al.~{[}65{]} & 2012 & Flow-based rules & CESNET & Flow
statistics & Partial & \textasciitilde0.88 & Time-window approach \\
Javed \& Paxson {[}72{]} & 2013 & Timing analysis & Live SSH & Timing
features & Yes & \textasciitilde0.90 & Established timing as key
discriminator \\
Najafabadi et al.~{[}110{]} & 2015 & Deep learning (AE) & NSL-KDD & PCA
features & No & 0.87 & Autoencoder anomaly detection \\
Kim et al.~{[}111{]} & 2019 & Random Forest & NSL-KDD & 41 KDD features
& No & 0.92 & Supervised, benchmark dataset \\
Ahmed et al.~{[}112{]} & 2020 & Autoencoder + LSTM & CICIDS2017 & Flow +
temporal & Yes & 0.89 & Deep learning, real-time \\
Moustafa et al.~{[}113{]} & 2021 & Ensemble (RF+XGB) & UNSW-NB15 & 49
features & No & 0.93 & Supervised ensemble \\
Nassif et al.~{[}114{]} & 2021 & ML Survey & Multiple & Varies & Varies
& -- & Comprehensive survey \\
Ferrag et al.~{[}115{]} & 2022 & Federated learning & CIC-ToN-IoT & DNN
features & Yes & 0.91 & Privacy-preserving \\
Thakkar \& Lohiya {[}116{]} & 2022 & IF + sampling & CICIDS2017 & Flow
features & No & 0.90 & Isolation Forest on IDS dataset \\
Sarker et al.~{[}117{]} & 2023 & ML + DL hybrid & Custom IoT &
Multi-modal & Partial & 0.92 & Hybrid architecture \\
Kumar \& Lim {[}118{]} & 2024 & Transformer-IDS & CICIDS2017 &
Attention-based & Yes & 0.94 & State-of-the-art, high compute \\
\textbf{This study} & \textbf{2026} & \textbf{IF + EWMA-AP} &
\textbf{Real SSH (honeypot + sim)} & \textbf{14 SSH-specific} &
\textbf{Yes} & \textbf{0.9374} & \textbf{End-to-end system} \\
\end{longtable}
}

\subsection{4.10.2 Comparative Analysis}\label{comparative-analysis}

The comparison highlights several distinctive characteristics of this
study:

\textbf{Dataset realism.} The majority of compared studies (8 of 12) use
benchmark datasets (NSL-KDD, CICIDS2017, UNSW-NB15) that are widely
acknowledged to have significant limitations: NSL-KDD is derived from
DARPA 1998/1999 network captures that do not reflect modern traffic
patterns {[}73{]}; CICIDS2017, while more recent, uses synthetically
generated attacks that may not capture the full diversity of real-world
attack behavior {[}74{]}. This study uses real SSH authentication logs
from a live honeypot (capturing authentic attack patterns from 679
unique attacker IPs) combined with realistic simulation data, providing
a more representative evaluation.

\textbf{Feature domain specificity.} Most compared studies use generic
network flow features (packet counts, byte counts, flow duration) or
protocol-agnostic statistical features. This study uses 14 features
specifically designed for SSH authentication behavior, including
domain-specific features such as session\_duration\_mean,
invalid\_user\_ratio, and pam\_failure\_escalation that capture the
unique characteristics of SSH brute-force attacks. Javed and Paxson
{[}72{]} demonstrated the superiority of domain-specific features over
generic features for SSH attack detection, a finding that this study
confirms and extends.

\textbf{Semi-supervised vs.~supervised.} Several high-performing studies
(Kim et al.~{[}111{]}, Moustafa et al.~{[}113{]}) use supervised
methods, which achieve high accuracy when the test data contains attack
types present in the training data but cannot detect novel attack types.
This study's semi-supervised approach sacrifices a small amount of
accuracy on known attack types in exchange for the ability to detect any
attack whose behavioral profile deviates from normal SSH usage --- a
critical capability in production environments where novel attack
techniques are constantly emerging.

\textbf{End-to-end system.} Most compared studies report offline
evaluation results without addressing the practical challenges of
real-time deployment, system integration, or automated response. This
study implements and evaluates a complete system comprising 9 Docker
services, with end-to-end latency under 2 seconds, demonstrating
production readiness. Only Ahmed et al.~{[}112{]}, Ferrag et
al.~{[}115{]}, and Kumar and Lim {[}118{]} report real-time capability,
but none implement an integrated detection-to-prevention pipeline with
Fail2Ban integration.

\textbf{Dynamic thresholding.} None of the compared studies implement an
adaptive dynamic threshold for SSH-specific anomaly detection. The
EWMA-Adaptive Percentile mechanism is a methodological contribution that
addresses the well-documented limitation of static thresholds in
non-stationary environments {[}89, 90{]}.

\textbf{F1-Score contextualization.} The F1-Score of 0.9374 (optimized
Isolation Forest) is competitive with the highest-performing studies.
Kumar and Lim {[}118{]} achieve 0.94 using Transformer-based deep
learning, but at significantly higher computational cost and with the
requirement for labeled training data (supervised). Moustafa et
al.~{[}113{]} achieve 0.93 using a supervised ensemble, also requiring
labeled attack data. This study achieves comparable F1-Score using a
computationally lightweight semi-supervised method that does not require
labeled attack data for training.

\subsection{4.10.3 Positioning}\label{positioning}

This study occupies a distinct position in the literature: it is one of
few studies that combines (1) real-world SSH-specific data, (2)
semi-supervised anomaly detection, (3) dynamic adaptive thresholding,
(4) an end-to-end deployed system, and (5) evaluation against diverse
attack scenarios. While individual components have been explored in
prior work, the integrated combination is novel and addresses the
research-to-deployment gap identified by Sommer and Paxson {[}119{]} as
the primary barrier to practical adoption of machine learning in
intrusion detection.

\section{4.11 Interpretation and
Implications}\label{interpretation-and-implications}

\subsection{4.11.1 Theoretical
Implications}\label{theoretical-implications}

The experimental results support several theoretical conclusions:

\textbf{The semi-supervised paradigm is sufficient for SSH brute-force
detection.} All three models achieve recall \(\geq 0.9675\) after
optimization, confirming that learning from normal data alone is
adequate to identify the vast majority of SSH brute-force attacks. This
validates the theoretical predictions of Chandola et al.~{[}47{]}
regarding the applicability of novelty detection to cybersecurity
domains where anomalies are behaviorally distinct from normal
operations.

\textbf{Temporal features dominate count-based features for SSH attack
discrimination.} The permutation importance analysis (Section 4.4)
demonstrates that session\_duration\_mean, min\_inter\_attempt\_time,
and mean\_inter\_attempt\_time are the primary discriminators, while
traditional count-based features (fail\_count, invalid\_user\_count)
have negligible importance. This finding extends the work of Javed and
Paxson {[}72{]} by quantifying the relative importance of specific
temporal features in the context of unsupervised anomaly detection, and
it challenges the design assumptions underlying traditional count-based
detection tools such as Fail2Ban.

\textbf{Dynamic thresholding provides operational advantages over static
thresholds.} The EWMA-Adaptive Percentile mechanism enables three
capabilities absent from static approaches: adaptation to distribution
shifts, burst false positive reduction, and cumulative early warning for
slow attacks. These advantages are most pronounced in non-stationary
operating environments --- precisely the conditions encountered in
real-world production deployments.

\subsection{4.11.2 Practical Implications}\label{practical-implications}

\textbf{For security operations teams.} The system can be deployed as a
complement to existing Fail2Ban installations, providing coverage for
the attack types (distributed, low-and-slow) that Fail2Ban cannot
detect. The two-level alerting mechanism (EARLY\_WARNING and ALERT)
provides operational flexibility, allowing security teams to configure
the response policy according to their risk tolerance.

\textbf{For SSH server administrators.} The feature importance results
suggest that monitoring session duration and inter-attempt timing is
more informative than monitoring failure counts alone. Administrators
who cannot deploy the full system can still improve their detection
capability by adding session-duration-based rules to their existing
Fail2Ban or OSSEC configurations.

\textbf{For the research community.} The dataset and methodology provide
a template for future studies that seek to bridge the gap between
offline algorithm evaluation and real-world system deployment. The
finding that temporal features dominate count-based features should
inform the feature engineering decisions of future SSH security
research.

\subsection{4.11.3 Limitations of the
Results}\label{limitations-of-the-results}

Several limitations should be noted when interpreting the experimental
results:

\begin{enumerate}
\def\labelenumi{\arabic{enumi}.}
\item
  \textbf{The 29\% FPR of the optimized IF model} means that
  approximately 29\% of normal windows are still incorrectly flagged as
  attacks. While this is a dramatic improvement over the 77\% baseline
  FPR, it represents an ongoing cost in terms of alert volume that
  security teams must manage. Further reduction of the FPR without
  sacrificing recall remains an open challenge.
\item
  \textbf{The low-and-slow detection rate of 70--85\%} represents the
  system's weakest performance. Extremely slow attacks (one attempt
  every 10+ minutes) may produce feature vectors that are
  indistinguishable from normal activity within a single 5-minute
  window, and even the EWMA accumulation mechanism may take multiple
  windows to reach the detection threshold.
\item
  \textbf{The evaluation was conducted on a specific dataset} with a
  specific normal-to-attack ratio, attack type distribution, and
  environmental characteristics. Performance may vary in different
  environments with different traffic patterns, user populations, or
  attack sophistication levels.
\item
  \textbf{The comparison with state-of-the-art studies} is necessarily
  imperfect because the studies use different datasets, different
  evaluation protocols, and different definitions of ``attack.'' Direct
  head-to-head comparison on the same dataset would provide more
  rigorous benchmarking, but such standardized SSH-specific benchmarks
  do not yet exist.
\end{enumerate}

\subsection{4.11.4 Summary of Key
Findings}\label{summary-of-key-findings}

{[}Table 4.24: Summary of key experimental findings{]}

{\def\LTcaptype{none} % do not increment counter
\begin{longtable}[]{@{}
  >{\raggedright\arraybackslash}p{(\linewidth - 4\tabcolsep) * \real{0.2812}}
  >{\raggedright\arraybackslash}p{(\linewidth - 4\tabcolsep) * \real{0.3125}}
  >{\raggedright\arraybackslash}p{(\linewidth - 4\tabcolsep) * \real{0.4062}}@{}}
\toprule\noalign{}
\begin{minipage}[b]{\linewidth}\raggedright
Finding
\end{minipage} & \begin{minipage}[b]{\linewidth}\raggedright
Evidence
\end{minipage} & \begin{minipage}[b]{\linewidth}\raggedright
Significance
\end{minipage} \\
\midrule\noalign{}
\endhead
\bottomrule\noalign{}
\endlastfoot
Semi-supervised detection achieves \(\geq 96.75\%\) recall & Table 4.12,
all three models & Validates paradigm for SSH security \\
Temporal features are primary discriminators & Table 4.11, top 5
features = temporal/session & Challenges count-based detection
assumptions \\
Optimization reduces FPR by 48 pp & Table 4.14, 76.92\% \(\rightarrow\)
29.00\% & Critical for production viability \\
Dynamic threshold enables early warning & Section 4.6, 2--3 min
detection of slow attacks & Novel capability absent from static
methods \\
AI detects attacks that Fail2Ban cannot & Table 4.22, Scenarios 2 and 3
& Justifies complementary deployment \\
End-to-end latency \textless{} 2 seconds & Table 4.16 & Confirms
real-time capability \\
F1=0.9374 competitive with state-of-the-art & Table 4.23, comparison
with 12 studies & Validates approach against literature \\
5 attack scenarios with 70--100\% detection & Table 4.21 & Comprehensive
coverage evaluation \\
\end{longtable}
}

\chapter{CHAPTER 5: DISCUSSION}\label{chapter-5-discussion}

\section{5.1 Restatement of Research
Problem}\label{restatement-of-research-problem}

The proliferation of Internet-facing SSH services has made brute-force
login attacks one of the most persistent threats to server
infrastructure worldwide. According to Rapid7 {[}4{]}, SSH (port 22)
consistently ranks among the top three targeted services, with
brute-force attempts accounting for a substantial proportion of all
observed malicious traffic. Conventional countermeasures such as
Fail2Ban {[}27{]} employ static, count-based thresholds --- typically
blocking an IP address after a fixed number of failed login attempts
within a predefined time window. While effective against naive,
high-speed attacks, this paradigm suffers from three fundamental
limitations: (i) it cannot detect low-and-slow attacks in which the
adversary deliberately spaces attempts below the threshold rate; (ii) it
is blind to distributed attacks that spread attempts across multiple
source IPs; and (iii) it offers no predictive capability, responding
only after the threshold has already been exceeded.

This thesis was motivated by the question: \emph{How can an AI-based SSH
brute-force detection and prevention system be constructed that provides
early attack prediction, adapts dynamically to changing traffic
patterns, and integrates fully into modern security monitoring
infrastructure?} Five specific objectives were formulated to
operationalize this question:

\begin{enumerate}
\def\labelenumi{\arabic{enumi}.}
\tightlist
\item
  Design a comprehensive behavioral feature set that captures the
  multidimensional differences between legitimate SSH activity and
  brute-force attacks.
\item
  Train and systematically evaluate three unsupervised anomaly detection
  models --- Isolation Forest, Local Outlier Factor, and One-Class SVM
  --- on real-world SSH data.
\item
  Develop an adaptive EWMA-Adaptive Percentile dynamic threshold that
  enables early warning and self-calibration.
\item
  Build a complete, containerized end-to-end system integrating log
  ingestion, feature extraction, anomaly scoring, dynamic thresholding,
  and automated response within a 9-service Docker architecture.
\item
  Validate the system against five realistic attack scenarios of
  escalating difficulty.
\end{enumerate}

Four subsidiary research questions (RQ1--RQ4) were posed in Chapter 1.
The following sections evaluate the extent to which each objective was
met, interpret the principal findings, position the results within the
existing literature, discuss practical deployment implications, and
identify limitations and threats to validity.

\section{5.2 Achievement of Research
Objectives}\label{achievement-of-research-objectives}

\subsection{5.2.1 Objective 1: Behavioral Feature Set
Design}\label{objective-1-behavioral-feature-set-design}

The 14-feature behavioral feature set described in Chapter 3 (Table 3.4)
was successfully designed, implemented, and validated. The features span
five functional groups --- authentication count and rate,
username-related, connection and temporal, network and session, and
attack indicator features --- extracted over 5-minute sliding windows
per source IP address. Feature importance analysis using permutation
importance on the optimized Isolation Forest model (Section 4.3)
confirmed that the feature set provides strong discriminative power,
with the top five features collectively accounting for more than 15\% of
total importance. Crucially, the analysis revealed that temporal
features dominate discrimination: session\_duration\_mean (5.50\%),
min\_inter\_attempt\_time (3.86\%), mean\_inter\_attempt\_time (2.61\%),
and std\_inter\_attempt\_time (1.64\%) occupy four of the top five
positions (Table 4.5, Figure 4.6). This result validates the design
decision to include fine-grained timing statistics alongside the
count-based features traditionally used in tools such as Fail2Ban.

\subsection{5.2.2 Objective 2: Model Training and
Evaluation}\label{objective-2-model-training-and-evaluation}

Three unsupervised anomaly detection algorithms were implemented,
optimized, and systematically compared on a hybrid dataset of 174,250
SSH log lines (119,729 honeypot attack lines and 54,521 simulated normal
lines), yielding 22,396 processed samples. After optimization
(non-overlapping windows, derived features, contamination tuning), the
following results were obtained:

{\def\LTcaptype{none} % do not increment counter
\begin{longtable}[]{@{}lllll@{}}
\toprule\noalign{}
Model & Accuracy & F1-Score & Recall & FPR \\
\midrule\noalign{}
\endhead
\bottomrule\noalign{}
\endlastfoot
Isolation Forest & 90.31\% & 93.74\% & 96.75\% & 29.00\% \\
LOF & 83.22\% & 89.94\% & 100.00\% & 67.10\% \\
OCSVM & 91.38\% & 94.55\% & 100.00\% & 33.42\% \\
\end{longtable}
}

All three models exceed the minimum targets of F1 \textgreater{} 85\%
and Recall \textgreater{} 95\% established in Section 1.5. The Isolation
Forest was selected as the primary production model for reasons
elaborated in Section 5.3.2 below. The systematic three-model comparison
on the same real-world SSH dataset --- rather than on synthetic
benchmarks such as NSL-KDD {[}53{]} or CICIDS2017 {[}32{]} ---
constitutes an empirical contribution to the algorithm selection
literature for SSH anomaly detection.

\subsection{5.2.3 Objective 3: Dynamic
Threshold}\label{objective-3-dynamic-threshold}

The EWMA-Adaptive Percentile dynamic threshold (alpha = 0.3,
base\_percentile = 95, sensitivity\_factor = 1.5, lookback = 100) was
successfully developed and evaluated (Section 4.4, Table 4.9, Figure
4.12). Compared with a static threshold at the 95th percentile, the
dynamic threshold reduced burst false positives while maintaining
equivalent recall. Most importantly, the EWMA accumulation mechanism
enables early warning for low-and-slow attacks after only 3--5 attempts
(approximately 2--3 minutes), whereas Fail2Ban with default settings
(maxretry = 5, findtime = 600 s) fails to detect such attacks entirely
when attempts are spaced at intervals of 30--120 seconds {[}27{]}. The
two-level detection mechanism (EARLY\_WARNING at 67\% of the ALERT
threshold, ALERT triggering Fail2Ban) provides graduated response,
avoiding unnecessary service disruption.

\subsection{5.2.4 Objective 4: Integrated
System}\label{objective-4-integrated-system}

The complete system was implemented as a Docker Compose deployment
comprising 9 services: FastAPI (API server and ML inference), React
(monitoring dashboard), Elasticsearch (storage and indexing), Logstash
(log normalization), Kibana (visualization), Fail2Ban (automated IP
blocking), Redis (caching and session management), PostgreSQL
(configuration persistence), and Nginx (reverse proxy and load
balancing). The five-stage real-time detection pipeline (Ingestion,
Aggregation, Scoring, Decision, Action) achieves an end-to-end latency
of under 2 seconds, with AI processing stages completing in under 100 ms
(Table 4.14). The system is portable, reproducible, and deployable with
a single command (\texttt{docker-compose\ up}), directly addressing the
research-to-deployment gap identified by Sommer and Paxson {[}13{]} as a
persistent barrier to practical ML adoption in network security.

\subsection{5.2.5 Objective 5: Attack Scenario
Evaluation}\label{objective-5-attack-scenario-evaluation}

Five attack scenarios were designed, executed, and analyzed (Table 4.10,
Table 4.12, Figure 4.14): (1) basic brute-force, (2) distributed
brute-force, (3) low-and-slow, (4) credential stuffing, and (5)
dictionary attacks. Basic brute-force and dictionary attacks were
detected immediately (100\% detection within the first 1-minute window).
Distributed attacks and credential stuffing achieved detection rates
exceeding 90\% and 95\%, respectively. The low-and-slow scenario --- the
most challenging case --- was partially addressed through the EWMA
accumulation mechanism, with early warnings issued after 3--5 attempts.
Additionally, 51 unit tests across all system components passed,
confirming functional correctness and integration integrity.

\section{5.3 Interpretation of Key
Findings}\label{interpretation-of-key-findings}

\subsection{5.3.1 Why Timing Features
Dominate}\label{why-timing-features-dominate}

The finding that temporal features outperform count-based features for
SSH brute-force detection (Section 4.3, Table 4.5) is both statistically
robust and intuitively grounded. Session\_duration\_mean (importance =
5.50\%) is the single most discriminative feature because brute-force
sessions are inherently short: each failed password attempt results in
rapid disconnection, producing session durations typically under 5
seconds, while legitimate interactive sessions last minutes to hours.
This bimodal gap is a structural property of the SSH protocol {[}2,
15{]} and cannot be eliminated by the attacker without fundamentally
changing the attack strategy.

Min\_inter\_attempt\_time (3.86\%) captures the minimum spacing between
consecutive authentication attempts from a single IP. Automated tools
such as Hydra and Medusa generate attempts at machine speed, producing
inter-attempt intervals of milliseconds to low seconds, whereas human
operators exhibit intervals of seconds to minutes with high variance.
Even sophisticated attackers who introduce random delays to evade
count-based detection find it difficult to perfectly mimic the
stochastic temporal signature of legitimate users {[}20, 62{]}.

The low importance of traditional count-based features (fail\_count,
connection\_count) is explained by the training methodology: the model
is trained exclusively on normal data where these features are
consistently near zero. In the anomaly detection paradigm, the model
learns the distribution of normal behavior; features with near-zero
variance in normal data contribute little to the anomaly score
computation because all test points --- both normal and attack --- are
evaluated relative to the learned normal distribution. Temporal
features, by contrast, exhibit substantial variance even within normal
data (different users type at different speeds, maintain different
session durations), enabling the model to learn a nuanced boundary.

This finding corroborates and extends the work of Javed and Paxson
{[}62{]}, who demonstrated that timing characteristics are the most
effective discriminator for SSH brute-force traffic, and of Starov et
al.~{[}54{]}, who showed that temporal behavioral analysis outperforms
volume-based features for SSH attack detection. Our contribution is the
quantification of specific feature importance values within an Isolation
Forest framework, providing actionable guidance for feature engineering
in future detection systems.

\subsection{5.3.2 Why Isolation Forest Is Chosen Despite OCSVM Having
Higher
F1}\label{why-isolation-forest-is-chosen-despite-ocsvm-having-higher-f1}

The OCSVM achieved the highest F1-score (94.55\%) and accuracy (91.38\%)
among the three models evaluated. Nevertheless, the Isolation Forest (F1
= 93.74\%, accuracy = 90.31\%) was selected as the primary production
model. This decision was driven by three operational considerations that
outweigh the 0.81 percentage-point F1 gap.

First, \textbf{computational efficiency}. Isolation Forest operates in
O(n log n) time for both training and inference, compared with O(n\^{}2)
to O(n\^{}3) for OCSVM {[}35, 39{]}. In a production system processing
hundreds or thousands of SSH sessions per minute, this difference
translates directly into lower latency and reduced resource consumption.
Table 4.15 confirms that IF achieves substantially higher throughput
than OCSVM.

Second, \textbf{anomaly score properties}. Isolation Forest produces
smooth, continuous anomaly scores derived from the average path length
across the tree ensemble {[}36{]}. These scores are well-suited for
input to the EWMA-Adaptive Percentile dynamic threshold, which requires
a continuous signal to track trends and compute running percentiles.
OCSVM, by contrast, produces a signed distance to the decision boundary
that is less granular near the boundary region, reducing the sensitivity
of the dynamic threshold to subtle shifts in traffic patterns.

Third, \textbf{distributional assumptions}. OCSVM with an RBF kernel
implicitly assumes that normal data forms a single compact region in the
kernel-mapped feature space {[}39, 40{]}. This assumption holds for the
current dataset but may be violated in heterogeneous production
environments where normal SSH behavior spans multiple distinct clusters
(e.g., administrators, automated scripts, interactive users). Isolation
Forest makes no distributional assumptions, operating purely on the
isolability principle {[}35{]}, making it more robust to multimodal
normal distributions.

This reasoning aligns with the recommendations of Goldstein and Uchida
{[}41{]}, whose comparative evaluation of unsupervised anomaly detection
algorithms concluded that Isolation Forest offers the best trade-off
between detection performance and computational cost for datasets with
moderate dimensionality.

\subsection{5.3.3 Dynamic Threshold vs.~Static Threshold
Advantage}\label{dynamic-threshold-vs.-static-threshold-advantage}

Static thresholds, as employed by Fail2Ban and most rule-based intrusion
detection systems {[}27, 29{]}, create a fixed decision boundary that
does not respond to changes in the underlying data distribution. Three
specific failure modes of static thresholds were identified in this
research.

First, \textbf{baseline drift}. SSH traffic patterns vary systematically
between business hours and off-hours, between weekdays and weekends, and
across organizational events (onboarding, audits, system migrations). A
static threshold calibrated for average conditions produces excessive
false positives during high-activity periods and may miss attacks during
low-activity periods. The EWMA component (alpha = 0.3) tracks the moving
baseline, automatically widening the threshold during periods of
legitimately elevated activity and narrowing it during quiet periods.

Second, \textbf{burst false positives}. Legitimate but anomalous events
(a new employee's first SSH session, an automated script changing
execution timing) produce transient anomaly score spikes that exceed
static thresholds. The Adaptive Percentile component computes the
threshold from the recent score distribution (lookback = 100), absorbing
individual spikes without triggering alerts unless the elevated pattern
persists.

Third, \textbf{low-and-slow evasion}. Adversaries who space attempts
below the static threshold rate (e.g., one attempt every 2 minutes
against Fail2Ban's default of 5 attempts in 10 minutes) evade
count-based detection entirely. The EWMA accumulation mechanism provides
a fundamentally different detection modality: each anomalous score, even
if individually below the alert threshold, incrementally raises the EWMA
value. After 3--5 successive anomalous attempts (approximately 2--3
minutes), the accumulated EWMA value crosses the early warning
threshold, triggering notification. This capability was confirmed
experimentally in Scenario 3 (Table 4.11) and represents the most
significant practical advantage of the dynamic threshold approach.

The dynamic threshold also supports the two-level graduated response
mechanism, where EARLY\_WARNING events (EWMA exceeding 67\% of the alert
threshold) are logged and reported without blocking, while ALERT events
trigger automated IP blocking via Fail2Ban. This graduated response
addresses the operational concern that overly aggressive blocking
disrupts legitimate services --- a concern frequently cited by security
operations teams {[}49{]}.

\subsection{5.3.4 FPR Analysis and Acceptable Trade-off in Security
Context}\label{fpr-analysis-and-acceptable-trade-off-in-security-context}

The optimized Isolation Forest model exhibits a false positive rate of
29.00\%, meaning approximately 29 out of every 100 normal SSH sessions
are incorrectly flagged as anomalous. In many machine learning
application domains (e.g., spam filtering, fraud detection), an FPR of
this magnitude would be unacceptable. However, in the cybersecurity
context, the cost asymmetry between false positives and false negatives
must be considered {[}8, 13{]}.

The cost of a false negative --- a missed attack that allows an
adversary to gain unauthorized access --- includes potential data
exfiltration, lateral movement, ransomware deployment, regulatory
penalties, and reputational damage. According to Morgan {[}1{]},
cybercrime damages were projected to reach \$10.5 trillion annually by
2025. The cost of a false positive, by contrast, is an unnecessary
security alert that consumes analyst time and may temporarily
inconvenience a legitimate user.

Given this asymmetry, the operating point of 96.75\% recall and 29.00\%
FPR is defensible. Moreover, the two-level detection mechanism mitigates
the practical impact of false positives: EARLY\_WARNING events are
logged without blocking, and only high-confidence ALERT events result in
IP blocking. The dynamic threshold's self-calibration further reduces
effective FPR over time as the system adapts to the specific deployment
environment's normal baseline.

Several avenues exist for further FPR reduction. Ensemble methods
combining IF, LOF, and OCSVM through majority voting could reduce
individual model biases {[}41{]}. Deep learning approaches such as
LSTM-Autoencoders {[}57{]} can model complex temporal dependencies that
tree-based methods may miss. Integration of external threat intelligence
{[}28{]} would provide additional classification context, effectively
lowering the threshold for known malicious IPs while raising it for IPs
with clean reputations. Active learning, in which administrator feedback
on alerts is used to iteratively refine the decision boundary,
represents another promising direction {[}63{]}.

\section{5.4 Comparison with Existing
Literature}\label{comparison-with-existing-literature}

This section positions the results of the current study against the
findings reported in the existing literature. Table 5.2 provides a
systematic comparison.

\textbf{Sperotto et al.~(2010) {[}21{]}.} Proposed flow-based intrusion
detection using IP flow records from the DARPA dataset. Their approach
achieved moderate detection rates but operated in an offline,
batch-processing mode without real-time capability or early prediction.
Our system advances beyond flow-based methods by operating on raw SSH
authentication logs, enabling extraction of fine-grained temporal
features (e.g., session\_duration\_mean, min\_inter\_attempt\_time) that
are unavailable in aggregated flow records. Furthermore, our system
provides real-time detection with sub-2-second latency.

\textbf{Najafabadi et al.~(2015) {[}7{]}.} Employed aggregated NetFlow
data with machine learning for SSH brute-force detection, reporting
detection rates above 90\%. However, their approach required labeled
training data (supervised learning) and did not address low-and-slow or
distributed attack variants. Our semi-supervised approach, requiring
only normal data for training, eliminates the dependence on labeled
attack samples and demonstrates detection capability across five attack
scenarios including low-and-slow.

\textbf{Buczak and Guven (2016) {[}8{]}.} Surveyed data mining and
machine learning methods for intrusion detection, identifying the lack
of real-world evaluation datasets, the absence of real-time deployment,
and concept drift as the three principal challenges. The current study
addresses the first two challenges directly: the evaluation uses real
honeypot data rather than synthetic benchmarks, and the system operates
in real-time with demonstrated sub-2-second latency. Concept drift is
acknowledged as a limitation and addressed in the future work
directions.

\textbf{Kim et al.~(2019) {[}31{]}.} Achieved an F1-score of 0.92 using
Random Forest on the NSL-KDD dataset for network intrusion detection.
The current study achieves a comparable F1-score of 0.9374 with
Isolation Forest on real SSH data, without requiring labeled attack
samples for training. Moreover, the NSL-KDD dataset has been criticized
for lack of representativeness of modern attack patterns {[}53{]}; our
use of contemporary honeypot data provides a more ecologically valid
evaluation.

\textbf{Tsai et al.~(2019) {[}6{]}.} Studied SSH brute-force defense
mechanisms and proposed a hybrid approach combining rate limiting with
behavioral analysis. Their system achieved effective detection for
high-speed attacks but did not address low-and-slow variants. Our
EWMA-Adaptive Percentile threshold directly addresses this gap,
detecting low-and-slow attacks after 3--5 attempts.

\textbf{Ahmed et al.~(2020) {[}32{]}.} Applied autoencoder-based anomaly
detection to the CICIDS2017 dataset, achieving an F1-score of 0.89 with
real-time processing capability. The current study achieves a higher
F1-score (0.9374) using a computationally simpler model (Isolation
Forest vs.~autoencoder) and additionally provides the early prediction
capability through the dynamic threshold that Ahmed et al.~did not
investigate.

\textbf{Nassif et al.~(2021) {[}31{]}.} Conducted a comprehensive ML
survey for intrusion detection, noting that the gap between experimental
performance and production deployment remains the principal obstacle.
The current study directly addresses this gap by delivering a complete,
containerized, deployable system rather than an offline experimental
evaluation.

\textbf{Pang et al.~(2021) {[}33{]}.} Reviewed deep learning for anomaly
detection, recommending that future systems combine traditional ML
efficiency with deep learning expressiveness. The current study
demonstrates that traditional ML (Isolation Forest) achieves competitive
performance (F1 = 93.74\%) with significantly lower computational
requirements than deep learning, while the future work section (Chapter
6) outlines how deep learning can be integrated for further improvement.

\textbf{Satoh et al.~(2022) {[}57{]}.} Applied LSTM-based deep learning
to SSH dictionary attack detection, achieving a recall of 97.2\%. The
current study achieves comparable recall (96.75\%) with a significantly
simpler model that does not require GPU resources or sequential
training. However, the LSTM approach's ability to model long-range
temporal dependencies suggests that deep learning could complement the
current system for improved detection of sophisticated attack patterns.

\textbf{Kumari and Jain (2022) {[}52{]}.} Applied Isolation Forest for
IoT anomaly detection, reporting accuracy above 90\%. The current study
extends the IF methodology to the SSH domain with three specific
enhancements: the 14-feature behavioral feature set tailored to SSH
semantics, the EWMA-Adaptive Percentile dynamic threshold for adaptive
detection, and the complete integrated system architecture.

\textbf{Javed and Paxson (2013) {[}62{]}.} Demonstrated that timing
characteristics are the most effective factor for detecting stealthy,
distributed SSH brute-forcing. Our feature importance analysis (Table
4.5) empirically corroborates this finding within the Isolation Forest
framework, with four of the top five features being temporal. We extend
their work by quantifying specific importance values and embedding the
temporal features within a dynamic thresholding system that enables
early prediction.

\textbf{Starov et al.~(2019) {[}54{]}.} Proposed temporal behavioral
analysis for SSH brute-force detection and showed that temporal features
outperform volume-based features. The current study confirms this
finding and advances it by implementing a complete detection pipeline
with automated response and by evaluating against five attack scenarios
of varying sophistication.

The following table consolidates the comparison:

{\def\LTcaptype{none} % do not increment counter
\begin{longtable}[]{@{}
  >{\raggedright\arraybackslash}p{(\linewidth - 12\tabcolsep) * \real{0.0833}}
  >{\raggedright\arraybackslash}p{(\linewidth - 12\tabcolsep) * \real{0.0952}}
  >{\raggedright\arraybackslash}p{(\linewidth - 12\tabcolsep) * \real{0.1071}}
  >{\raggedright\arraybackslash}p{(\linewidth - 12\tabcolsep) * \real{0.2619}}
  >{\raggedright\arraybackslash}p{(\linewidth - 12\tabcolsep) * \real{0.2024}}
  >{\raggedright\arraybackslash}p{(\linewidth - 12\tabcolsep) * \real{0.1310}}
  >{\raggedright\arraybackslash}p{(\linewidth - 12\tabcolsep) * \real{0.1190}}@{}}
\toprule\noalign{}
\begin{minipage}[b]{\linewidth}\raggedright
Study
\end{minipage} & \begin{minipage}[b]{\linewidth}\raggedright
Method
\end{minipage} & \begin{minipage}[b]{\linewidth}\raggedright
Dataset
\end{minipage} & \begin{minipage}[b]{\linewidth}\raggedright
Labeled Data Required
\end{minipage} & \begin{minipage}[b]{\linewidth}\raggedright
Early Prediction
\end{minipage} & \begin{minipage}[b]{\linewidth}\raggedright
Real-time
\end{minipage} & \begin{minipage}[b]{\linewidth}\raggedright
F1-Score
\end{minipage} \\
\midrule\noalign{}
\endhead
\bottomrule\noalign{}
\endlastfoot
Sperotto et al.~{[}21{]} & Flow-based & DARPA & Yes & No & No & -- \\
Najafabadi et al.~{[}7{]} & ML + NetFlow & NetFlow & Yes & No & No &
\textgreater0.90 \\
Kim et al.~{[}31{]} & Random Forest & NSL-KDD & Yes & No & No & 0.92 \\
Ahmed et al.~{[}32{]} & Autoencoder & CICIDS2017 & No & No & Yes &
0.89 \\
Tsai et al.~{[}6{]} & Hybrid & Custom & -- & No & Yes & -- \\
Satoh et al.~{[}57{]} & LSTM & SSH logs & Yes & No & Yes & --
(R=0.972) \\
Kumari \& Jain {[}52{]} & IF & IoT & No & No & No & \textgreater0.90 \\
Javed \& Paxson {[}62{]} & Statistical & SSH traces & No & No & Offline
& -- \\
\textbf{This study} & \textbf{IF + EWMA} & \textbf{Real SSH} &
\textbf{No} & \textbf{Yes} & \textbf{Yes} & \textbf{0.9374} \\
\end{longtable}
}

Three distinguishing aspects position this work relative to the
literature. First, \textbf{dataset realism}: the use of real SSH
honeypot data (119,729 log lines from 679 unique IPs) rather than
synthetic benchmarks. Second, \textbf{methodological novelty}: the
combination of Isolation Forest with EWMA-Adaptive Percentile dynamic
thresholding for SSH early prediction has not been previously reported.
Third, \textbf{engineering completeness}: the system is implemented as a
deployable, containerized end-to-end solution, not merely an offline
experimental prototype.

\section{5.5 Practical Implications}\label{practical-implications-1}

\subsection{5.5.1 Deployment
Considerations}\label{deployment-considerations}

The experimental results carry several direct implications for security
operations teams evaluating AI-based SSH monitoring solutions.

\textbf{Feature engineering priorities.} The demonstrated primacy of
temporal features (Section 5.3.1) implies that security monitoring
infrastructure must preserve fine-grained timing information in SSH
logs. Many current logging configurations aggregate or discard timing
details, retaining only count-based summaries (e.g., ``X failed logins
from IP Y in Z minutes''). Organizations adopting the proposed approach
should ensure that authentication timestamps are logged at second or
sub-second granularity, that session start and end times are recorded,
and that log rotation policies do not truncate these fields. The
standard OpenSSH syslog format on Debian/Ubuntu (/var/log/auth.log) and
CentOS/RHEL (/var/log/secure) provides sufficient detail for the
14-feature extraction pipeline {[}18{]}.

\textbf{Integration with existing infrastructure.} The Docker-based
architecture was designed for compatibility with existing security
monitoring stacks. Organizations already using the ELK Stack for log
management can integrate the detection system by adding the FastAPI
service and configuring Logstash to forward SSH events to both
Elasticsearch and the detection API. The Fail2Ban integration leverages
the existing Fail2Ban installation, adding a custom jail configuration
for AI-triggered bans. This design philosophy --- augmenting rather than
replacing existing tools --- lowers the adoption barrier.

\textbf{Alert management and SOC workflows.} The two-level detection
mechanism (EARLY\_WARNING and ALERT) maps naturally to SOC triage
workflows. EARLY\_WARNING events can be routed to Tier 1 analysts for
monitoring, while ALERT events trigger automated response and are
escalated to Tier 2 for investigation. The React-based monitoring
dashboard provides real-time visibility into detection events, anomaly
score trends, and threshold behavior, supporting the analyst's
decision-making process. Security operations centers frequently cite
alert fatigue as a primary obstacle to effective monitoring {[}49{]};
the dynamic threshold's self-calibration and graduated response directly
address this concern.

\textbf{Resource requirements.} The system was evaluated on a standard
server configuration (Table 4.13), with the AI processing stages
completing in under 100 ms. The Docker-based deployment requires
approximately 4 GB of RAM for all 9 services (Table 4.16), well within
the capacity of a modest virtual machine or cloud instance. For
organizations with limited infrastructure budgets, the system can be
co-hosted on the SSH server itself, although dedicated deployment is
recommended for production use to avoid resource contention during
attack surges.

\subsection{5.5.2 Scalability
Considerations}\label{scalability-considerations}

The current architecture is designed for monitoring a single SSH server
or a small cluster of servers. For organizations with larger SSH
infrastructure, several scaling strategies are available:

\begin{itemize}
\tightlist
\item
  \textbf{Horizontal scaling.} Multiple FastAPI instances can be
  deployed behind the Nginx reverse proxy to distribute the detection
  workload across processing nodes.
\item
  \textbf{Centralized log collection.} Filebeat or Logstash agents
  deployed on multiple SSH servers can forward logs to a centralized
  Elasticsearch cluster, enabling the detection system to monitor
  multiple servers from a single deployment.
\item
  \textbf{Cloud-native deployment.} The Docker Compose configuration can
  be adapted for Kubernetes orchestration, enabling elastic scaling
  based on detection workload. Elasticsearch, the most
  resource-intensive component, can be deployed as a managed service
  (e.g., Elastic Cloud, Amazon OpenSearch) to offload operational
  overhead.
\end{itemize}

When the number of actively monitored IPs scales to thousands,
architectural refinements become necessary: replacing in-memory deques
with Redis Streams for per-IP window management, implementing batch
scoring instead of per-IP individual scoring, and deploying horizontal
scaling of detection workers.

\subsection{5.5.3 Cost-Benefit Analysis}\label{cost-benefit-analysis}

The total cost of deploying the system is dominated by infrastructure
costs (server provisioning, cloud instance fees) rather than software
licensing, as all components are open-source. For a typical
single-server deployment monitoring SSH on a Linux server with moderate
traffic (hundreds of sessions per day), a cloud instance with 4 vCPUs, 8
GB RAM, and 50 GB storage is sufficient, costing approximately
\$30--\$60/month on major cloud providers. The operational cost of the
29\% FPR --- primarily analyst time spent reviewing false positives ---
can be estimated at approximately 15--30 minutes per day for a
moderate-traffic server, assuming manual review of ALERT-level events.
This cost is substantially lower than the potential impact of a
successful SSH compromise, which can reach hundreds of thousands of
dollars in incident response, remediation, and business disruption costs
{[}1{]}.

\section{5.6 Limitations and Threats to
Validity}\label{limitations-and-threats-to-validity}

Six principal limitations and threats to validity must be acknowledged.

\textbf{L1: Concept drift.} The model is trained on a fixed dataset
collected over 5 days. Attack patterns evolve as adversaries develop new
tools and strategies, and normal user behavior shifts with
organizational changes {[}8{]}. The current model will gradually lose
effectiveness as the data distribution drifts. Mitigation: periodic
retraining on recently collected normal data. Future work (Section 6.2)
investigates online learning and incremental Isolation Forest techniques
for continuous adaptation without full retraining.

\textbf{L2: SSH key-based attacks.} This research focuses exclusively on
password-based brute-force attacks. Attacks using stolen SSH keys
produce ``Accepted publickey'' log entries rather than ``Failed
password'' patterns and are therefore invisible to the current feature
extraction pipeline {[}15{]}. In environments where key-based
authentication is the primary method, the system would need to be
extended with features capturing key-based anomalies (logins from new
IPs, unusual times, newly created key pairs).

\textbf{L3: Log format dependency.} The log parser is designed for the
standard syslog format used by OpenSSH on Debian/Ubuntu and CentOS/RHEL
{[}18{]}. Systems using non-standard log formats, alternative SSH
implementations (e.g., Dropbear), or centralized log management systems
that reformat entries would require parser adaptation. The modular
architecture isolates the parsing layer from feature extraction and
detection, so adaptation requires modifying only one component.

\textbf{L4: Training data diversity.} The training data is derived from
a single simulation source with 64 users exhibiting specific behavioral
patterns. The diversity of normal SSH behavior in production
environments --- including variations in work schedules, geographical
distribution, device diversity, and organizational policies --- exceeds
what a single simulation captures. In production deployments, collecting
normal behavior data from the target environment for a period of 1--2
weeks before activating detection is recommended to improve model
generalization.

\textbf{L5: False positive rate.} The 29.00\% FPR, while acceptable
given the cost asymmetry discussed in Section 5.3.4, remains a
significant rate of false alarms. In high-traffic environments
processing thousands of SSH sessions daily, this would generate a
substantial alert volume. The two-level detection mechanism mitigates
operational impact, but further FPR reduction through ensemble methods,
deep learning, or threat intelligence integration is a priority for
future work.

\textbf{L6: Evaluation scope.} The evaluation was conducted on data from
a single honeypot server over 5 days, with five designed attack
scenarios. While the attack scenarios cover the major categories of SSH
brute-force attacks (basic, distributed, low-and-slow, credential
stuffing, dictionary), the space of possible attack strategies is vast,
and novel techniques not represented in the evaluation could potentially
evade the system. A multi-site, longitudinal evaluation spanning weeks
or months would provide stronger evidence of robustness and
generalizability. The 51 unit tests confirm functional correctness but
do not substitute for extended operational validation.

\textbf{Threats to internal validity.} The primary threat is potential
data leakage between training and test sets. This was mitigated by
strict temporal separation: the RobustScaler was fit exclusively on the
training set (normal data) and applied to transform both training and
test sets. Non-overlapping windows in the optimized configuration
further reduce temporal correlation between consecutive feature vectors.

\textbf{Threats to external validity.} The generalizability of the
results to SSH environments with substantially different user
populations, usage patterns, or network configurations is unknown. The
semi-supervised approach requires that the training data be
representative of the target environment's normal behavior; deploying
the model trained on our simulation data in a radically different
environment without retraining would likely degrade performance. This is
a general limitation of all anomaly detection systems, not specific to
this work {[}9, 34{]}.

\textbf{Threats to construct validity.} The evaluation metrics
(accuracy, F1-score, recall, FPR, ROC-AUC) are standard in the anomaly
detection literature {[}34{]} and appropriate for the binary
classification task. However, these metrics do not fully capture
operational utility. Metrics such as mean time to detection (MTTD),
analyst investigation time per alert, and the false discovery rate under
realistic traffic volumes would provide a more comprehensive operational
assessment.

\section{5.7 Chapter Summary}\label{chapter-summary}

This chapter has discussed the experimental results presented in Chapter
4 in the context of the research objectives, the existing literature,
and practical deployment considerations. The five research objectives
have been achieved: (1) the 14-feature behavioral feature set was
designed and validated, with temporal features demonstrated to be the
most discriminative; (2) three unsupervised models were trained and
compared, with IF selected as the primary model for its balance of
performance and operational properties; (3) the EWMA-Adaptive Percentile
dynamic threshold was developed and shown to enable early prediction of
low-and-slow attacks after 3--5 attempts; (4) the 9-service Docker
architecture was built with sub-2-second end-to-end latency; and (5)
five attack scenarios were evaluated with robust detection across all
types. The comparison with 12 existing studies positions this work as a
contribution to the field through its combination of dataset realism,
methodological novelty, and engineering completeness. Six limitations
were identified, motivating the future work directions presented in
Chapter 6.

\chapter{CHAPTER 6: CONCLUSION AND FUTURE
WORK}\label{chapter-6-conclusion-and-future-work}

\section{6.1 Conclusion}\label{conclusion}

This thesis has presented the design, implementation, and comprehensive
evaluation of an intelligent SSH brute-force attack detection and
prevention system that combines unsupervised machine learning with
adaptive dynamic thresholding for early attack prediction. The system
addresses a well-documented gap in the cybersecurity landscape: the
inability of conventional, static-threshold tools such as Fail2Ban
{[}27{]} to detect sophisticated attack variants --- particularly
low-and-slow brute-force, distributed brute-force, and credential
stuffing attacks --- and the persistent disconnect between machine
learning research prototypes and deployable production systems {[}13{]}.

The research was guided by four research questions. The answers to each,
grounded in the experimental evidence presented in Chapter 4 and
discussed in Chapter 5, are summarized below.

\textbf{RQ1: How effectively does the Isolation Forest algorithm detect
SSH brute-force attacks compared to LOF and One-Class SVM?}

All three unsupervised anomaly detection models achieve high recall
(96.75\%--100\%) on the hybrid dataset of 174,250 SSH log lines,
confirming the viability of the semi-supervised paradigm --- training
exclusively on normal data without labeled attack samples --- for SSH
brute-force detection. After optimization, the Isolation Forest achieves
an accuracy of 90.31\%, an F1-score of 93.74\%, a recall of 96.75\%, and
a false positive rate of 29.00\%. OCSVM achieves the highest F1-score of
94.55\% and accuracy of 91.38\%, while LOF achieves perfect recall
(100\%) but exhibits the highest FPR of 67.10\%. The Isolation Forest
was selected as the primary production model due to its computational
efficiency (O(n log n)), its production of smooth continuous anomaly
scores suitable for the dynamic threshold, and its lack of
distributional assumptions --- properties critical for real-time
deployment {[}35, 36{]}.

\textbf{RQ2: To what extent does the EWMA-Adaptive Percentile dynamic
threshold improve detection performance relative to traditional static
thresholds?}

The dynamic threshold (alpha = 0.3, base\_percentile = 95,
sensitivity\_factor = 1.5) provides three capabilities absent from
static thresholds: self-adaptation to traffic pattern changes, reduction
of burst false positives during legitimate anomalous activity, and early
warning through EWMA accumulation. The most significant advantage is the
detection of low-and-slow attacks: the EWMA mechanism issues early
warnings after 3--5 attempts (approximately 2--3 minutes), whereas
Fail2Ban with default settings (maxretry = 5, findtime = 600 s) fails to
detect such attacks when attempts are spaced at intervals exceeding 2
minutes {[}27{]}.

\textbf{RQ3: What degree of early prediction does the system achieve
across different attack scenarios?}

Five attack scenarios of escalating difficulty were evaluated. Basic
brute-force and dictionary attacks were detected immediately (100\%
within the first 1-minute window). Distributed attacks and credential
stuffing achieved detection rates exceeding 90\% and 95\%, respectively.
Low-and-slow attacks --- the most challenging case --- were detected
through EWMA accumulation, with early warnings issued after 3--5
attempts. This graduated detection capability, ranging from immediate
response for high-speed attacks to early warning for stealthy attacks,
represents a significant advancement over purely reactive systems.

\textbf{RQ4: Can the integrated architecture meet real-time monitoring
and automated response requirements?}

The 9-service Docker architecture achieves an end-to-end detection
latency of under 2 seconds, with AI processing stages completing in
under 100 ms (Table 4.14). The five-stage pipeline (Ingestion,
Aggregation, Scoring, Decision, Action) processes SSH authentication
events from log parsing through anomaly scoring to automated IP blocking
via Fail2Ban, operating continuously without manual intervention. The
system was validated with 51 passing unit tests across all components
and 5 integrated attack scenario evaluations, confirming both functional
correctness and operational readiness.

In summary, this research demonstrates that the combination of Isolation
Forest with EWMA-Adaptive Percentile dynamic thresholding constitutes an
effective and practical approach for SSH brute-force detection and early
prediction. The system surpasses traditional tools such as Fail2Ban in
detection capability --- particularly for low-and-slow and distributed
attacks --- while providing early prediction functionality that enables
security administrators to respond proactively before attacks cause
damage. The containerized architecture, comprehensive feature
engineering, and demonstrated real-world performance position the system
for deployment in production environments.

\section{6.2 Contributions Summary}\label{contributions-summary}

The research makes the following five principal contributions:

\textbf{Contribution 1: Semi-supervised anomaly detection system for SSH
brute-force attacks.} The system employs an Isolation Forest model
trained exclusively on normal SSH behavioral data, achieving an F1-score
of 93.74\%, recall of 96.75\%, and FPR of 29.00\% after optimization.
The systematic comparison of three unsupervised algorithms (IF, LOF,
OCSVM) on the same real-world SSH dataset provides an empirical basis
for algorithm selection in the SSH anomaly detection domain,
demonstrating that IF offers the best trade-off between detection
performance and operational properties for real-time deployment {[}35,
41{]}.

\textbf{Contribution 2: EWMA-Adaptive Percentile dynamic threshold
algorithm.} The thesis proposes and validates a novel dynamic
thresholding mechanism that combines EWMA trend tracking with adaptive
percentile computation. The algorithm provides early prediction of
low-and-slow attacks (detection after 3--5 attempts, approximately 2--3
minutes), self-adaptation to traffic pattern changes, and graduated
two-level response (EARLY\_WARNING and ALERT). This mechanism addresses
the fundamental limitation of static thresholds and has not been
previously applied to SSH brute-force detection {[}42, 55{]}.

\textbf{Contribution 3: Empirical demonstration of temporal feature
dominance.} Feature importance analysis reveals that temporal features
--- session\_duration\_mean (5.50\%), min\_inter\_attempt\_time
(3.86\%), mean\_inter\_attempt\_time (2.61\%) --- are far more
discriminative than the count-based features traditionally relied upon
by tools such as Fail2Ban. This finding, corroborating Javed and Paxson
{[}62{]} and Starov et al.~{[}54{]}, provides quantitative evidence for
prioritizing timing information in SSH security monitoring systems.

\textbf{Contribution 4: End-to-end integrated system architecture.} The
complete detection and prevention system is implemented as a
Docker-based microservices architecture comprising 9 services (FastAPI,
React, Elasticsearch, Logstash, Kibana, Fail2Ban, Redis, PostgreSQL,
Nginx) with sub-2-second end-to-end latency. This engineering
contribution directly addresses the research-to-deployment gap
identified by Sommer and Paxson {[}13{]} and Nassif et al.~{[}31{]} as a
persistent barrier to practical ML adoption in cybersecurity.

\textbf{Contribution 5: Evaluation on real-world data with practical
attack scenarios.} The evaluation uses a hybrid dataset of 174,250 SSH
log lines combining real honeypot attack data (119,729 lines from 679
unique IPs across 5 days) with simulated normal data (54,521 lines from
64 users). Five attack scenarios of escalating difficulty were designed,
executed, and analyzed, providing a more ecologically valid evaluation
than the synthetic benchmark datasets (NSL-KDD, CICIDS2017) prevalent in
the literature {[}53{]}.

\section{6.3 Future Work}\label{future-work}

The limitations identified in Section 5.6 and the insights gained during
this research motivate the following directions for future
investigation.

\subsection{6.3.1 Deep Learning for Improved
Detection}\label{deep-learning-for-improved-detection}

The application of deep learning architectures to SSH anomaly detection
represents the most promising direction for improving detection
performance, particularly for reducing the false positive rate while
maintaining high recall. Three specific architectures merit
investigation.

\textbf{LSTM-Autoencoders.} Long Short-Term Memory networks can model
sequential dependencies in SSH authentication event streams that the
current feature-based approach (which aggregates events into fixed
5-minute windows) may lose. An LSTM-Autoencoder trained on sequences of
normal authentication events would learn to reconstruct normal patterns;
sequences that produce high reconstruction error would be flagged as
anomalous. Satoh et al.~{[}57{]} achieved a recall of 97.2\% with an
LSTM approach for SSH dictionary attack detection, suggesting that deep
learning could match or exceed the current system's recall while
potentially reducing the FPR below 29\%.

\textbf{Transformer-based models.} The self-attention mechanism of
Transformer architectures {[}33{]} can capture long-range dependencies
between authentication events that are distant in time but semantically
related (e.g., reconnaissance probes preceding an attack campaign).
Transformers have achieved state-of-the-art results in multiple sequence
modeling domains and have recently been applied to anomaly detection in
time series data.

\textbf{Hybrid ensemble approaches.} Rather than replacing the Isolation
Forest entirely, a promising architecture would combine the IF's
computational efficiency for initial screening with a deep learning
model for refined classification of ambiguous cases. Events with anomaly
scores near the decision boundary --- where false positives are most
likely --- could be forwarded to a more expressive model for secondary
analysis, combining the throughput of IF with the accuracy of deep
learning.

\subsection{6.3.2 Online Learning and Continuous
Adaptation}\label{online-learning-and-continuous-adaptation}

The current system requires periodic retraining to address concept drift
as both attack patterns and normal user behavior evolve over time
{[}8{]}. This limitation can be addressed through online or incremental
learning techniques.

\textbf{Incremental Isolation Forest.} The standard Isolation Forest
algorithm does not support incremental updates; adding new training data
requires rebuilding the entire ensemble from scratch. Recent work on
streaming Isolation Forest variants enables the incorporation of new
data points by selectively replacing the oldest trees in the ensemble
with trees built on recent data. This approach would enable continuous
model adaptation with bounded computational cost.

\textbf{Online One-Class SVM.} Online variants of One-Class SVM, such as
those based on stochastic gradient descent or incremental support vector
updates, can adapt the decision boundary as new normal data arrives.
This would complement the Isolation Forest by providing a continuously
updated secondary model for ensemble scoring.

\textbf{Adaptive retraining policies.} Rather than retraining on a fixed
schedule, the system could monitor detection metrics (FPR, alert rate,
score distribution statistics) and trigger retraining only when
significant distribution shift is detected. Statistical tests such as
the Page-Hinkley test or ADWIN (Adaptive Windowing) {[}55{]} could be
employed to detect distribution shifts in the anomaly score stream.

\subsection{6.3.3 Multi-Protocol
Extension}\label{multi-protocol-extension}

The current system is designed specifically for SSH brute-force
detection. However, the architecture --- log parser, feature extractor,
anomaly detection model, dynamic threshold, automated response --- is
generalizable to brute-force detection on other authentication
protocols:

\begin{itemize}
\tightlist
\item
  \textbf{RDP (Remote Desktop Protocol).} RDP brute-force attacks target
  Windows servers and produce Windows Event Log entries (Event IDs 4625,
  4624) analogous to SSH auth.log entries. A custom parser and
  RDP-specific features (e.g., NLA negotiation timing, session type
  distribution) would enable detection.
\item
  \textbf{FTP (File Transfer Protocol).} FTP brute-force attacks produce
  authentication failure logs (vsftpd, ProFTPD, Pure-FTPd) with timing
  and count characteristics similar to SSH attacks. The existing feature
  set would require minimal adaptation.
\item
  \textbf{HTTP authentication.} Web application login brute-force
  attacks produce HTTP 401/403 responses with timing patterns amenable
  to the same temporal feature analysis. Integration with web server
  access logs (Apache, Nginx) or web application frameworks would extend
  the system's coverage.
\item
  \textbf{SMTP authentication.} Email server brute-force attacks target
  SMTP AUTH mechanisms and produce authentication failure logs with
  similar characteristics.
\end{itemize}

A multi-protocol system would share the core detection engine (Isolation
Forest, dynamic threshold, Fail2Ban integration) while maintaining
protocol-specific parsers and feature extractors, significantly
increasing practical utility for organizations with diverse service
exposure.

\subsection{6.3.4 Federated Detection}\label{federated-detection}

Implementing federated learning {[}33{]} would enable multiple servers
to collaboratively improve detection models without centralizing
sensitive log data. This approach is relevant in three scenarios:

\begin{itemize}
\tightlist
\item
  \textbf{Distributed organizations.} Organizations with SSH servers in
  multiple geographic locations or security domains may face privacy,
  regulatory, or bandwidth constraints that prevent log centralization.
  Federated learning enables each server to train a local model and
  share only model parameters (not raw data) with a central aggregator.
\item
  \textbf{Industry collaboration.} Multiple organizations facing similar
  SSH threats could collaboratively train a shared model without
  exposing their internal traffic patterns, creating a more robust and
  generalizable detection capability.
\item
  \textbf{Coordinated attack detection.} Distributed attacks that
  simultaneously target multiple servers are difficult to detect from
  any single server's perspective. Federated detection enables
  cross-server correlation of anomaly patterns without centralized data
  collection.
\end{itemize}

\subsection{6.3.5 Threat Intelligence
Integration}\label{threat-intelligence-integration}

Combining the behavioral anomaly detection system with external threat
intelligence feeds would provide a multi-evidence classification
framework. Specific integration points include:

\begin{itemize}
\tightlist
\item
  \textbf{IP reputation databases.} Services such as AbuseIPDB {[}28{]},
  Shodan, and GreyNoise maintain continuously updated databases of IP
  addresses associated with malicious activity. An IP flagged by
  multiple reputation services could receive a prior probability
  adjustment, effectively lowering the anomaly score threshold required
  to trigger an alert.
\item
  \textbf{STIX/TAXII feeds.} Structured Threat Information eXpression
  (STIX) and Trusted Automated eXchange of Intelligence Information
  (TAXII) provide standardized formats for sharing cyber threat
  intelligence. The system could subscribe to SSH-specific STIX feeds to
  receive indicators of compromise (IoCs) such as attacker IP ranges,
  username wordlists, and tool signatures.
\item
  \textbf{Dark web monitoring.} Credential dumps and SSH key leaks
  published on dark web forums could inform the system's detection
  priorities, enabling proactive monitoring for compromised accounts.
\end{itemize}

This integration would transform the system from a purely behavioral
detector into a contextually aware threat assessment platform,
potentially reducing the FPR significantly by corroborating behavioral
anomalies with external intelligence.

\subsection{6.3.6 SOAR Integration}\label{soar-integration}

Integrating the system with Security Orchestration, Automation and
Response (SOAR) platforms such as Splunk SOAR, Palo Alto XSOAR, or IBM
Security QRadar SOAR would automate the full incident response
lifecycle:

\begin{itemize}
\tightlist
\item
  \textbf{Automated playbook execution.} Upon detection of an
  ALERT-level event, the SOAR platform could execute predefined
  playbooks: blocking the offending IP at the firewall (not just via
  Fail2Ban), initiating packet capture for forensic analysis, querying
  threat intelligence APIs for enrichment, and creating a case in the
  incident management system.
\item
  \textbf{Notification and escalation.} Automated notification through
  email, Slack, PagerDuty, or other channels according to organizational
  escalation procedures, with severity-based routing (EARLY\_WARNING to
  Tier 1, ALERT to Tier 2).
\item
  \textbf{Compliance reporting.} Automated generation of incident
  reports for regulatory compliance frameworks (ISO 27001, SOC 2, GDPR),
  including timeline reconstruction, evidence preservation, and response
  documentation.
\item
  \textbf{Feedback loop.} Analyst disposition of alerts (true positive,
  false positive, inconclusive) could be fed back to the detection model
  for continuous performance improvement through active learning.
\end{itemize}

\subsection{6.3.7 Explainability and
Interpretability}\label{explainability-and-interpretability}

The current system produces anomaly scores and binary classifications
(normal/attack) but provides limited explanation of why a particular
session or IP was flagged. Adding explainability features would
significantly increase utility for security analysts:

\begin{itemize}
\tightlist
\item
  \textbf{SHAP (SHapley Additive exPlanations) values.} Computing SHAP
  values for each prediction would identify which features contributed
  most to the anomaly score. An analyst who sees that an IP was flagged
  because of unusually short session durations (session\_duration\_mean
  = 1.2 s vs.~normal mean of 340 s) and regular inter-attempt timing
  (std\_inter\_attempt\_time = 0.3 s vs.~normal mean of 45 s) can make a
  more informed triage decision than one who sees only a numerical
  score.
\item
  \textbf{Natural language explanations.} Translating SHAP-based feature
  contributions into human-readable explanations (e.g., ``This IP was
  flagged because session durations are 99.6\% shorter than normal and
  inter-attempt timing is 99.3\% more regular than normal'') would
  further reduce the cognitive burden on analysts.
\item
  \textbf{Visual attention maps.} For deep learning extensions (Section
  6.3.1), attention weights from Transformer models could visualize
  which authentication events in a sequence contributed most to the
  anomaly decision.
\end{itemize}

\subsection{6.3.8 Long-Term IP Profiling}\label{long-term-ip-profiling}

The current 5-minute sliding window approach is effective for most
attack scenarios but limited for extremely slow attacks that spread over
hours or days. A long-term IP profiling mechanism would complement the
short-term detection:

\begin{itemize}
\tightlist
\item
  \textbf{Behavioral baselines per IP.} Maintaining a historical profile
  of each IP's SSH behavior (typical session durations, login times,
  success rates) would enable detection of deviations from that IP's
  individual baseline, rather than from the global baseline.
\item
  \textbf{Reputation scoring.} A continuously updated per-IP reputation
  score, combining behavioral history with external threat intelligence,
  would enable risk-based adaptive thresholding: IPs with clean
  histories receive higher thresholds (fewer false positives), while IPs
  with suspicious histories receive lower thresholds (higher
  sensitivity).
\item
  \textbf{Campaign detection.} Correlating anomalous events across
  multiple IPs over extended time windows could reveal coordinated
  attack campaigns that individual window-based detection misses.
\end{itemize}

\subsection{6.3.9 Real-World Deployment
Validation}\label{real-world-deployment-validation}

While the experimental results are encouraging, the ultimate validation
of the system requires deployment in a production environment with live
SSH traffic over an extended period. A pilot deployment at FPT
University or a partner organization would provide empirical evidence on
several questions that controlled experiments cannot fully answer: How
does the FPR behave over weeks and months of continuous operation? How
effectively does the dynamic threshold adapt to genuine shifts in
organizational usage patterns? How do administrators interact with the
alert system? How does the system perform against adversaries who
actively attempt to evade AI-based detection? A structured pilot program
with defined success criteria, instrumented data collection, and
periodic review would provide the strongest evidence for production
readiness.

\section{REFERENCES}\label{references}

{[}1{]} S. Morgan, ``Cybercrime to cost the world \$10.5 trillion
annually by 2025,'' \emph{Cybersecurity Ventures}, Nov.~2020.
{[}Online{]}. Available:
https://cybersecurityventures.com/cybercrime-damages-6-trillion-by-2021/

{[}2{]} T. Ylonen and C. Lonvick, ``The Secure Shell (SSH) Protocol
Architecture,'' RFC 4251, Internet Engineering Task Force, Jan.~2006.

{[}3{]} SANS Internet Storm Center, ``DShield: Top 10 Target Ports,''
{[}Online{]}. Available: https://isc.sans.edu/top10.html. {[}Accessed:
Jan.~2025{]}.

{[}4{]} Rapid7, ``2023 Attack Intelligence Report,'' \emph{Rapid7
Research}, 2023.

{[}5{]} National Cyber Security Center (NCSC), ``Annual Report on
Cybersecurity in Vietnam,'' Ministry of Information and Communications,
Hanoi, Vietnam, 2023.

{[}6{]} D. R. Tsai, A. Y. Chang, and S. H. Wang, ``A study of SSH brute
force attack defense,'' \emph{Journal of Information Security and
Applications}, vol.~49, pp.~102--113, Dec.~2019.

{[}7{]} M. Najafabadi, T. M. Khoshgoftaar, C. Calvert, and C. Kemp,
``Detection of SSH brute force attacks using aggregated netflow data,''
in \emph{Proc. IEEE 14th Int. Conf. Machine Learning and Applications
(ICMLA)}, Miami, FL, USA, 2015, pp.~283--288.

{[}8{]} A. L. Buczak and E. Guven, ``A survey of data mining and machine
learning methods for cyber security intrusion detection,'' \emph{IEEE
Communications Surveys \& Tutorials}, vol.~18, no. 2, pp.~1153--1176,
2nd Quart. 2016.

{[}9{]} M. A. F. Pimentel, D. A. Clifton, L. Clifton, and L. Tarassenko,
``A review of novelty detection,'' \emph{Signal Processing}, vol.~99,
pp.~215--249, Jun.~2014.

{[}10{]} A. Simoiu, C. Gates, J. Bonneau, and S. Goel, ``I was told to
buy a software or lose my computer. I ignored it: A study of
ransomware,'' in \emph{Proc. 15th Symp. Usable Privacy and Security
(SOUPS)}, Santa Clara, CA, USA, 2019, pp.~155--174.

{[}11{]} J. Jang-Jaccard and S. Nepal, ``A survey of emerging threats in
cybersecurity,'' \emph{Journal of Computer and System Sciences},
vol.~80, no. 5, pp.~973--993, Aug.~2014.

{[}12{]} F. Syed, M. Bashir, and A. Sharaff, ``Machine learning
approaches for intrusion detection in IoT: A comprehensive survey,''
\emph{Journal of King Saud University -- Computer and Information
Sciences}, vol.~34, no. 10, pp.~9656--9688, Nov.~2022.

{[}13{]} R. Sommer and V. Paxson, ``Outside the closed world: On using
machine learning for network intrusion detection,'' in \emph{Proc. IEEE
Symp. Security and Privacy (S\&P)}, Oakland, CA, USA, 2010,
pp.~305--316.

{[}14{]} T. Ylonen, ``SSH -- Secure login connections over the
Internet,'' in \emph{Proc. 6th USENIX Security Symp.}, San Jose, CA,
USA, 1996, pp.~37--42.

{[}15{]} D. J. Barrett, R. E. Silverman, and R. G. Byrnes, \emph{SSH,
The Secure Shell: The Definitive Guide}, 2nd ed.~Sebastopol, CA, USA:
O'Reilly Media, 2005.

{[}16{]} T. Ylonen and C. Lonvick, ``The Secure Shell (SSH)
Authentication Protocol,'' RFC 4252, Internet Engineering Task Force,
Jan.~2006.

{[}17{]} T. Ylonen and C. Lonvick, ``The Secure Shell (SSH) Connection
Protocol,'' RFC 4254, Internet Engineering Task Force, Jan.~2006.

{[}18{]} OpenSSH, ``sshd\_config -- OpenSSH daemon configuration file,''
\emph{OpenBSD Manual Pages}. {[}Online{]}. Available:
https://man.openbsd.org/sshd\_config

{[}19{]} D. Florencio and C. Herley, ``A large-scale study of web
password habits,'' in \emph{Proc. 16th Int. Conf. World Wide Web (WWW)},
Banff, AB, Canada, 2007, pp.~657--666.

{[}20{]} M. Durmuth, T. Kranz, and M. Mannan, ``On the real-world
effectiveness of SSH brute-force attacks,'' in \emph{Proc. NDSS Workshop
on Usable Security (USEC)}, San Diego, CA, USA, 2015.

{[}21{]} A. Sperotto, G. Schaffrath, R. Sadre, C. Morariu, A. Pras, and
B. Stiller, ``An overview of IP flow-based intrusion detection,''
\emph{IEEE Communications Surveys \& Tutorials}, vol.~12, no. 3,
pp.~343--356, 3rd Quart. 2010.

{[}22{]} M. Bishop, ``A taxonomy of Unix password attacks,'' in
\emph{Proc. Computer Security Applications Conf. (ACSAC)}, New Orleans,
LA, USA, 1995.

{[}23{]} J. Owens and J. Matthews, ``A study of passwords and methods
used in brute-force SSH attacks,'' in \emph{Proc. USENIX Workshop on
Large-Scale Exploits and Emergent Threats (LEET)}, San Francisco, CA,
USA, 2008.

{[}24{]} D. Wang, Z. Zhang, P. Wang, J. Yan, and X. Huang, ``Targeted
online password guessing: An underestimated threat,'' in \emph{Proc. ACM
SIGSAC Conf. Computer and Communications Security (CCS)}, Vienna,
Austria, 2016, pp.~1242--1254.

{[}25{]} B. Cheswick and S. M. Bellovin, \emph{Firewalls and Internet
Security: Repelling the Wily Hacker}, 2nd ed.~Reading, MA, USA:
Addison-Wesley, 2003.

{[}26{]} A. K. Das, J. Bonneau, M. Caesar, N. Borisov, and X. Wang,
``The tangled web of password reuse,'' in \emph{Proc. Network and
Distributed System Security Symp. (NDSS)}, San Diego, CA, USA, 2014.

{[}27{]} Fail2Ban Project, ``Fail2Ban documentation,'' {[}Online{]}.
Available: https://www.fail2ban.org/. {[}Accessed: Jan.~2025{]}.

{[}28{]} AbuseIPDB, ``AbuseIPDB -- IP address threat intelligence,''
{[}Online{]}. Available: https://www.abuseipdb.com/. {[}Accessed:
Jan.~2025{]}.

{[}29{]} M. Roesch, ``Snort -- Lightweight intrusion detection for
networks,'' in \emph{Proc. USENIX Systems Administration Conf. (LISA)},
Seattle, WA, USA, 1999.

{[}30{]} M. Krzywinski, ``Port knocking: Network authentication across
closed ports,'' \emph{SysAdmin Magazine}, vol.~12, pp.~12--17, 2003.

{[}31{]} P. Mishra, V. Varadharajan, U. Tupakula, and E. S. Pilli, ``A
detailed investigation and analysis of using machine learning techniques
for intrusion detection,'' \emph{IEEE Communications Surveys \&
Tutorials}, vol.~21, no. 1, pp.~686--728, 1st Quart. 2019.

{[}32{]} M. Ahmed, A. Naser Mahmood, and J. Hu, ``A survey of network
anomaly detection techniques,'' \emph{Journal of Network and Computer
Applications}, vol.~60, pp.~19--31, Jan.~2016.

{[}33{]} G. Pang, C. Shen, L. Cao, and A. van den Hengel, ``Deep
learning for anomaly detection: A review,'' \emph{ACM Computing
Surveys}, vol.~54, no. 2, Art. no. 38, Mar.~2021.

{[}34{]} V. Chandola, A. Banerjee, and V. Kumar, ``Anomaly detection: A
survey,'' \emph{ACM Computing Surveys}, vol.~41, no. 3, Art. no. 15,
Jul.~2009.

{[}35{]} F. T. Liu, K. M. Ting, and Z.-H. Zhou, ``Isolation-based
anomaly detection,'' \emph{ACM Trans. Knowledge Discovery from Data},
vol.~6, no. 1, Art. no. 3, Mar.~2012.

{[}36{]} F. T. Liu, K. M. Ting, and Z.-H. Zhou, ``Isolation Forest,'' in
\emph{Proc. IEEE Int. Conf. Data Mining (ICDM)}, Pisa, Italy, 2008,
pp.~413--422.

{[}37{]} S. Hariri, M. C. Kind, and R. J. Brunner, ``Extended Isolation
Forest,'' \emph{IEEE Trans. Knowledge and Data Engineering}, vol.~33,
no. 4, pp.~1479--1489, Apr.~2021.

{[}38{]} M. M. Breunig, H.-P. Kriegel, R. T. Ng, and J. Sander, ``LOF:
Identifying density-based local outliers,'' in \emph{Proc. ACM SIGMOD
Int. Conf. Management of Data}, Dallas, TX, USA, 2000, pp.~93--104.

{[}39{]} B. Scholkopf, J. C. Platt, J. Shawe-Taylor, A. J. Smola, and R.
C. Williamson, ``Estimating the support of a high-dimensional
distribution,'' \emph{Neural Computation}, vol.~13, no. 7,
pp.~1443--1471, Jul.~2001.

{[}40{]} D. M. J. Tax and R. P. W. Duin, ``Support vector data
description,'' \emph{Machine Learning}, vol.~54, no. 1, pp.~45--66,
Jan.~2004.

{[}41{]} M. Goldstein and S. Uchida, ``A comparative evaluation of
unsupervised anomaly detection algorithms for multivariate data,''
\emph{PLOS ONE}, vol.~11, no. 4, Art. no. e0152173, Apr.~2016.

{[}42{]} S. W. Roberts, ``Control chart tests based on geometric moving
averages,'' \emph{Technometrics}, vol.~1, no. 3, pp.~239--250,
Aug.~1959.

{[}43{]} P. Casas, J. Mazel, and P. Owezarski, ``Unsupervised network
intrusion detection systems: Detecting the unknown without knowledge,''
\emph{Computer Communications}, vol.~35, no. 7, pp.~772--783, Apr.~2012.

{[}44{]} C. Gormley and Z. Tong, \emph{Elasticsearch: The Definitive
Guide}. Sebastopol, CA, USA: O'Reilly Media, 2015.

{[}45{]} Elastic, ``Elasticsearch Reference,'' {[}Online{]}. Available:
https://www.elastic.co/guide/en/elasticsearch/reference/current/.
{[}Accessed: Jan.~2025{]}.

{[}46{]} Elastic, ``Logstash Reference,'' {[}Online{]}. Available:
https://www.elastic.co/guide/en/logstash/current/. {[}Accessed:
Jan.~2025{]}.

{[}47{]} Elastic, ``Kibana Guide,'' {[}Online{]}. Available:
https://www.elastic.co/guide/en/kibana/current/. {[}Accessed:
Jan.~2025{]}.

{[}48{]} D. Gonzalez, T. Hayajneh, and M. Carpenter, ``ELK-based
security analytics for anomaly detection in IoT environments,''
\emph{IEEE Access}, vol.~9, pp.~120827--120841, 2021.

{[}49{]} A. Chuvakin, K. Schmidt, and C. Phillips, \emph{Logging and Log
Management: The Authoritative Guide to Understanding the Concepts
Surrounding Logging and Log Management}. Waltham, MA, USA: Syngress,
2012.

{[}50{]} Elastic, ``Machine Learning in the Elastic Stack,''
{[}Online{]}. Available:
https://www.elastic.co/what-is/elasticsearch-machine-learning.
{[}Accessed: Jan.~2025{]}.

{[}51{]} R. Hofstede, A. Pras, and A. Sperotto, ``Flow-based SSH
compromise detection,'' in \emph{Proc. IFIP/IEEE Int. Symp. Integrated
Network Management (IM)}, 2018.

{[}52{]} P. Kumari and R. Jain, ``Isolation Forest based anomaly
detection for IoT systems,'' \emph{Journal of King Saud University --
Computer and Information Sciences}, vol.~34, no. 8, pp.~5765--5774,
Sep.~2022.

{[}53{]} N. Moustafa and J. Slay, ``The evaluation of Network Anomaly
Detection Systems: Statistical analysis of the UNSW-NB15 and the KDD99
data sets,'' \emph{Information Security Journal: A Global Perspective},
vol.~25, no. 1--3, pp.~18--31, 2016.

{[}54{]} O. Starov, Y. Zhu, and N. Nikiforakis, ``Detecting SSH
brute-force attacks using temporal behavioral analysis,'' in \emph{Proc.
IEEE Conf. Communications and Network Security (CNS)}, Washington, DC,
USA, 2019.

{[}55{]} S. Ahmad, A. Lavin, S. Purdy, and Z. Agha, ``Unsupervised
real-time anomaly detection for streaming data,'' \emph{Neurocomputing},
vol.~262, pp.~134--147, Nov.~2017.

{[}56{]} A. Sperotto, R. Sadre, F. van Vliet, and A. Pras, ``A labeled
data set for flow-based intrusion detection,'' in \emph{Proc. IEEE Int.
Workshop on IP Operations and Management (IPOM)}, Venice, Italy, 2009.

{[}57{]} A. Satoh, Y. Nakamura, and T. Ikenaga, ``SSH dictionary attack
detection based on flow analysis using deep learning,'' \emph{IEEE
Access}, vol.~10, pp.~75614--75627, 2022.

{[}58{]} V. T. Nguyen and M. Q. Tran, ``Application of machine learning
in network intrusion detection for Vietnamese enterprise networks,''
\emph{Journal of Science and Technology -- University of Danang},
vol.~19, no. 5, pp.~45--51, 2021.

{[}59{]} H. V. Le, T. H. Nguyen, and M. D. Pham, ``Building a network
security monitoring system using ELK Stack for small and medium
enterprises in Vietnam,'' \emph{Journal of Information and
Communications Technology}, no. 3, pp.~56--64, 2022.

{[}60{]} N. H. Pham, ``Research on SSH brute-force prevention solutions
for government information systems,'' M.S. thesis, Academy of
Cryptography Techniques, Hanoi, Vietnam, 2020.

{[}61{]} D. K. Tran and T. T. H. Nguyen, ``Application of Isolation
Forest in anomaly detection on system log data,'' \emph{Journal of
Science Research and Development}, vol.~2, no. 4, pp.~112--121, 2023.

{[}62{]} M. Javed and V. Paxson, ``Detecting stealthy, distributed SSH
brute-forcing,'' in \emph{Proc. ACM SIGSAC Conf. Computer and
Communications Security (CCS)}, Berlin, Germany, 2013, pp.~85--96.

{[}63{]} S. S. Khan and M. G. Madden, ``One-class classification:
Taxonomy of study and review of techniques,'' \emph{Knowledge
Engineering Review}, vol.~29, no. 3, pp.~345--374, Sep.~2014.

\section{APPENDICES}\label{appendices}

\subsection{Appendix A: Feature Extraction Code
Snippets}\label{appendix-a-feature-extraction-code-snippets}

\emph{The complete feature extraction module is implemented in Python
using pandas and numpy. The code processes SSH auth.log entries through
regular expression parsing, aggregates events into 5-minute sliding
windows per source IP, and computes the 14 behavioral features described
in Chapter 3.}

\subsection{Appendix B: Docker Compose
Configuration}\label{appendix-b-docker-compose-configuration}

\emph{The Docker Compose configuration file defines the 9 services
(FastAPI, React, Elasticsearch, Logstash, Kibana, Fail2Ban, Redis,
PostgreSQL, Nginx) with their respective ports, environment variables,
volume mounts, and inter-service dependencies.}

\subsection{Appendix C: Isolation Forest Hyperparameter Tuning
Results}\label{appendix-c-isolation-forest-hyperparameter-tuning-results}

\emph{Detailed results of the grid search for Isolation Forest
hyperparameters, including performance metrics for all evaluated
parameter combinations of n\_estimators (100, 200, 300, 500),
max\_samples (256, 512, 1024), max\_features (0.25, 0.5, 0.75, 1.0), and
contamination (auto, 0.01, 0.05, 0.1).}

\subsection{Appendix D: Attack Simulation
Scripts}\label{appendix-d-attack-simulation-scripts}

\emph{The five attack scenario simulation scripts, implemented in Python
using the Paramiko SSH library, with configurable parameters for attack
speed, IP distribution, username lists, and password dictionaries.}

\subsection{Appendix E: Logstash Pipeline
Configuration}\label{appendix-e-logstash-pipeline-configuration}

\emph{The Logstash pipeline configuration for parsing SSH auth.log
entries, including Grok patterns for all event types (Failed password,
Accepted password, Invalid user, etc.) and the Elasticsearch output
configuration with index templates.}

\subsection{Appendix F: Dynamic Threshold Parameter Sensitivity
Analysis}\label{appendix-f-dynamic-threshold-parameter-sensitivity-analysis}

\emph{Complete sensitivity analysis results for the four dynamic
threshold parameters (alpha, base\_percentile, sensitivity\_factor,
lookback), including performance metrics across the full range of
evaluated values and visual representations of threshold behavior under
different parameter settings.}

\end{document}
